\documentclass[conference]{IEEEtran}
\IEEEoverridecommandlockouts
\renewcommand\IEEEkeywordsname{Keywords}
% The preceding line is only needed to identify funding in the first footnote. If that is unneeded, please comment it out.

%for tables
\usepackage{booktabs}
\usepackage{multirow}
%end for tables

\usepackage{cite}
\usepackage{algorithm}
\usepackage{array}
\usepackage{longtable}
\usepackage{amsmath,amssymb,amsfonts} %Mathematical Symbol,Equation
\usepackage{algorithm,algpseudocode}
\usepackage{hyperref}
\hypersetup{
    colorlinks=true,
    linkcolor=blue,
    filecolor=magenta,      
    urlcolor=cyan,
    citecolor=blue,
    pdftitle={Unifying Event Evolution Graphs and Reasoning Traces via Semantic Entailment: A Dual-Phase Architecture for Judicial GraphRAG},
    pdfpagemode=FullScreen,
}

\usepackage{graphicx} %Adding picture & figure
\usepackage{float} %Figure Alignment
\usepackage{textcomp} %Text formatting
\usepackage{eso-pic}
\usepackage{xcolor}
\usepackage{comment}
\usepackage{csquotes}

% Vector Graphics & Diagrams (TikZ & PGFPlots & Standalone)
\usepackage{standalone}
% \usepackage{lmodern} % Removed to prevent dual-font clash with IEEEtran standard font
\usepackage{anyfontsize}
\usepackage{microtype}
\usepackage{tikz}
\usetikzlibrary{arrows.meta, positioning, shapes.geometric, shadows, fit, calc, backgrounds}
\usepackage{pgfplots}
\pgfplotsset{compat=1.14}

\def\BibTeX{{\rm B\kern-.05em{\sc i\kern-.025em b}\kern-.08em
    T\kern-.1667em\lower.7ex\hbox{E}\kern-.125emX}}
    
\makeatletter
%%%%%%%%%%%%%%%%%%%%%%%%%%%%%% User specified LaTeX commands.
\newcommand{\snippet}[1]{``\texttt{#1}''}

\usepackage{eso-pic}
\newcommand\AtPageUpperMyright[1]{\AtPageUpperLeft{%
 \put(\LenToUnit{0.5\paperwidth},\LenToUnit{-1cm}){%
     \parbox{0.5\textwidth}{\raggedleft\fontsize{9}{11}\selectfont #1}}%
 }}%
\newcommand{\conf}[1]{%
\AddToShipoutPictureBG*{%
\AtPageUpperMyright{#1}
}
}
\usepackage{fancyhdr}

\setlength {\marginparwidth }{2cm} 
\setlength{\dbltextfloatsep}{5pt plus 1pt minus 1pt}
\setlength{\textfloatsep}{5pt plus 1pt minus 1pt} 

\makeatletter
\def\footnoterule{\kern-3\p@
  \hrule \@width 2in \kern 2.6\p@} % the \hrule is .4pt high
\makeatother

% Auto-generated canonical statistics macros
\input{generated_stats}
\input{generated_obscure_stats}
\input{generated_corpus_stats_20k}
\input{phase3_corrections} % Phase 3 Batch 01 overrides (R35; Category D Holm family)
\input{phase3_corrections_b2} % Phase 3 Batch 02 overrides (R28; R10)
\input{phase3_corrections_b3} % Phase 3 Batch 03 overrides (R13; Table XXI, section VII-M)

\begin{document}

% Header settings
\title{Unifying Event Evolution Graphs and Reasoning Traces via Semantic Entailment: A Dual-Phase Architecture for Judicial GraphRAG}

%Authors
\author{\IEEEauthorblockN{Md. Emam Zafor Saadik and Md. Raihan Molla}
\IEEEauthorblockA{\textit{Department of Computer Science and Engineering}\\ 
\textit{Bangladesh University of Business and Technology (BUBT), Dhaka, Bangladesh}}}

\maketitle

% The paper headers and footers
\thispagestyle{fancyplain}
\fancyhf{}
\lhead{\fancyplain{}{}}
\lfoot{\fancyplain{}{}}
\renewcommand{\headrulewidth}{0pt}

%%%%%%%%%%%%%%%%%%%%%%%%%%%%%%%%%%%%%%%%%%%%%%%%%%%%%%%%%%%%%%%%%%%%%%%%%%%%%%%%
\begin{abstract}
Retrieval-Augmented Generation (RAG) in law must trace multi-hop precedent chains and must not present overruled doctrine as current law, yet common evaluators reward entity overlap rather than grounded reasoning. We introduce TERRA (Temporal-Evolution and Reasoning-Trace RAG), an experimental dual-phase GraphRAG architecture. Offline, it builds an Event Evolution Graph over \CorpusTotalOpinions\ Supreme Court opinions connected by \CorpusPrecedentEdges\ directed citation edges, plus a procedural reasoning-trace index. Online, a small-model Traffic Cop router sends single-point lookups down a fast path and multi-hop questions through citation traversal, and an autonomous semantic entailment grader evaluates retrieved evidence before generation. We evaluate on a \numUniqueQueries-query benchmark against four baselines (\numTotalEvals\ evaluations) using two non-Llama judges, paired Wilcoxon tests and Benjamini--Hochberg correction. TERRA scores higher overall faithfulness ($\TerraOverallFaithfulnessMean$ vs.\ $\HippoRAGOverallFaithfulnessMean$ for the next best system, overlapping 95\% CIs) and answer-only faithfulness ($\TerraAnswerOnlyFaithfulnessMean$, computed on differing answered subsets), and significantly outperforms all three retrieval baselines on Category~B multi-hop reasoning ($\Delta$ from $\CatBTerraVsHippoRAGDelta$ [HippoRAG] and $\CatBTerraVsFlatDelta$ [Flat] to $\CatBTerraVsGraphRAGDelta$ [GraphRAG]; FDR $q<0.05$), although a no-retrieval Direct LLM scores significantly higher on Category~B ($\Delta=\CatBTerraVsDirectLLMDelta$, $q=\CatBTerraVsDirectLLMAdjQ$). On adversarial Category~D queries TERRA safely handles \TerraCatDSafeHandlingRate\ of traps (\TerraCatDRejectionRate\ explicit refusals, \TerraCatDRefutationRate\ verified refutations) with \TerraCatDFatalErrors\ fatal errors, but refuses \TerraCatAFalseRefusals\ Category~A and \TerraCatBFalseRefusals\ Category~B queries, at \TerraLatencyMedianSeconds~s median latency. On \ObscCohortSize\ obscure overruled precedents, Direct LLM safe handling falls to \ObscDirectSafePct\ versus \ObscFullSafePct\ for full TERRA, but TERRA without curated overruling labels drops to \ObscUnlabeledSafePct, not significantly different from flat retrieval ($p\,\ObscPrimaryMcNemarPRel$); safety depends on curated annotations. Absent licensed attorney evaluation, human review remains an essential deployment prerequisite.
\end{abstract}

\begin{IEEEkeywords}
\textcolor{black}{\textit{Retrieval-Augmented Generation, Large Language Models, Temporal Knowledge Graph, Agentic Reasoning, Semantic Entailment, Explainable AI, Data Leakage Mitigation}}
\end{IEEEkeywords}

\section{Introduction}
\label{sec:Introduction}
Retrieval-Augmented Generation (RAG) has rapidly evolved into a foundational architecture for knowledge-intensive artificial intelligence. By explicitly decoupling the parametric memory of Large Language Models (LLMs) from dynamic external knowledge bases, RAG enables real-time factual grounding and transparent source attribution without the prohibitive computational costs of continuous model retraining. In high-stakes, knowledge-intensive environments---where factual errors and hallucinations carry severe real-world consequences---the fidelity of the retrieved context is not merely an optimization preference, but a strict operational requirement. As enterprise applications scale, the early paradigm of naive ``retrieve-then-generate'' pipelines characterized by static vector databases and fixed-length document chunking has proven structurally inadequate.

The fundamental bottleneck in modern RAG systems arises when navigating complex data landscapes characterized by continuous knowledge drift and multi-step reasoning demands. Real-world data evolves dynamically; consequently, static vector indices inevitably accumulate conflicting, time-stamped versions of the same fact. When queried, standard retrievers surface all of these conflicting versions simultaneously based on semantic similarity, leaving the frozen generative model unable to resolve the chronological authority, which induces severe temporal hallucinations. Furthermore, when processing complex logical derivations, standard pipelines retrieve fragmented textual snippets that sever precedential and doctrinal dependencies, proving entirely insufficient for tasks that require deep, multi-hop reasoning.

Recent advancements have attempted to address these bottlenecks in isolated silos, exposing critical failure modes. While structural graph-based retrievers organize facts chronologically, they act merely as static data foundations lacking the procedural logic to guide multi-hop generative reasoning. Conversely, prior paradigms retrieving intermediate structured reasoning traces significantly enhance logical deduction but remain temporally blind and have been primarily restricted to quantitative domains like mathematics or coding. More critically, current Agentic RAG frameworks rely heavily on corrective evaluators that fail fundamentally on conceptual tasks. Mathematical explainability analyses, such as SHAP, have proven that standard evaluators act merely as superficial named entity alignment detectors, failing to comprehend genuine semantic relevance. Compounding this, existing query routers are largely restricted to rigid, binary-channel systems (e.g., pure dense versus pure sparse search), resulting in severe evidence-intent mismatches when unrestricted search floods the context window with noisy, irrelevant data.

To bridge these isolated gaps, this paper introduces an integrated dual-phase GraphRAG architecture purpose-built for dynamically evolving, knowledge-intensive environments (with empirical graph traversal stability across temporal snapshot expansions evaluated in Section~\ref{subsec:temporal-update-stability}). The proposed system confines all computationally heavy operations to an offline preprocessing phase, constructing a dual-index corpus that cross-links an Event Evolution Graph (EEG)---which encodes chronological precedent relationships across \CorpusTotalOpinions\ authentic judicial decisions and \CorpusPrecedentEdges\ directed citation edges---with a node-level procedural reasoning trace database that captures structured legal analysis trajectories (constitutional issue, procedural history, legal test, holding). During online inference, a small language model acts as a Traffic Cop router, dynamically allocating retrieval priority across text chunks, graph subgraphs, and procedural reasoning traces, enabling a bimodal execution profile (gross median \TerraFastPathMedianSeconds~s on Category~A factual lookups vs. \TerraHardPathMedianSeconds~s on Category~B multi-hop precedent chains). Before the final response generation, an autonomous semantic entailment grader replaces flawed keyword-matching heuristics, verifying the true conceptual and structural relevance of the retrieved evidence.

The proposed framework was rigorously evaluated against comparative baselines, including Direct LLM generation and Flat Vector RAG, alongside component-level ablation studies utilizing comprehensive temporal and multi-perspective reasoning benchmarks. Experimental results demonstrate that the model significantly mitigates the attention dilution and ``context cliff'' limitations inherent in processing overly broad retrieved contexts. By seamlessly aligning chronological facts with procedural reasoning, the model substantially reduces temporal hallucinations in grounded retrieval, providing an experimental research framework for context fidelity in temporal legal AI on historical SCOTUS jurisprudence.

\subsection{Summary of Contributions}
The primary contributions of this work are as follows:
\begin{itemize}
    \item \textbf{An Authentic Dual-Index Architecture:} We design and implement an authentic dual-index judicial GraphRAG corpus that cross-links a structural Event Evolution Graph (EEG) of \CorpusTotalOpinions\ authentic SCOTUS decisions (\CorpusPrecedentEdges\ citation edges) with a node-level procedural reasoning trace database, supporting joint chronological and logical evidence alignment with zero synthetic cases.
    \item \textbf{Traffic Cop Routing with Coarse Domain Pre-filtering:} We introduce an SLM-driven complexity router paired with an optional coarse proximity flag ($d_{L2} < 1.3$ + Title Match), calibrated offline as a scale-dependent domain-boundary signal that avoids static allowlists; in the canonical run this flag terminated no query before retrieval (calibration and measured stage activity are detailed in Section~\ref{subsec:computational-budget} and Section~\ref{subsec:firewall-analysis}).
    \item \textbf{NLI-Based Entailment Evaluation and XAI Attribution:} We replace superficial keyword-matching heuristics with a specialized Natural Language Inference (NLI) grader that evaluates retrieved evidence to intercept ungrounded generations, providing qualitative reasoning traces and citation provenance.
    \item \textbf{Data Leakage Mitigation Framework:} We formalize explicit safeguards against temporal precedent inversion, benchmark retrieval contamination, and cross-domain feature bleeding.
    \item \textbf{Empirical Rigor and Statistical Validation:} Through benchmarking (\numTotalEvals\ evaluations across \numUniqueQueries\ queries $\times$ \numComparativePipelines\ pipelines at $T=0.0$) with paired Wilcoxon tests, Cohen's $d_z$, and 95\% bootstrap CIs, TERRA achieves the highest Answer-Only Faithfulness among RAG architectures ($\TerraAnswerOnlyFaithfulness$ vs.\ HippoRAG's $\HippoRAGAnswerOnlyFaithfulness$), outperforms Flat Vector RAG ($\Delta = \CatBTerraVsFlatDelta$, $p = \CatBTerraVsFlatWilcoxonP$, $d_z = \CatBTerraVsFlatDz$) and HippoRAG ($\Delta = \CatBTerraVsHippoRAGDelta$, $p = \CatBTerraVsHippoRAGWilcoxonP$, $d_z = \CatBTerraVsHippoRAGDz$) on Category~B reasoning, delivers $\TerraCatDSafeHandlingRate$ adversarial safe handling ($\TerraCatDSafeHandling$: \TerraCatDRejection\ refusals, \TerraCatDRefutations\ refutations) with \TerraCatDFatalErrors\ fatal errors vs.\ $\BaselineFatalMin$--$\BaselineFatalMax$ for baselines, and incurs \TerraCatAFalseRefusals\ false refusals on Category~A (\TerraCatBFalseRefusals\ on Category~B). Evaluations rely on automated LLM judges without licensed attorneys, establishing human review as an essential prerequisite (Section~\ref{subsec:limitations}).

\end{itemize}


\section{Related Work}
\label{sec:related-work}
The foundational paradigm of Retrieval-Augmented Generation (RAG) \cite{lewis2020retrieval, guu2020realm} decoupled parametric LLM memory from external knowledge bases. However, standard RAG suffers from severe vulnerabilities when managing continuous knowledge drift and multi-step reasoning. To structure retrieved evidence, recent work has pivoted toward graph-based architectures. Microsoft Research introduced GraphRAG \cite{edge2024graphrag}, which constructs hierarchical entity-relation communities for global query summarization. Similarly, Sun et al. \cite{sun2024thinkongraph} developed Think-on-Graph (ToG) to execute iterative LLM-guided reasoning over knowledge graphs. In temporal domains, Liu and Cao \cite{liu2026raglearningunderstandinglimits} introduced Chronos to track dynamic events in an Event Evolution Graph (EEG), Zhang et al. \cite{zhang2025respectingtemporalcausalconsistencyentityevent} proposed E2RAG to preserve temporal-causal consistency across entity-event subgraphs, and Zhu et al. \cite{zhu2025rightanswerrighttime} developed STAR-RAG for time-sensitive rule graph summarization. Within the legal domain, recent architectures have explored structured representations: Chen et al.~\cite{chen2026legalgraphrag} proposed LegalGraphRAG, deploying multi-agent reasoning over a tripartite legal knowledge graph (ontology, fact, and rule layers) for statutory civil and criminal law, while de Martim~\cite{demartim2025satgraphrag} introduced SAT-Graph RAG to trace diachronic statutory norm versions across formal constitutional amendments. However, these systems model codified statutory rules or formal legislative amendment events; in contrast, common-law judicial systems require tracking the longitudinal evolution of precedent across citation networks, where earlier decisions remain presumptive authority unless explicitly distinguished, limited, or overruled. Furthermore, static graph models retrieve structural summaries but lack the procedural logic required to execute multi-hop reasoning over evolving precedents.

To bridge this gap in procedural logic, recent research has explored augmenting retrieval with intermediate reasoning trajectories rather than static factual snippets alone. Li et al. \cite{li2025cotragintegratingchainthought} demonstrated this through the CoT-RAG framework, establishing that retrieving structured reasoning sub-cases and decision trees significantly mitigates logical biases compared to unstructured text retrieval. The broader transition toward System 2 agentic reasoning has been mapped by Lin et al. \cite{lin2025comprehensivesurveyreinforcementlearningbased} and Liang et al. \cite{liang2025reasoningrag12}, who detail multi-stage pipelines where autonomous agents dynamically identify knowledge gaps and generate intermediate reasoning traces to inform subsequent search actions. In parallel, the legal AI community has established rigorous evaluation benchmarks to measure procedural reasoning capabilities: CaseHOLD \cite{zheng2021casehold} evaluates judicial holding identification, while the Common Law Case Law Retrieval Benchmark (CLERC) \cite{hou2024clerc} measures multi-step precedent retrieval. These benchmarks confirm that standard lexical or dense retrievers frequently retrieve irrelevant or superseded authority unless procedural reasoning context is explicitly preserved. However, existing trace-retrieval frameworks have been predominantly evaluated on quantitative domains or static statutory collections, leaving them largely untested in longitudinal common-law environments where precedent evolution alters substantive legal authority over time.

Beyond the challenge of retrieving complex reasoning traces and temporal data, the successful integration of this multi-faceted information relies entirely on the ``connective tissue'' of the RAG pipeline: query routers and autonomous evaluators. Hashmi \cite{hashmi2026adaptivequeryroutingtierbased} highlighted the rigidity of standard single-channel retrieval, introducing an Adaptive Hybrid Retrieval framework to route queries dynamically, noting that pure vector similarity frequently fails to capture structural navigation. More critically, the evaluation mechanisms governing these routers are fundamentally flawed. Yalavarthi \cite{yalavarthi2026opensourcereproductionexplainabilityanalysis} conducted a SHAP explainability analysis on the widely adopted Corrective RAG (CRAG) architecture, mathematically proving that its standard T5-based evaluator functions purely as a named entity alignment detector. Consequently, the CRAG evaluator misclassified $88.3\%$ of conceptual science tasks as ambiguous and achieved only $5.0\%$ accuracy on non-entity queries. The real-world consequences of these shallow heuristics were exposed by Putta et al. \cite{putta2026casefactsbenchmarklegalfactchecking} via the CaseFacts benchmark, which demonstrated that augmenting models with unrestricted web search actively degrades performance in legal fact-checking compared to closed-book baselines. Because standard retrievers rely on lexical matching, they flood the context window with noisy, overruled precedents, failing to recognize when a fact has been logically overturned. These findings highlight a critical architectural gap: current systems lack the genuine Natural Language Inference (NLI) capabilities required for true semantic entailment.

Ultimately, a comprehensive review of this literature reveals that addressing these vulnerabilities in isolation is insufficient; no prior architecture unifies temporal precedent management, node-level procedural reasoning trace retrieval, and semantic-entailment evaluation into a single, cohesive pipeline. Graph models track chronological topology but lack procedural reasoning logic; trace models enhance reasoning but lack temporal awareness; and current routers rely on superficial keyword-matching that triggers catastrophic failure in high-stakes domains. The proposed architecture directly addresses this convergence gap, as summarized in the comprehensive architectural and doctrinal feature comparison of Table~\ref{tab:system_comparison}. By cross-linking a temporal Event Evolution Graph with a node-level procedural reasoning trace database, it ensures that structural precedent topology (what happened and when) and operative judicial analysis (issue, test, and holding) are retrieved as a semantically aligned pair. Furthermore, it replaces flawed entity-alignment evaluators with an autonomous NLI entailment grader and utilizes a continuous three-way Traffic Cop router to dynamically balance textual, graphical, and procedural data. This design establishes a dual-phase framework capable of addressing continuous knowledge drift (empirically examined under incremental corpus snapshot expansions in Section~\ref{subsec:temporal-update-stability}) and multi-hop reasoning without requiring continuous foundational model retraining.

\input{tables/table_system_comparison}

\section{System Architecture and Methodology}
\label{sec:methodology}

TERRA (Temporal-Evolution and Reasoning-Trace RAG) is designed as a two-phase, closed-loop GraphRAG architecture for temporal legal reasoning on historical common-law jurisprudence. The core engineering philosophy strictly separates offline corpus ingestion and knowledge structure construction from online query routing and inference. This separation ensures that all computationally intensive graph traversals, citation link extractions, and reasoning trace generations are performed offline, maintaining responsive online inference with a bimodal latency profile (median \TerraFastPathMedianSeconds~s on factual lookups vs. \TerraHardPathMedianSeconds~s on multi-hop citation chains).


\begin{figure*}[t]
    \centering
    \resizebox{0.98\textwidth}{!}{\input{figure_1_system_architecture}}
    \caption{TERRA two-phase architecture. Phase 1 (offline) constructs the Event Evolution Graph (EEG) and ChromaDB vector index from the \CorpusTotalOpinions-opinion authentic SCOTUS citation corpus. Phase 2 (online) routes user queries through a Traffic Cop Router with an optional proximity flag, performs BFS citation-graph traversal, validates context via NLI entailment, and conditionally generates grounded answers or triggers a safety refusal.}
    \label{fig:system_architecture}
\end{figure*}

\subsection{Phase 1: Offline Ingestion and Knowledge Graph Construction}
\label{subsec:phase1-ingestion}

\subsubsection{Authentic Judicial Opinion Corpus (\CorpusTotalOpinions\ Cases)}
To ensure legal authenticity and reflect the true topological complexity of judicial precedent, TERRA operates on a scaled corpus of \CorpusTotalOpinions\ authentic, full-text Supreme Court of the United States (SCOTUS) opinions spanning \CorpusTemporalSpan\ (Volumes 50--448). The corpus was curated from an initial harvest of 23,010 public-domain federal decisions acquired from the open CourtListener and Harvard Caselaw repositories. Under 17 U.S.C. \S~105, all federal judicial opinions are public domain. To eliminate non-precedential noise while preserving doctrinal completeness, opinions were filtered under a substantive length mandate ($\text{character count} \ge 4{,}000$). Across the \IssueTwentyTotalVolumesScanned\ scanned U.S. Reports volumes (\IssueTwentyTotalDecisionsScanned\ total decisions), \IssueTwentySubFiveHundredCharCount\ records are shorter than 500 characters (\IssueTwentySubFiveHundredCharPct, representing procedural orders, summary dispositions, and docket records, subject to audit verification). However, filtering strictly by character length induces corpus length-filtering bias: it excludes \IssueTwentyIntermediateExcludedCount\ intermediate decisions ($500$--$3{,}999$ characters, \IssueTwentyIntermediateExcludedPct), among which \IssueTwentyIntermediatePageRankGeNinetyNineFive\ rank in the top $0.5\%$ of citation authority under Harvard CAP PageRank (including procedural landmarks such as \textit{Haines v. Kerner} on pro se pleadings and \textit{United States v. Diebold, Inc.} on summary judgment inferences, competency standards such as \textit{Dusky v. United States}, tax milestones such as \textit{Lucas v. Earl}, and desegregation orders such as \textit{Sipuel v. Board of Regents} and \textit{Gayle v. Browder}). The production corpus contains zero synthetic cases ($0\%$ synthetic supplement). Within this global network, a core set of 34 landmark civil rights decisions---spanning the constitutional evolution from \textit{Dred Scott v. Sandford} (1857), through \textit{Plessy v. Ferguson} (1896), \textit{Missouri ex rel. Gaines v. Canada} (1938), \textit{Brown v. Board of Education} (1954), to \textit{Swann v. Charlotte-Mecklenburg} (1971)---serves as the primary ground-truth anchor set for benchmarking multi-hop doctrinal reasoning.

\subsubsection{Citation Parsing and Directed Edge Classification}
Inter-case legal citations embedded within opinion texts are extracted using the \texttt{eyecite} library~\cite{cushman2021eyecite}. We emphasize that \texttt{eyecite} serves strictly as a token-level citation extraction and reporter normalization utility; it establishes citation existence ($A \xrightarrow{\text{CITES}} B$) rather than inferring substantive precedential treatment, doctrinal approval, or negative authority. Citation strings are normalized via \texttt{cite.corrected\_citation()} (e.g., \snippet{347 U.S. 483}) and cross-referenced against exact opinion identifier hashes. Across \CorpusTotalOpinions\ decisions, the verified Event Evolution Graph contains \CorpusPrecedentEdges\ directed edges (\IssueTwentyOneTotalCitesEdges\ precedential \texttt{CITES} links and \IssueTwentyOneOverrulesEdges\ \texttt{OVERRULES} relations). Precedential citations strictly maintain forward chronological validity at calendar-year granularity with zero cross-year forward violations (\IssueTwentyOneCrossYearViolationsPct); \IssueTwentyOneSameYearEdges\ edges (\IssueTwentyOneSameYearEdgesPct) connect opinions within the same calendar year (including \IssueTwentyOneIntraYearReversals\ intra-year day-of-month reporting artifacts). The legacy claim of ``99.91\% chronological validity / 0.09\% noise'' has unresolved provenance in project logs (likely an unverified decimal approximation of full-date ordering, $\IssueTwentyOneStrictlyPriorDateEdges/\IssueTwentyOneTotalCitesEdges = {}$\IssueTwentyOneStrictlyPriorDatePct). A reproducible stratified audit across 14 decades ($N=\IssueTwentyOneAuditSampleN$, seed \IssueTwentyOneAuditSeed) confirms verbatim citation presence in \IssueTwentyOneAuditCitationPresentPct\ of sampled edges (95\% CI: \IssueTwentyOneAuditCitationPresentCI), with \IssueTwentyOneAuditTargetConfirmedPct\ verified direct precedential reliance (95\% CI: \IssueTwentyOneAuditTargetConfirmedCI) and \IssueTwentyOneAuditQuotationPct\ citations within third-party quotations.

The graph contains \IssueTwentyOneTotalCitesEdges\ generic \texttt{CITES} edges and \IssueTwentyOneOverrulesEdges\ authoritative typed \texttt{OVERRULES} edges imported from Congressional Research Service (CRS) annotations ($A \xrightarrow{\text{OVERRULES}} B$, denoting explicit judicial abrogation of prior doctrine). The 19 Category-D pairs (e.g., \textit{Brown v. Board of Education} explicitly invalidating \textit{Plessy v. Ferguson}'s separate-but-equal doctrine) are a curated evaluation subset of those \IssueTwentyOneOverrulesEdges\ graph relations, populating the Category~D adversarial benchmark (Section~\ref{subsec:benchmark-design}) to evaluate temporal reasoning and deception defense.

The resulting Event Evolution Graph (EEG) is serialized as a NetworkX directed graph (\texttt{DiGraph}) stored in \texttt{dataset\_20k\_cap/terra\_eeg\_index\_20k.json}.

\subsubsection{Procedural Reasoning Trace Synthesis and Vector Serialization}
To provide generation with structured procedural legal context rather than static text snippets alone, each case is paired with a node-level Structured Legal Reasoning Trace (SLRT), formalizing procedural case briefs grounded in Harvard Caselaw Access Project (CAP) records:
\begin{enumerate}
    \item \textbf{Curated Landmark Prototypes ($N=\TraceCuratedCount$):} During early system development, \TraceCuratedCount\ landmark constitutional decisions received curated procedural reasoning traces generated via instruction-tuned models (\texttt{gemini-3.1-flash-lite} baseline; \texttt{mistral.ministral-3-8b-instruct} validated variants for active benchmark seeds) at pinned temperature $T=0.0$. Traces were constrained to a 150-word ceiling and strictly followed a five-point doctrinal schema: (1) core constitutional question, (2) governing constitutional provisions and review standard, (3) central legal holding, (4) direct precedent connection (built upon, distinguished, or overruled), and (5) longitudinal doctrinal significance.
    \item \textbf{Deterministic Procedural Substrate at Scale ($N=\TraceProceduralCount$, \TraceProceduralPct):} When scaling to the substantive \CorpusTotalOpinions-case retrieval corpus, to completely eliminate generative hallucination risk from the retrieval index, all \TraceProceduralCount\ opinions are serialized deterministically via a procedural function (\texttt{build\_structured\_trace}) parameterized directly by authenticated CAP metadata: official case title (\texttt{name\_abbreviation}), decision year (\texttt{decision\_date}), and legal subject taxonomy (\texttt{category}). This function generates a five-step analytical workflow (identifying the operative legal question under the specific domain jurisprudence, scrutinizing governing Supreme Court precedent, evaluating the factual record under appellate standards, harmonizing conflicting doctrines, and pronouncing the holding). Thus, the substantive retrieval substrate contains zero generative LLM hallucinations.
\end{enumerate}

Each vector document serializes the case title header, its SLRT (under legacy internal key \texttt{"Thinking Trace:"}), and an authentic opinion preamble snippet:
\begin{equation}
\begin{aligned}
\text{Document} = \; & \texttt{\{title\} (\{year\})} \\
& \oplus \; \texttt{"\textbackslash nThinking Trace: \{trace\}"} \\
& \oplus \; \texttt{"\textbackslash nOpinion Snippet: \{text\_snip[:1000]\}"}
\end{aligned}
\label{eq:trace_doc}
\end{equation}
All \CorpusTotalOpinions\ documents are embedded using Sentence-Transformers (\texttt{sentence-transformers/all-MiniLM-L6-v2}~\cite{reimers2019sentencetransformers}, \TraceEmbeddingDimensions-dimensional dense vectors, normalized Euclidean $L_2$ metric) and indexed into a local ChromaDB \texttt{PersistentClient} vector store (\texttt{dataset\_20k\_cap/vector\_db\_20k}), establishing an aligned dual-index architecture.

To verify evidence integrity before treating the trace store as a trusted retrieval layer, we executed an exhaustive empirical audit across all \TraceTotalNodes\ index nodes. Forensic substring verification against raw CAP casebody JSON files confirms an aggregate span grounding rate of \TraceSpanGroundingRate\ (\TraceSpanGroundingCount\ exact verbatim matches), decisively surpassing the pre-registered $\ge \TraceSpanGroundingThreshold$ acceptance floor (expected 100.0\%, with exactly \TraceSpanDiscrepanciesCount\ discrepancies, \TraceSpanDiscrepanciesPct, attributable to short historical per curiam docket orders lacking opinion text bodies in CAP). Furthermore, a dual-critic factual consistency and temporal anachronism audit across an $N=\TraceAuditSampleCount$ verification cohort---comprising a 100.0\% complete census of all curated landmark prototypes indexed in the 20,000-case graph ($N=\TraceCuratedCensusCount$, with 3 short historical per curiam orders omitted by CAP's $\ge 4,000$ character length filter) and a stratified random sample of the procedural majority ($N=\TraceStratifiedProcCount$ across all seven subject categories)---confirmed \TracePopWeightedAnachronismFreePct\ population-weighted temporal integrity (zero cited precedents post-date the target decision's ruling date; automated check $\text{decision\_date}(\text{precedent}) \le \text{decision\_date}(\text{target})$) and \TracePopWeightedFactualPct\ population-weighted factual consistency. Complete per-case landmark provenance tags, generator prompts, and audit ledgers are published in the open-source release (\texttt{docs/REASONING\_TRACE\_SPECIFICATION.md}).

\subsection{Phase 2: Online Query Routing and Inference Engine}
\label{subsec:phase2-inference}

\subsubsection{Traffic Cop Router and Multi-Tier Domain Boundary Filtering}
Incoming user queries $q$ are initially processed by the Traffic Cop Router, an SLM-driven classification module powered by Google Gemma 3 12B Instruction-tuned (\texttt{\RouterModel}) operating at temperature $T=0.0$ under structured Pydantic schema validation (\texttt{RoutingDecision}):
\begin{equation}
\begin{aligned}
\text{Route}(q) = \begin{cases} 
\text{EASY}, & \text{if } \text{Low Complexity} \land \text{In-Domain} \\
\text{HARD}, & \text{if } \text{High Complexity} \lor \text{Out-of-Domain}
\end{cases}
\end{aligned}
\label{eq:routing}
\end{equation}
EASY queries (e.g., single-fact decision years or case author lookups) bypass multi-hop graph expansion and route directly to the Grounded Fast Path, trading complex retrieval for a $\AblationRouterLatencySlowdownFactor\times$ operational speedup in the pilot ablation (Table~\ref{tab:ablation}) while eliminating context bloat from tangential citing precedents. As evaluated standalone in Section~\ref{subsec:router-standalone-eval} and Table~\ref{tab:router_eval}, the router achieves $\RouterAccuracy$ classification accuracy and $\RouterMacroFOne$ macro-$F_1$ across the \numUniqueQueries-query benchmark suite.

To prevent hallucinations on out-of-domain topics without maintaining brittle, static allowlists, TERRA implements a multi-tier \textbf{Domain Boundary Safety Framework} structured across three defense-in-depth layers:
\begin{enumerate}
    \item \textbf{Tier 1: Lexical Title Matching:} Extracts indexed landmark case titles from the active NetworkX EEG nodes and checks whether $q$ contains an exact title substring.
    \item \textbf{Tier 2: Optional Geometric Pre-Filter (Scale-Dependent):} Evaluates the L2 embedding distance $d_{L2}$ of the nearest ChromaDB reasoning trace against an empirical boundary threshold:
    \begin{equation}
    d_{L2}(q, v_{\text{nearest}}) < 1.3
    \label{eq:l2threshold}
    \end{equation}
    Crucially, embedding distance does not prove legal validity or domain membership; rather, it is intended strictly as an optional coarse pre-filter to conserve retrieval and grading compute on obvious out-of-domain inputs, with zero in-domain false refusals in the offline calibration. Empirical auditing on the \CorpusTotalOpinions-case substantive index (Table~\ref{tab:computational_budget}, Panel~B) establishes that in-domain legal queries cluster tightly (mean $0.365 \pm 0.082$, max $0.509$), allowing $\tau = 1.30$ to achieve zero false refusal (\TerraGateFPRTauOneThreeZero, \TerraGateFPRCountTauOneThreeZero) while flagging \TerraGateTPRTauOneThreeZero\ (\TerraGateTPRCountTauOneThreeZero) of out-of-domain probes by distance; in the canonical run no query was terminated before retrieval.
    \item \textbf{Tier 3: Semantic Entailment Verification (Downstream Check):} Because geometric proximity alone cannot verify semantic authority or prevent doctrinal misapplication, a downstream check is applied by the autonomous NLI Smart Grader. If retrieved legal subgraphs fail logical entailment across two expansion attempts, generation terminates deterministically in a safety refusal. In the canonical run, 81 of the 139 refusals followed an empty retrieved context (74 of the 75 Category~C probes) and 58 followed a non-empty context (Section~\ref{subsec:firewall-analysis}).
\end{enumerate}

Within this multi-tier verification framework, a \textit{fatal constraint violation} is formally defined as generating an affirmative, factual-sounding response that validates an overruled legal doctrine, misattributes controlling constitutional authority, or credulously adopts a legally impossible premise embedded within an adversarial prompt (Category D). Conversely, a \textit{false refusal} is formally defined as issuing an unwarranted epistemic refusal on an in-domain factual lookup query (Category A) where controlling historical precedent is fully present and unambiguous in the corpus.

\subsubsection{Multi-Hop Citation Graph Traversal}
For HARD-routed queries requiring legal synthesis, retrieval proceeds via a coordinated vector-graph traversal:
\begin{enumerate}
    \item \textbf{Seed Retrieval and Error Propagation Insulation:} ChromaDB retrieves top-$k$ (initialized to $k=2$) reasoning trace vectors matching $q$, establishing initial entry seed nodes $V_{\text{seed}}$. Crucially, vector retrieval over procedural traces is deployed strictly as a coarse seed discovery mechanism rather than the terminal evidence substrate. Downstream precedent discovery is governed by topological citation graph traversal ($G_{\text{EEG}}$) and substantive holding injection (Section~\ref{subsec:retrieval-quality}), structurally insulating generative synthesis from vector seeding inaccuracies: while dense vector seeding alone achieves only $\CatBFlatMacroNodeRecall$ precedent node recall and $\CatBFlatFullPathRecall$ path recovery, two-hop citation traversal lifts precedent recall to $\CatBTerraMacroNodeRecall$ (\TraceTopologicalAbsLift, \TraceTopologicalRelMultiplier\ relative multiplier) and path recovery to $\CatBTerraFullPathRecall$ (\TracePathAbsLift, \TracePathRelMultiplier\ multiplier). Before context compilation reaches the generator, the NLI Smart Grader is designed to enforce semantic entailment boundaries on retrieved traces; in the canonical run Category~C showed \CatCTerraCleanRefusals\ clean refusals (74 of the 75 probes had an empty retrieved context) and Category~D showed \TerraCatDSafeHandlingRate\ safe handling.
    \item \textbf{BFS Traversal with Hub Capping:} Starting from $V_{\text{seed}}$, Breadth-First Search (BFS) traverses the EEG up to depth $d=2$ across both predecessors and successors:
    \begin{equation}
    \mathcal{C}(q) = V_{\text{seed}} \cup \bigcup_{v \in V_{\text{seed}}} \text{BFS}(v, G_{\text{EEG}}, d=2)
    \label{eq:bfs}
    \end{equation}
    To prevent super-hub nodes (e.g., cases with degree $>50$) from flooding the context, hub expansion is dynamically capped. Traversed node summaries and topological trajectories are compiled chronologically into a unified evidence block $\mathcal{C}(q)$.
\end{enumerate}

\subsubsection{NLI Smart Grader and Self-Correction Expansion Loop}
\label{subsubsec:terra-nli-grader}
Before answer generation, context $\mathcal{C}(q)$ is evaluated by an autonomous NLI Smart Grader. The Smart Grader is implemented as a structured-prompt LLM call to \texttt{databricks-meta-llama-3-3-70b-instruct}~\cite{dubey2024llama3} (temperature $T=0.0$), presented with the user query and retrieved legal context, and instructed to output a JSON object with fields \texttt{entails}~(boolean), \texttt{confidence\_score}~(float), and \texttt{reasoning}~(string). The Smart Grader operates as a specialized verification prompt on the primary 70B parameter model at $T=0.0$, while front-line complexity routing is decoupled into the independent 12B parameter Gemma 3 SLM (\texttt{\RouterModel}), enabling lightweight single-API-call verification without deploying a separate specialized NLI cross-encoder (e.g., DeBERTa~\cite{he2023debertav3}). No independent labeled validation set was constructed for the NLI verifier; the primary empirical evidence of its reliability is the ablation result: removing the entailment gate (Ablation~2, Table~\ref{tab:ablation}) causes adversarial safety rejection to collapse from \AblationFullCatDSafetyRate\ to \AblationTwoCatDSafetyRate, confirming that the gate is the operative safety mechanism. False-non-entailment errors (insufficient context incorrectly passed to the generator) constitute the principal unquantified failure mode; we identify this as a limitation (Section~\ref{subsec:limitations}) and a target for future validation with a labeled entailment test set. The Smart Grader operates with structured JSON validation as follows:
\begin{itemize}
    \item \textbf{Entailment Verified (\texttt{entails == True}):} The verified context is forwarded directly to the generator.
    \item \textbf{Entailment Failed (\texttt{entails == False}):} The engine triggers an autonomous \textbf{Self-Correction Expansion Loop}. Retrieval breadth expands ($k \leftarrow k + 2$), ChromaDB is re-queried, and BFS traversal re-executes over the wider citation neighborhood.
    \item \textbf{Double Failure Termination:} If entailment fails after two expansion attempts, the \textbf{Safety Refusal Firewall} halts generation, outputting a deterministic refusal: \textit{``I apologize, but I do not have sufficient validated legal context in my databases to answer this query accurately without risking hallucination.''}
\end{itemize}

\subsubsection{Data Leakage Mitigation Framework}
\label{subsubsec:data-leakage}
In legal temporal GraphRAG, data leakage represents a major methodological failure mode. TERRA implements four strict safeguards:
\begin{enumerate}
    \item \textbf{Temporal Directionality and Precedent Inversion Protection:} In legal precedent networks, an earlier case cannot cite future jurisprudence. Directed EEG edges enforce chronological integrity: forward \texttt{CITES} edges strictly require $t(v_{\text{cite}}) \ge t(v_{\text{cited}})$. Retrospective invalidations are evaluated via authoritative \texttt{OVERRULES} relations imported from CRS annotations (Section~\ref{subsec:phase1-ingestion}) rather than unrestricted backward edges, preventing the retriever from anachronistically projecting modern legal standards onto historical inquiries.
    \item \textbf{Vector Store Contamination Isolation:} Benchmark test queries, evaluation rubrics, and reference answers are strictly segregated from the retrieval index. ChromaDB contains exclusively authentic opinion text and procedural reasoning traces; zero evaluation queries or answers exist in the vector database.
    \item \textbf{Cross-Domain Feature Bleed Prevention:} Out-of-domain concepts are prevented from contaminating generated answers through a defense-in-depth hierarchy: title matching and a scale-dependent geometric proximity flag are designed as coarse first-line signals (the flag terminated no query in the canonical run), while downstream NLI entailment verification is designed to suppress non-entailed retrieved subgraphs before answer synthesis.
    \item \textbf{Temporal Parametric Containment Firewall:} While modern foundation models encode extensive post-training parametric knowledge, historical legal reasoning requires strict fidelity to jurisprudence as of a specific temporal horizon $T$. To prevent temporal parametric leakage---wherein the generator inadvertently incorporates subsequent abrogations, modern statutory standards, or post-$T$ precedents into historical evaluations---TERRA implements a dual-layer containment firewall. First, graph traversals and retrieved opinion corpora are strictly bounded by provenance timestamps ($\text{KnowledgeProvenance} \le T$), ensuring that retrieved evidence contains no anachronistic materials. Second, candidate generations $\hat{y}_{\text{candidate}}$ are evaluated by an entailment-based NLI gate against contemporaneous opinion evidence prior to output release. As demonstrated in Section~\ref{subsec:temporal-leakage}, this mechanism intercepts anachronistic doctrine and ungrounded hallucinations, reducing operative temporal leakage while preserving historical correctness without requiring costly model retraining or fine-tuning.
\end{enumerate}

\subsubsection{Explainable AI (XAI) and Topological Provenance}
\label{subsubsec:xai-provenance}
Transparent reasoning paths are essential in judicial AI research, particularly where legal conclusions must be traceable to cited authority. TERRA enforces four distinct explainability layers:
\begin{enumerate}
    \item \textbf{Graph Traversal Provenance:} Every synthesized answer is rooted in an explicit multi-hop citation path $\mathcal{P} = \langle v_1, e_1, v_2, \dots, v_k \rangle$ traversed through the EEG, providing a full legal audit trail from landmark anchor to modern precedent.
    \item \textbf{Dynamic Hyperlinked Attribution:} In generated answers, referenced cases possessing indexed node identifiers are dynamically formatted as interactive relative markdown links (e.g., \texttt{[Brown v. Board of Education](/case/347\_0483-01)}), enabling instant auditability against primary legal texts. Across benchmark completions, TERRA emitted \TerraTotalMarkdownLinks\ interactive markdown links with a \TerraMarkdownLinkResolutionRate\ corpus resolution rate (\TerraResolvedMarkdownLinks/\TerraTotalMarkdownLinks) and \TerraMarkdownLinkFidelityRate\ semantic title fidelity (evaluated standalone in Section~\ref{subsec:citation-correctness} and Table~\ref{tab:citation_eval}), while comparative baselines emit zero interactive links.
    \item \textbf{NLI Entailment Transparency:} The Smart Grader outputs structured qualitative justifications (\texttt{faithfulness\_reasoning} and \texttt{confidence\_score}), detailing why retrieved context is legally sufficient or deficient.
    \item \textbf{Routing Decision Transparency:} The Traffic Cop emits explicit reasoning strings explaining the classification into the EASY or HARD execution path.
\end{enumerate}

\subsubsection{Grounded Answer Generation}
When context passes NLI grading, the final answer is generated by \texttt{databricks-meta-llama-3-3-70b-instruct}~\cite{dubey2024llama3} under a strict grounding instruction. Generation temperature is pinned to:
\begin{equation}
T = 0.0
\label{eq:temperature}
\end{equation}
enforcing greedy decoding ($T=0.0$). We note that remote hosted inference APIs may exhibit minor output variability across serving sessions, even at $T=0.0$ (Section~\ref{subsec:limitations}, Limitation~\ref{item:hosted-api-determinism}).

\subsection{Implementation Details}
\label{subsec:implementation-details}
The research prototype is implemented in Python 3.11 using \texttt{openai}, \texttt{chromadb}, \texttt{networkx}, and \texttt{fastapi}. Online execution is served via a FastAPI/Uvicorn endpoint with an in-memory query cache. Telemetry logging captures per-request routing latency, NLI confidence, and firewall status into \texttt{terra\_telemetry\_traces.json}.


\section{Dataset and Corpus Analysis}
\label{sec:dataset-analysis}

To validate the structural and topological integrity of the Event Evolution Graph (EEG), we perform network analysis on the pre-built dual-index corpus. This section outlines the corpus composition, topological graph metrics, and centrality hub properties.

\subsection{Corpus Statistics and Doctrinal Taxonomy}
\label{subsec:corpus-stats}
The indexed legal knowledge base contains \CorpusTotalOpinions\ authentic substantive judicial opinions connected by \CorpusPrecedentEdges\ directed citation edges across \CorpusTemporalSpan\ (Volumes 50--448), partitioned across seven doctrinal categories. To structure graph traversal and domain-routed retrieval, each opinion is assigned to exactly one mutually exclusive primary doctrinal domain via an automated deterministic rule-based taxonomy comprising an exact \IssueTwentyTwoTotalLexiconPhrases-phrase legal lexicon across six substantive domains: Substantive Private Law (\IssueTwentyTwoSubstantivePrivateLawPhrases\ phrases, \IssueTwentyTwoSubstantivePrivateLawCorpus\ opinions, \IssueTwentyTwoSubstantivePrivateLawCorpusPct), Administrative Law and Federal Courts (\IssueTwentyTwoAdminLawPhrases\ phrases, \IssueTwentyTwoAdminLawCorpus\ opinions, \IssueTwentyTwoAdminLawCorpusPct), Criminal Procedure (\IssueTwentyTwoCriminalProcedurePhrases\ phrases, \IssueTwentyTwoCriminalProcedureCorpus\ opinions, \IssueTwentyTwoCriminalProcedureCorpusPct), Commerce Clause and Economic Regulation (\IssueTwentyTwoCommerceClausePhrases\ phrases, \IssueTwentyTwoCommerceClauseCorpus\ opinions, \IssueTwentyTwoCommerceClauseCorpusPct), Civil Rights and Equal Protection (\IssueTwentyTwoCivilRightsPhrases\ phrases, \IssueTwentyTwoCivilRightsCorpus\ opinions, \IssueTwentyTwoCivilRightsCorpusPct), and First Amendment (\IssueTwentyTwoFirstAmendmentPhrases\ phrases, \IssueTwentyTwoFirstAmendmentCorpus\ opinions, \IssueTwentyTwoFirstAmendmentCorpusPct). 

The seventh domain, General Constitutional and Statutory (\IssueTwentyTwoGeneralConstitutionalCorpus\ opinions, \IssueTwentyTwoGeneralConstitutionalCorpusPct), functions as an unkeyed residual fallback category assigned strictly when an opinion contains zero substantive keyword hits across all six explicit lexicons. Positive-score ties between substantive domains are resolved deterministically by documented dictionary insertion order; ties never assign the residual fallback category. In addition to mutually exclusive primary domain partitioning, opinions independently support secondary multi-label doctrinal tags capturing multidimensional jurisprudence.

To audit taxonomy quality without circularity or manual inflation, we evaluated a stratified sample of $N=\IssueTwentyTwoTotalSampleN$ opinions (\IssueTwentyTwoBaselineSampleN\ baseline stratified cases and \IssueTwentyTwoEdgeSampleN\ edge cases: \IssueTwentyTwoTieSampleN\ ties, \IssueTwentyTwoMultiTagSampleN\ multi-tag cases, and \IssueTwentyTwoFallbackSampleN\ residual fallback cases; random seed \IssueTwentyTwoSamplingSeed). Validation was conducted using two independent automated evaluators: Judge~A (a deterministic statutory/constitutional parser scanning opinion texts and official headnotes) and Judge~B (\texttt{databricks-gpt-oss-120b} with logged fallback to \texttt{databricks-qwen3-next-80b-a3b-instruct}; zero LLaMA calls). Inter-evaluator agreement yielded raw agreement of \IssueTwentyTwoIEARawAgreement\ (Cohen's $\kappa = \IssueTwentyTwoIEACohensKappa$). Disagreements and edge cases were adjudicated by a third independent automated adjudicator, Judge~C (\IssueTwentyTwoAdjudicatorModel), evaluated strictly blinded to the assigned taxonomy and prior evaluator votes, with strictly zero endpoint overlap with Judge~B (\IssueTwentyTwoAdjudicatorOverlapCount/\IssueTwentyTwoTotalSampleN). We explicitly disclose that this protocol represents automated model adjudication and parser evaluation, not human legal-expert validation. Against this independent model-derived reference adjudication standard, the primary-domain taxonomy achieves a Horvitz-Thompson prevalence-weighted accuracy of \IssueTwentyTwoHorvitzThompsonAccuracy\ (SE: \IssueTwentyTwoHorvitzThompsonSE, 95\% CI: \IssueTwentyTwoHorvitzThompsonCI). Across the complete \IssueTwentyTwoTotalSampleN-case audit, the taxonomy exhibits \IssueTwentyTwoStrictCoherentPct\ strict coherence (\IssueTwentyTwoStrictCoherentCount/\IssueTwentyTwoTotalSampleN), \IssueTwentyTwoPartialCoherentPct\ partial alignment (\IssueTwentyTwoPartialCoherentCount/\IssueTwentyTwoTotalSampleN), yielding an overall substantively grounded rate of \IssueTwentyTwoSubstantivelyGroundedPct\ (\IssueTwentyTwoSubstantivelyGroundedCount/\IssueTwentyTwoTotalSampleN), with \IssueTwentyTwoIncoherentPct\ incoherent classifications (\IssueTwentyTwoIncoherentCount/\IssueTwentyTwoTotalSampleN). In edge-case strata, tie resolution achieved \IssueTwentyTwoTieAccuracy\ coherence, secondary multi-label tagging reached \IssueTwentyTwoMultiTagPrecision\ precision, \IssueTwentyTwoMultiTagRecall\ recall (F1 score \IssueTwentyTwoMultiTagFOne), and residual fallback assignment exhibited \IssueTwentyTwoFallbackValidityPct\ validity as an unkeyed default. Table~\ref{tab:corpus_stats} summarizes corpus composition and substantive domain coverage across the complete graph structure.

\begin{table}[htbp]
\caption{TERRA Event Evolution Graph (EEG) Corpus Statistics and Doctrinal Taxonomy.}
\label{tab:corpus_stats}
\centering
\resizebox{0.98\columnwidth}{!}{%
\begin{tabular}{lcp{4.1cm}}
\toprule
\textbf{Doctrinal Domain} & \textbf{Opinions} & \textbf{Substantive Coverage (\CorpusTemporalSpan)} \\
\midrule
Substantive Private Law & \CorpusPrivateLawOpinions & Patents, Tax, Admiralty, Commercial Contracts \\
General Constitutional\textsuperscript{*} & \CorpusGeneralConstitutionalOpinions & Separation of Powers, Judicial Review, Supremacy \\
Administrative Law & \CorpusAdminLawOpinions & Regulatory Agencies, Federal Jurisdiction, Procedure \\
Criminal Procedure & \CorpusCriminalProcedureOpinions & 4th, 5th, 6th, and 8th Amendment Jurisprudence \\
Commerce Clause & \CorpusCommerceClauseOpinions & Interstate Commerce, Economic Regulation, Antitrust \\
Civil Rights & \CorpusCivilRightsOpinions & Equal Protection, Due Process, Voting Rights Core \\
First Amendment & \CorpusFirstAmendmentOpinions & Free Speech, Press, Exercise, Establishment Clause \\
\midrule
\textbf{Total Opinions ($|V|$)} & \textbf{\CorpusTotalOpinions} & 100\% Authentic Substantive Opinions (0 Procedural) \\
\textbf{Citation Edges ($|E|$)} & \textbf{\CorpusPrecedentEdges} & Directed Precedent Links (\IssueTwentyOneAuditTargetConfirmedPct\ Direct Precision, 0 Cross-Year Violations) \\
\bottomrule
\end{tabular}%
}
\vspace{1mm}
\footnotesize{\textsuperscript{*}Primary-domain partition: Each opinion is assigned to exactly one mutually exclusive primary domain via a deterministic \IssueTwentyTwoTotalLexiconPhrases-phrase rule-based lexicon, with \texttt{general\_constitutional\_statutory} serving as the residual fallback category for decisions with zero substantive keyword hits (\IssueTwentyTwoGeneralConstitutionalCorpus\ opinions, \IssueTwentyTwoGeneralConstitutionalCorpusPct). Opinions independently support secondary multi-label doctrinal tags.}
\end{table}

\subsection{Graph Topology Analysis}
\label{subsec:graph-topology}
The EEG exhibits a sparse topology with high-degree authority hubs and a heavy-tailed empirical degree distribution, consistent with judicial citation networks. The largest weakly connected component (WCC) spans \CorpusLargestComponentSize\ out of \CorpusTotalOpinions\ nodes (\CorpusLargestWccPercent\ coverage), confirming robust macro-level graph connectivity and ensuring that multi-hop retrieval paths remain navigable across historical eras. Table~\ref{tab:graph_metrics} details the network metrics computed via the \texttt{graph\_analytics} module.

\begin{table}[htbp]
\caption{Event Evolution Graph Topological Properties.}
\label{tab:graph_metrics}
\centering
\resizebox{0.98\columnwidth}{!}{%
\begin{tabular}{lc}
\toprule
\textbf{Metric} & \textbf{Value} \\
\midrule
Total Nodes ($|V|$) & \CorpusTotalOpinions \\
Total Edges ($|E|$) & \CorpusPrecedentEdges \\
Graph Density ($\rho$) & \CorpusGraphDensity \\
Mean Node Degree ($\bar{k}$) & \CorpusMeanNodeDegree \\
Maximum Degree (\CorpusMaxDegreeCase) & \CorpusMaxDegree \\
Average Clustering Coefficient ($C$) & \CorpusAverageClustering \\
Graph Diameter ($d_{\text{max}}$; undirected, largest WCC) & \CorpusGraphDiameter\ hops \\
Average Shortest Path Length ($L$; undirected, largest WCC) & \CorpusAvgShortestPath\ hops \\
Largest Weakly Connected Component (WCC) & \CorpusLargestComponentSize\ (\CorpusLargestWccPercent) \\
Largest Strongly Connected Component (SCC) & \CorpusLargestSccSize \\
\bottomrule
\end{tabular}%
}
\end{table}

The graph diameter of \CorpusGraphDiameter\ hops with an average shortest path length of \CorpusAvgShortestPath\ hops across the \CorpusLargestComponentSize-node component demonstrates compact topological reachability across federal appellate history. This confirms that setting the BFS graph traversal depth to $d=2$ in Phase 2 effectively captures dense local citation neighborhoods without flooding the context window with distant nodes.

\subsection{Hub Node Centrality Analysis}
\label{subsec:hub-analysis}
Betweenness centrality ($C_B$) measures the extent to which a case node lies on the shortest paths between other pairs of cases, identifying structural bridge cases in legal history:
\begin{equation}
C_B(v) = \sum_{s \neq v \neq t} \frac{\sigma_{st}(v)}{\sigma_{st}}
\label{eq:betweenness}
\end{equation}
where $\sigma_{st}$ is the total number of shortest paths from node $s$ to node $t$ and $\sigma_{st}(v)$ is the number of those paths that pass through $v$.

\begin{figure}[htbp]
    \centering
    \resizebox{0.98\columnwidth}{!}{\input{figure_2_hub_centrality}}
    \caption{Top-5 EEG hub nodes by betweenness centrality ($C_B$) in the \CorpusTotalOpinions-node network. All top structural hubs correspond to major federal constitutional precedents, demonstrating that historical bridge cases naturally emerge as topological centers.}
    \label{fig:hub_centrality}
\end{figure}

As shown in Fig.~\ref{fig:hub_centrality}, betweenness centrality analysis on the \CorpusTotalOpinions-node EEG identifies \textit{Simpson v. Shepard} ($C_B = 0.0043$), \textit{Fay v. Noia} ($C_B = 0.0036$), \textit{Ashwander v. Tennessee Valley Authority} ($C_B = 0.0033$), \textit{Crowell v. Benson} ($C_B = 0.0032$), and \textit{Dahnke-Walker Milling Co. v. Bondurant} ($C_B = 0.0032$) as the primary structural bridges in the citation network. High-degree hubs such as \CorpusMaxDegreeCase\ (degree \CorpusMaxDegree), \textit{Baker v. Carr} (degree 245), \textit{Crowell v. Benson} (degree 213), \textit{Nebbia v. New York} (degree 211), and \textit{Fay v. Noia} (degree 204) function as critical statutory and constitutional cross-roads linking disparate doctrinal lines across the historical graph.


\section{Experimental Setup}
\label{sec:experimental-setup}

To rigorously assess TERRA's performance in legal temporal reasoning and safety firewalling, we establish a comprehensive comparative benchmarking suite. This section outlines the benchmark query suite design, baseline architectures, and evaluation metrics.

\subsection{Benchmark Suite Design}
\label{subsec:benchmark-design}
The evaluation suite comprises an expanded benchmark of \numUniqueQueries\ unique queries evaluated across \numComparativePipelines\ comparative pipelines (\numTotalEvals\ total evaluations; see Table~\ref{tab:run_provenance} for execution dates and model identifiers), alongside a foundational $N=\numPilotQueries$ unique-query pilot ablation set (\numPilotRepeatedRuns\ repeated stability trials per query, \numPilotEvals\ total executions across four configurations) structured across four distinct categories designed to isolate factual accuracy, multi-hop reasoning, domain safety, and adversarial robustness. Temperature $T{=}0.0$ was used throughout, but does not guarantee identical outputs across runs (see Section~\ref{subsec:limitations} and Table~\ref{tab:run_provenance}). Quantitative paired text generation metrics (Faithfulness and ROUGE-L) are evaluated across in-domain queries (Categories A and B) that produce generative text outputs, while safety firewall performance is evaluated across out-of-domain and deceptive queries (Categories C and D):
\begin{itemize}
    \item \textbf{Category A --- Single-Point Factual:} Single-fact retrieval queries (e.g., decision years, specific party names) evaluated against researcher-constructed ground-truth reference answers using ROUGE-L.
    \item \textbf{Category B --- Multi-Hop Evolutionary:} Complex legal queries requiring multi-decade precedent tracking across citation chains, evaluated against researcher-constructed reference answers using ROUGE-L. To test non-trivial multi-hop graph traversal while reflecting the small-world hub topology of judicial citation networks, all Category B chains are constrained to 3--5 nodes (2--4 intermediate hops), comprising 63 three-node, 32 four-node, and 5 five-node evolutionary arcs with zero 2-node direct citation pairs.
    \item \textbf{Category C --- Out-of-Context Safety:} Queries covering topics demonstrably absent from indexed U.S. Supreme Court case law (\CorpusTemporalSpan), spanning five distinct non-indexed domains ($N=15$ each, total 75): (1) Natural and Physical Sciences, (2) Foreign Constitutional and International Treaty Law (e.g., European Court of Justice, 1951 Geneva Convention, German Federal Constitutional Court), (3) Modern State and Municipal Administrative Ordinances (e.g., municipal zoning codes, city curb-cut regulations), (4) Contemporary Software Systems and Post-2020 Culture, and (5) Corporate Financial Accounting (GAAP/IFRS arithmetic). Because no controlling SCOTUS precedent exists for these inquiries, the ground-truth system requirement is strict safety refusal.
    \item \textbf{Category D --- Adversarial Deception:} Queries that embed real in-domain case titles alongside fabricated legal premises or temporal falsehoods (e.g., asserting that the 1905 \textit{Lochner} ruling was governed by the 1938 Fair Labor Standards Act, or asking for Section 5 preclearance provisions under the 1883 \textit{Civil Rights Cases}), designed to trick entity-matching routers. The acceptable safe response is either explicit safety refusal or verified substantive legal refutation that actively exposes and refutes the false premise citing authentic authority without endorsing invalidated doctrine.
\end{itemize}

\noindent\textbf{Task Scope vs. Background Corpus Scope Partitioning.} To evaluate multi-step legal temporal reasoning without artificial domain simplifications, we establish a deliberate structural distinction between the broad background retrieval corpus and the evaluated benchmark task scope. While the pre-built EEG encompasses the full \CorpusTotalOpinions-case Supreme Court docket across all seven substantive areas (including \CorpusPrivateLawOpinions\ Substantive Private Law decisions and \CorpusGeneralConstitutionalOpinions\ General Statutory decisions; Table~\ref{tab:corpus_stats}), the multi-hop evolutionary reasoning benchmark (Category~B) intentionally concentrates exclusively on constitutional and federal public-law precedent trajectories (20 queries each across Civil Rights, Criminal Procedure, First Amendment, Commerce Clause, and Administrative Law; \IssueTwentyFiveCatBConstitutionalPublicPct). Across all \IssueTwentyFiveBenchInDomainTotal\ in-domain queries (Categories~A, B, and D), constitutional and public-law questions comprise \IssueTwentyFiveBenchConstitutionalPublicCount\ queries (\IssueTwentyFiveBenchConstitutionalPublicPct), while Substantive Private Law comprises \IssueTwentyFiveBenchPrivateLawCount\ (\IssueTwentyFiveBenchPrivateLawPct) and General Statutory comprises \IssueTwentyFiveBenchGeneralStatutoryCount\ (\IssueTwentyFiveBenchGeneralStatutoryPct). This partitioning reflects our task-design rationale: private-law disputes (e.g., patent claim construction, maritime salvage, diversity contract disputes) were not modeled as 3--5 node multi-decade citation chains under our Category~B construction criteria, whereas constitutional jurisprudence exhibits longitudinal precedent evolution and explicit doctrinal overruling. Retaining the complete 20,000-case repository provides a realistic, high-distractor topological environment where constitutional reasoning must navigate past \IssueTwentyFiveCorpusCombinedPrivateStatutoryPct\ commercial and statutory background nodes without false-positive retrieval.

\noindent\textbf{Category~C Construction Protocol.} The 75 out-of-context queries were pre-registered across the five non-SCOTUS domains ($N=15$ each, Subtypes C1--C5) prior to model evaluation, spanning physical sciences, foreign jurisprudence, municipal ordinances, software systems, and corporate accounting arithmetic. An automated corpus-isolation audit against the 20,000-opinion indexed repository verified 0/75 case title, reporter citation, or substantive holding overlaps. All 75 queries and metadata were locked into \texttt{terra\_350\_benchmark.json} (SHA256: \texttt{8144fa81...}), and the formal domain-labeling protocol and 75-query catalog are published in supplementary documentation.

\noindent\textbf{Category~D Construction Protocol.} The 75 adversarial queries were constructed by the research team using a pre-registered four-subtype taxonomy prior to any model evaluation: Subtype D1 ($N=19$): explicit overruled precedent traps; Subtype D2 ($N=19$): title-substitution traps; Subtype D3 ($N=19$): temporal inversion traps enforcing a $\ge 5$-year historical gap; and Subtype D4 ($N=18$): procedural posture deceptions. Query sampling used \texttt{RANDOM\_SEED = 42} from the 19 verified doctrine-invalidation pairs. All queries were locked into \texttt{terra\_350\_benchmark.json} (SHA256: \texttt{8144fa81...}) before any pipeline was evaluated, with zero queries or reference answers indexed into ChromaDB (Section~\ref{subsubsec:data-leakage}).

In addition to the comparative benchmark, a standalone stress test of 15 queries (10 Out-of-Context and 5 Adversarial) is executed via \texttt{stress\_test.py} to verify firewall boundary consistency.

\enlargethispage{3\baselineskip}
\paragraph{Obscure Overruled Precedent Suite.} To test whether Category~D safety reflects retrieval or the pretraining memorization of famous landmarks, we constructed a supplementary suite of \ObscCohortSize\ obscure overruled-precedent queries. Candidates were the CRS \texttt{OVERRULES} edges of the Event Evolution Graph not used in Category~D (forty-seven edges); three edges whose overruling opinions never mention the target case were excluded as corrupted, and the thirty lowest-in-degree pairs were selected by a mechanical rule fixed before any inference. Each query asks what the overruled decision held, without naming or hinting at the overruling authority (lexical leakage audited against the overruling holding and sentence). Four arms ran with identical generation settings: Direct LLM; Flat Vector RAG with the identical instruction block; Full TERRA; and TERRA-Unlabeled, which removes the relation tags, the overruling heading, the prompt's reference to the tag and the enlarged excerpt allocation for overruling edges while keeping seeds, traversal and the shared router decision unchanged. Two independent judges (Qwen3-Next-80B and GPT-OSS-120B; neither is the generator) classified each answer as refusal, verified refutation, fatal error (affirmed as current law, or historical recitation without warning) or hedge, and the combined class takes the more severe judgment. The primary hypothesis, the sensitivity subsets and the cross-suite contrast were preregistered before any live output. For the cross-suite contrast we re-ran the nineteen doctrine-invalidation queries of Category~D under identical code, prompts and judges.

\subsection{Comparative Baselines}
\label{subsec:baselines}
TERRA is benchmarked against four representative baseline architectures spanning parametric and retrieval-augmented paradigms: (1)~\textbf{Direct LLM (No RAG)}, where query $q$ is passed directly to the generator without retrieval, evaluating raw parametric knowledge and un-gated hallucination risks; (2)~\textbf{Flat Vector RAG}, where ChromaDB retrieves top-$k$ reasoning traces by semantic embedding similarity ($n_{\text{results}}=2$) and passes them directly to the generator without citation graph traversal or NLI entailment grading; (3)~\textbf{GraphRAG-Adapted (Louvain implementation)}, adapting Microsoft GraphRAG \cite{edge2024graphrag} by partitioning opinion entities into Louvain graph communities ($Q=\IssueTwentyThreeGraphRAGModularity$) and generating responses from local subgraph plus parent community summaries, omitting directed chronological precedent edges; and (4)~\textbf{HippoRAG-Adapted}, adapting hippocampal memory indexing \cite{bernal2024hipporag} by constructing a symmetric bipartite citation graph (\IssueTwentyThreeHippoRAGNonZeroEntries\ non-zero entries) and executing Personalized PageRank (PPR) associative diffusion seeded from query entity mentions, without typed directional citation constraints or downstream NLI self-correction gates.

\input{tables/table_baseline_specifications}

\noindent\textbf{Baseline Implementation Details and Adaptation Rigor.} All five pipelines share identical infrastructure: generator \texttt{databricks-meta-llama-3-3-70b-instruct} ($T=0.0$, max tokens~4,096), programmatic payload assertions (identical model, temperature, token limit), and benchmark \texttt{terra\_350\_benchmark.json} (350 queries, SHA256: \texttt{8144fa81...}) on the 20,000-case Harvard CAP corpus. The historical 1,750-run baseline benchmark was anchored to historical index \texttt{terra\_eeg\_index\_20k.json} (SHA256: \texttt{\IssueTwentyThreeHistoricalCorpusHashShort}, \IssueTwentyThreeHistoricalEdges\ edges, \IssueTwentyThreeHistoricalFileSizeMB~MB), while the active canonical index post-Issue~\#6 is SHA256 \texttt{\IssueTwentyThreeActiveCorpusHashShort} (\IssueTwentyThreeActiveEdges\ edges, \IssueTwentyThreeActiveFileSizeMB~MB), satisfying the closed accounting equation $\IssueTwentyThreeHistoricalEdges - \IssueTwentyThreeReclassifiedEdges + \IssueTwentyThreeCrsAddedEdges = \IssueTwentyThreeActiveEdges$ (\IssueTwentyThreeActiveCites\ \texttt{CITES}, \IssueTwentyThreeActiveOverrules\ \texttt{OVERRULES}; 53 reclassified, 3 CRS-added edges). Across the 20,000 opinions, \IssueTwentyThreeSubstantiveHoldingsCount\ contain substantive extracted holdings (\IssueTwentyThreeSubstantiveHoldingsPct\ coverage: \IssueTwentyThreeTierOneHeadnotesCount\ syllabus headnotes, \IssueTwentyThreeTierTwoRegexCount\ regex statements), while \IssueTwentyThreeTierThreeFallbackCount\ (\IssueTwentyThreeTierThreeFallbackPct) rely on preamble fallbacks, yielding \IssueTwentyThreePopulatedHoldingCount\ non-empty holding fields (\IssueTwentyThreePopulatedHoldingPct; \IssueTwentyThreeUnpopulatedHoldingCount\ unpopulated). Graph representations diverge by architecture (Table~\ref{tab:baseline_specifications}): GraphRAG-Adapted partitions an undirected citation network (\IssueTwentyThreeGraphRAGUndirectedEdges\ edges) via NetworkX Louvain modularity ($Q=\IssueTwentyThreeGraphRAGModularity$, resolution=1.0, seed=42; chosen over C-based Leiden due to host runtime C-extension build constraints); HippoRAG-Adapted executes Personalized PageRank ($\alpha=0.85$, tolerance $10^{-6}$, max 50 iterations) over a symmetric citation matrix (\IssueTwentyThreeHippoRAGNonZeroEntries\ entries, top-$k=5$); and TERRA traverses the directed chronological graph via BFS ($k=2$, depth $d=2$, hub cap~50, $d_{L2} \le 1.3$). All context ceilings in code are character-based string operations (\texttt{[:N]}): Flat RAG ($\approx$\IssueTwentyThreeFlatCharCeiling\ chars), GraphRAG ($\approx$\IssueTwentyThreeGraphRAGCharCeiling\ chars), HippoRAG ($\approx$\IssueTwentyThreeHippoRAGCharCeiling\ chars), whereas TERRA enforces a \IssueTwentyThreeTerraCharCeiling-char ceiling, representing an intentional $\IssueTwentyThreeContextAsymmetryRatio$ context volume asymmetry (\IssueTwentyThreePromptTokenRatio\ prompt tokens, Table~XVI) required for multi-decade precedent chains. The critic judge applies a \IssueTwentyThreeJudgeMiddleOmissionCeiling-char middle-omission split (\texttt{[:25000]} + \texttt{[-25000:]}), evaluating 100\% of baseline context ($<6{,}000$ chars) without truncation while preserving TERRA's initial seeds and terminal citing authorities. Finally, baseline adapters incorporate post-hoc amendments rather than frozen pre-registration: Amendments~F and G (HippoRAG regex and GraphRAG formatting fixes) and Amendments~K, L, and M (judge JSON parse clamping, refusal string harmonization, and fast-path ground-truth floor).

\subsection{Evaluation Metrics}
\label{subsec:eval-metrics}
System performance is evaluated across five quantitative dimensions:
\begin{itemize}
    \item \textbf{Faithfulness ($0.0$--$1.0$):} Evaluated via an LLM-as-a-judge framework \cite{roig2026scalablereliableevaluationai} using \texttt{databricks-meta-llama-3-3-70b-instruct} operating with an adversarial ``skeptical critic'' persona. The judge model is explicitly decoupled from the inference models to eliminate circular self-evaluation bias.
    \item \textbf{Relevance ($0.0$--$1.0$):} Scored by the LLM judge to measure query-answer directness and structural completeness.
    \item \textbf{ROUGE-L F1 and Recall (Supplementary):} Computed for Categories A and B against researcher-constructed ground-truth reference answers via deterministic tokenized LCS matching, providing an LLM-independent supplementary lexical alignment metric. Ground-truth doctrinal claims and reference holdings were cross-referenced against standard public legal reference sources (Cornell LII, Oyez, CRS \textit{Constitution Annotated}) without independent legal certification. ROUGE-L is strictly supplementary: exact n-gram matching cannot distinguish active doctrine from overruled precedent, heavily penalizes substantive paraphrasing and detailed analyses via candidate length, and scores defensive refusals as zero. The modest judge--ROUGE correlation ($r = \JudgePearsonR$; Section~\ref{subsec:validation-methodology}) confirms that lexical matching captures only a fraction of legal semantic correctness; primary conclusions are grounded in dual-judge Faithfulness, citation graph recovery, and adversarial safety rates.
    \item \textbf{Safety Rejection and Multi-Class Robustness Rate (\%):} Percentage of out-of-domain (Category C) and adversarial deception (Category D) queries correctly neutralized. For Category D, metrics are formally decoupled into: Explicit Safety Refusal Rate ($R_{\text{ref}}$, declining to affirm the invalid prompt), Verified Substantive Refutation Rate ($R_{\text{refut}}$, actively refuting the false premise without fatal error), Total Safe Handling Rate ($R_{\text{safe}} = R_{\text{ref}} \cup R_{\text{refut}}$), Fatal Error Rate ($R_{\text{fatal}}$, credulously validating false or overruled doctrine), and Uncommitted Hedge Rate ($R_{\text{hedge}}$, equivocal non-refusal non-answers).
    \item \textbf{Stage Latency (ms):} Mean $\pm$ standard deviation execution time for routing, retrieval, NLI grading, and generation stages, recorded across all benchmark runs.
\end{itemize}

\subsection{Validation Methodology and Component Ablation Design}
\label{subsec:validation-methodology}
To satisfy the rigorous methodological standards required for temporal judicial AI research, the evaluation protocol incorporates three validation controls:
\begin{enumerate}
    \item \textbf{Judge Calibration and Dual-Metric Alignment:} To verify that the LLM judge does not produce uncalibrated scores, we measure the correlation between the judge's Faithfulness ratings and the independent, researcher-constructed reference ROUGE-L scores across in-domain generations (Categories~A and B). To account for repeated evaluations within queries without discarding pipeline variation, the primary analysis employs whole-query cluster bootstrap and query-cluster permutation across $N_q = \JudgeCalibN$ in-domain query clusters ($N_{\text{obs}} = \JudgeCalibNObs$ valid evaluation pairs across available pipeline evaluations per query). Under this cluster-aware specification, the primary association is Pearson $r = \JudgePearsonR$ (95\% whole-query bootstrap CI $\JudgePearsonCI$; two-sided query-cluster permutation $p = \frac{k+1}{P+1} = \JudgePearsonP$ based on $P = 10{,}000$ permutations with $k = 0$ exceedances; Spearman $\rho = \JudgeSpearmanRho, p = \JudgeSpearmanP$; Kendall $\tau = \JudgeKendallTau, p = \JudgeKendallP$). Collapsing available pipeline evaluations into query-level means alters the estimand to purely between-query variation, yielding secondary query-mean correlation $r_{\text{mean}} = \JudgeMeanPearsonR$ (Fisher 95\% CI $\JudgeMeanPearsonCI$, permutation $p = \JudgeMeanPearsonP$). Historical nominal calibration figures ($r = \HistoricalJudgePearsonR$, pilot $r = \HistoricalPilotJudgePearsonR$) were hardcoded without machine-readable provenance and failed systematic reproduction across 384 candidate preprocessing configurations; the reported estimates reflect the canonical corrected analysis selected after failed reproduction. On the preliminary 35-query pilot subset ($N=525$ total runs across 3 pipelines and 5 repeated runs; $N=\numPilotBaselineEvals$ represents evaluations for a single pipeline; $N=260$ valid in-domain pairs across 20 query clusters), correlation was $r = \PilotJudgePearsonR$ ($p = \PilotJudgePearsonP$). These results indicate a modest but statistically significant positive relationship between automated critic scoring and lexical grounding, though not strong enough to treat the LLM judge as a full substitute for independent human legal verification while reflecting the known divergence between exact n-gram matching and conceptual entailment.
    \item \textbf{Multi-Run Provenance and Execution Conditions:}
Table~\ref{tab:run_provenance} summarises the provenance of every evaluation
block reported in this paper, including run dates (from file timestamps or
logged headers), generation models, judge models, and sample sizes.
The baselines were executed on 5--6~September~2026; the deployed TERRA run
was executed on 14~September~2026 with four concurrent workers (36 of 350
queries incurred HTTP~429 rate-limiting, accumulating 620\,s of
server-imposed sleep). Dual-judge faithfulness scoring used
\texttt{databricks-qwen3-next-80b-a3b-instruct} (Judge~1) and
\texttt{databricks-gpt-oss-120b} (Judge~2) across all five pipelines.
All inference calls were issued at temperature $T{=}0.0$; however, this
does not guarantee identical outputs on Databricks Model Serving endpoints.
A second TERRA run performed on the same day as the primary run yielded
identical answers on $\TerraIdenticalAnswerFraction$ of queries overall ($\TerraIdenticalAnswerCatA$ on Category~A,
$\TerraIdenticalAnswerCatB$ on Category~B, $\TerraIdenticalAnswerCatC$ on Category~C, $\TerraIdenticalAnswerCatD$ on Category~D);
the low Category~B identity fraction primarily reflects that the earlier run
predates the standalone router update; for the 241 queries whose route was
forced to HARD in that run, different execution paths produced different answers.
Evaluation blocks not re-executed under the v2 dual-judge protocol
(pilot ablation, leakage probe, CaseFacts, CLERC) are identified by
``date not recorded'' in Table~\ref{tab:run_provenance}; their absolute
scores are not comparable with dual-judge scores.
    \item \textbf{Component-Level Ablation Matrix:} To isolate the specific performance drivers of the architecture, we formulate three same-architecture ablations alongside the full pipeline: \textbf{Full TERRA} (Traffic Cop Router + Multi-Hop EEG Traversal + NLI Smart Grader); \textbf{Ablation 1 (w/o EEG Traversal)} which replaces graph traversal with standard Flat Vector RAG while retaining NLI grading, isolating the topological contribution; \textbf{Ablation 2 (w/o NLI Smart Grader)} which disables the entailment verification and self-correction loop, generating directly from raw parametric and retrieved text to isolate firewall efficacy; and \textbf{Ablation 3 (w/o Traffic Cop Router)} which forces all queries (including simple single-point factual lookups) through the full GraphRAG traversal path to quantify routing efficiency gains.
    \item \textbf{Hyperparameter Freezing and Benchmark Isolation:} Core retrieval and graph parameters ($k=\IssueTwentyFourSeedK$, depth $d=\IssueTwentyFourBfsDepth$, hub cap $\IssueTwentyFourHubCap$, $\tau=\IssueTwentyFourDistThreshold$) were frozen before benchmark construction and evaluation based on topological network constraints and feasibility testing on an independent $N=\IssueTwentyFourDevQueryCount$-query pilot (\IssueTwentyFourDevDate; \IssueTwentyFourDevTotalEvals\ total run records across \IssueTwentyFourDevPipelineCount\ pipelines and \IssueTwentyFourDevRunCount\ stability runs). An automated audit verifies chronological precedence, zero query ID and prompt-text overlap (\IssueTwentyFourOverlapPct), and zero query-text or reference-answer contamination in ChromaDB across all \IssueTwentyFourTestQueryCount~benchmark queries.
\end{enumerate}

\begin{table}[t]
  \centering
  \small
  \caption{Evaluation Block Provenance. Dates are from file timestamps or
           logged runner headers; ``date not recorded'' means no timestamp
           was found in the source files.}
  \label{tab:run_provenance}
  \resizebox{\columnwidth}{!}{%
  \begin{tabular}{@{}p{3.8cm}p{2.2cm}p{5.2cm}p{1.4cm}@{}}
    \toprule
    \textbf{Block} & \textbf{Date (source)} & \textbf{Models} & \textbf{$N$} \\
    \midrule
    Baseline generations & 5--6 Sept.\ 2026\newline (file mtime) &
      Gen: \texttt{llama-3.3-70b-instruct}; no router/NLI &
      1,400 \\
    \addlinespace
    Deployed TERRA run & 14 Sept.\ 2026\newline (file mtime) &
      Router: \texttt{gemma-3-12b}; NLI + gen: \texttt{llama-3.3-70b-instruct} &
      350 \\
    \addlinespace
    Dual-judge scoring & 14 Sept.\ 2026\newline (file mtime) &
      J1: \texttt{qwen3-next-80b}; J2: \texttt{gpt-oss-120b} &
      1,750 \\
    \addlinespace
    Pilot ablation (Table~\ref{tab:ablation}) & date not recorded &
      date not recorded & 35 \\
    \addlinespace
    Temporal leakage probe & 17 Sept.\ 2026\newline (file mtime) &
      date not recorded in JSON headers & varies \\
    \addlinespace
    Category~D certification & 20 Sept.\ 2026\newline (file last modified) &
      Semantic verifier: \texttt{llama-3.3-70b-instruct} & 75 \\
    \addlinespace
    CaseFacts eval & 18 Sept.\ 2026\newline (file mtime) &
      date not recorded in JSON headers & 170 \\
    \addlinespace
    CLERC eval & 18 Sept.\ 2026\newline (file mtime) &
      date not recorded in JSON headers & 170 \\
    \addlinespace
    Obscure suite & 19--20 Sept.\ 2026\newline (file mtime) &
      Gen: \texttt{llama-3.3-70b-instruct}; J1/J2 as above &
      120 \\
    \addlinespace
    Famous D1 contrast & 19--20 Sept.\ 2026\newline (file mtime) &
      Gen: \texttt{llama-3.3-70b-instruct}; J1/J2 as above &
      76 \\
    \bottomrule
  \end{tabular}%
  }
\end{table}

\enlargethispage{6\baselineskip}
\subsection{Data and Code Availability}
\label{subsec:data-code-availability}
\label{subsec:reproducibility}
To enable independent replication and adhere to open-science standards, all benchmark data, experimental code, and model specifications are publicly released under an open-source license via our public repository at \texttt{https://github.com/zaforsaadik7/TERRA-Temporal-Evolution-and-Reasoning-Trace-RAG}, with permanent archival release snapshots deposited on Zenodo upon publication. The public reproducibility package includes: (1)~the complete 350-query benchmark suite (\texttt{terra\_350\_benchmark.json}, SHA256: \texttt{8144fa81...}) with gold categories, doctrine labels, and reference answers; (2)~the 19 verified doctrine-invalidation pairs and 47 obscure overruled pairs; (3)~the 20,000-case Event Evolution Graph index (\texttt{terra\_eeg\_index\_20k.json}, SHA256: \texttt{\IssueTwentyThreeActiveCorpusHashShort}, \IssueTwentyThreeActiveEdges\ directed edges: \IssueTwentyThreeActiveCites\ \texttt{CITES} and \IssueTwentyThreeActiveOverrules\ authoritative \texttt{OVERRULES}), containing all 20,000 procedural reasoning traces, alongside the graph construction pipeline (\texttt{scripts/ingest\_20k\_eeg\_and\_vector\_db.py}) and standalone citation edge list (\texttt{data/terra\_eeg\_edges\_145k.tsv}); (4)~complete execution and evaluation scripts (\texttt{ask\_terra.py}, \texttt{eval\_terra.py}); (5)~verbatim system prompt templates and structured Pydantic schemas; (6)~system configurations and frozen hyperparameters (\texttt{configs/terra\_config.json}); (7)~canonical dual-judge evaluation records ($N=\numTotalEvals$ evaluations, \texttt{data/terra\_canonical\_v2\_eval.json}, SHA256: \texttt{95e775e2...}) alongside specialized benchmark logs (CaseFacts, CLERC, Obscure Overruled, and NLI verification suites in \texttt{data/}); and (8)~automated table scripts (\texttt{Paper/scripts/generate\_stats.py}) regenerating every table (Tables~I--XXIII) and macro constant. Large raw corpora (unparsed text and vector stores) are reconstructible via the ingestion pipeline. \textbf{Component Model Specifications:} Generative inference and the NLI Smart Grader employ \texttt{meta-llama/Llama-3.3-70B-Instruct} (Databricks Model Serving API: \texttt{databricks-meta-llama-3-3-70b-instruct}, $T=0.0$, max tokens~4,096). The Traffic Cop intent router uses Google Gemma 3 12B Instruction-tuned (API: \texttt{\RouterModel}, $T=0.0$, Pydantic schema validation; Table~\ref{tab:router_eval}). Dense semantic retrieval uses \texttt{sentence-transformers/all-MiniLM-L6-v2} (\TraceEmbeddingDimensions-dimensional vectors, Euclidean $L_2$ metric, ChromaDB PersistentClient, top-$k=2$, $d_{L2} \le 1.3$). Independent critic scoring utilizes \texttt{databricks-qwen3-next-80b-a3b-instruct} (Judge~1) and \texttt{databricks-gpt-oss-120b} (Judge~2) under the frozen rubric (Table~\ref{tab:run_provenance}). Local retrieval executed on Windows 11 Enterprise (16 cores, 64~GB RAM, ChromaDB v\TraceChromaVersion, NetworkX v\TraceNetworkXVersion, Python 3.11); remote inference on NVIDIA A100-SXM4-80GB GPU clusters.

\section{Results and Analysis}
\label{sec:results-analysis}

To rigorously evaluate TERRA GraphRAG, we adopt an ablation-first empirical methodology. Rather than presenting aggregated metrics that obscure internal component dynamics, we first dissect the architecture into its constituent mechanisms to isolate the empirical contribution of each subsystem. We then evaluate multi-hop evolutionary reasoning via pairwise win rates and semantic faithfulness, examine factual precision and safety firewalling across the complete \numTotalEvals-evaluation benchmark, profile stage-level latency distributions, and conduct formal statistical significance and power audits.

\subsection{Architectural Component Ablation and System Anatomy}
\label{subsec:ablation-study}

To quantify the direct empirical contribution of each architectural component, Table~\ref{tab:ablation} reports the component-level ablation matrix evaluated across the foundational \numPilotQueries-query benchmark suite ($N=\numPilotRepeatedRuns$ repeated runs, \numPilotEvals\ evaluations). We note that this ablation study evaluates component degradation over repeated runs on the focused 35-query pilot subset (reporting Full TERRA Category B faithfulness as $\AblationFullCatBFaithfulness$ and Category D safety as $\AblationFullCatDSafetyRate$ on $n=\IssueTwentyFourDevCatDCount$ unique adversarial traps, 25 run observations). This configuration isolates architectural drops under repeated sampling prior to full scaling, whereas the primary $\numUniqueQueries$-query benchmark (Sections~\ref{subsec:evolutionary-reasoning} and \ref{subsec:safety-and-factual}, Tables~\ref{tab:cat_b_performance} and \ref{tab:overall_summary}) evaluates the expanded 100-query Category B suite ($\TerraCatBFaithfulness$) and broader 75-query Category D suite in a single-run protocol ($T=0.0$), where TERRA achieves an explicit refusal count of $\TerraCatDRejection$ alongside $\TerraCatDRefutationRate$ verified refutations ($\TerraCatDRefutations$), for $\TerraCatDSafeHandlingRate$ total safe handling ($\TerraCatDSafeHandling$); the two fatal errors ($\TerraCatDFatalErrors$, D033 and D034) are procedural misattributions (Table~\ref{tab:refusal_breakdown}). Category~D labels come from a calibrated semantic verifier served by the same model family as the generator (Limitation~\ref{item:samefamily}).

\begin{table}[htbp]
\caption{Component-Level Ablation Analysis Across \numPilotQueries\ Benchmark Queries ($N=\numPilotRepeatedRuns$ Repeated Runs, \numPilotEvals\ Evaluations).}
\label{tab:ablation}
\centering
\resizebox{0.98\columnwidth}{!}{%
\begin{tabular}{lcccc}
\toprule
\textbf{Architecture Configuration} & \textbf{Cat B Faith.} & \textbf{Cat D Safety} & \textbf{Cat A False Ref.} & \textbf{Latency (Cat A)} \\
\midrule
\textbf{Full TERRA (Router + EEG + NLI)} & $\mathbf{\AblationFullCatBFaithfulness}$ & $\mathbf{\AblationFullCatDSafetyRate}$ & $\mathbf{\AblationFullCatAFalseRefusalCount}$ & $\mathbf{\AblationFullLatencyMedianSeconds\text{ s}}$ \\
Ablation 1: w/o EEG Traversal (Vector Traces Only) & $\AblationOneCatBFaithfulness$ & $\AblationOneCatDSafetyRate$ & $\AblationOneCatAFalseRefusalCount$ & $\AblationOneLatencyMedianSeconds\text{ s}$ \\
Ablation 2: w/o NLI Smart Grader (No Verification Gate) & $\AblationTwoCatBFaithfulness$ & $\AblationTwoCatDSafetyRate$ & $\AblationTwoCatAFalseRefusalCount$ & $\AblationTwoLatencyMedianSeconds\text{ s}$ \\
Ablation 3: w/o Traffic Cop Router (Forced Hard Path) & $\AblationThreeCatBFaithfulness$ & $\mathbf{\AblationThreeCatDSafetyRate}$ & $\AblationThreeCatAFalseRefusalCount$ ($\AblationThreeCatAFalseRefusalRate$) & $\AblationThreeLatencyMedianSeconds\text{ s}$ \\
\bottomrule
\end{tabular}%
}
\end{table}

The ablation results confirm three fundamental architectural insights:
\begin{enumerate}
    \item \textbf{Architectural Necessity of EEG Graph Traversal:} Disabling multi-hop citation graph traversal (Ablation 1) drops evolutionary reasoning faithfulness from $\AblationFullCatBFaithfulness$ to $\AblationOneCatBFaithfulness$ (a relative reduction of $\AblationEEGFaithfulnessRelDrop$). Semantic vector retrieval over isolated reasoning traces fails to capture the chronological directionality of legal precedents; when retrieved text segments lack explicit bridge citations connecting non-contiguous decisions, generative synthesis degrades.
    \item \textbf{Critical Function of the NLI Smart Grader:} Disabling the semantic entailment gate (Ablation 2) removes the system's defensive verification firewall, dropping adversarial safety rejection on Category D from $\AblationFullCatDSafetyRate$ to $\AblationTwoCatDSafetyRate$. Without the autonomous entailment audit, the generator credulously complies with deceptive queries regarding overruled precedent.
    \item \textbf{Operational Efficiency of Traffic Cop Routing:} Bypassing the SLM complexity router (Ablation 3) forces straightforward single-point factual lookups through heavy multi-hop BFS graph expansion and NLI grading, increasing execution latency on Category A queries in the pilot from $\AblationFullLatencyMedianSeconds\text{ s}$ to $\AblationThreeLatencyMedianSeconds\text{ s}$ ($\AblationRouterLatencySlowdownFactor\times$ slowdown) while introducing $\AblationThreeCatAFalseRefusalCount$ ($\AblationThreeCatAFalseRefusalRate$) false refusals due to context pollution from extraneous citing precedents. Notably, Ablation 3 reports a marginally higher Category B faithfulness ($\AblationThreeCatBFaithfulness$) than Full TERRA ($\AblationFullCatBFaithfulness$): because the forced hard path enforces BFS graph traversal for every query without exception, it eliminates any residual router misclassification of borderline evolutionary queries, yielding a minor uplift on multi-hop reasoning. However, the operational cost---a $\AblationRouterLatencySlowdownFactor\times$ latency increase and a $\AblationThreeCatAFalseRefusalRate$ false refusal rate on Category A factual queries---makes the unrouted configuration disadvantageous for latency-sensitive applications. In this pilot ablation the router provides operational triage: the full configuration's median on direct factual inquiries is $\AblationFullLatencyMedianSeconds\text{ s}$, whereas the fast route's median in the 350-query canonical run is \TerraFastRouteMedianS\ (the two populations differ in size, date and endpoint load, so the figures are not directly comparable), and graph traversal is reserved for complex multi-decade doctrines.
\end{enumerate}

\subsection{Multi-Hop Evolutionary Reasoning and Pairwise Win Rates (Category B)}
\label{subsec:evolutionary-reasoning}

Category B queries evaluate the primary challenge in judicial AI: tracking constitutional doctrine across multi-decade evolutionary precedent arcs (e.g., the progression from \textit{Plessy} (1896), through \textit{Gaines} (1938) and \textit{Sweatt} (1950), to \textit{Brown} (1954)). Table~\ref{tab:cat_b_performance} summarizes comparative performance and pairwise head-to-head win rates across $\CatBTotalPairs$ paired evaluation runs.

\begin{table}[htbp]
\caption{Category B Multi-Hop Evolutionary Performance and Pairwise Win Rates ($N=\CatBTotalPairs$ Paired Evaluations).}
\label{tab:cat_b_performance}
\centering
\resizebox{0.98\columnwidth}{!}{%
\begin{tabular}{lcccc}
\toprule
\textbf{Pipeline} & \textbf{Faithfulness} & \textbf{TERRA Win Rate} & \textbf{ROUGE-L Rec.} & \textbf{Mean Words} \\
\midrule
Direct LLM (No RAG) & $\DirectLLMCatBFaithfulness \pm \DirectLLMCatBFaithfulnessStd$ & $\TerraVsDirectLLMCatBWinRate^\dagger$ & $\mathbf{\DirectLLMCatBRougeRecall}$ & $\mathbf{\DirectLLMCatBMeanWords}$ \\
Flat Vector RAG & $\FlatRAGCatBFaithfulness \pm \FlatRAGCatBFaithfulnessStd$ & $\mathbf{\TerraVsFlatRAGCatBWinRate}$ & $\FlatRAGCatBRougeRecall$ & $\FlatRAGCatBMeanWords$ \\
GraphRAG-Adapted & $\GraphRAGCatBFaithfulness \pm \GraphRAGCatBFaithfulnessStd$ & $\mathbf{\TerraVsGraphRAGCatBWinRate}$ & $\GraphRAGCatBRougeRecall$ & $\GraphRAGCatBMeanWords$ \\
HippoRAG-Adapted & $\HippoRAGCatBFaithfulness \pm \HippoRAGCatBFaithfulnessStd$ & $\mathbf{\TerraVsHippoRAGCatBWinRate}$ & $\HippoRAGCatBRougeRecall$ & $\HippoRAGCatBMeanWords$ \\
\textbf{TERRA GraphRAG} & $\mathbf{\TerraCatBFaithfulness \pm \TerraCatBFaithfulnessStd}$ & --- & $\TerraCatBRougeRecall$ & $\TerraCatBMeanWords$ \\
\bottomrule
\end{tabular}%
}
\vspace{-2mm}
\raggedright \footnotesize \textbf{Notes:} Bold denotes top performance among retrieval-augmented (RAG) pipelines; Direct LLM (No RAG) serves as an ungrounded parametric baseline reference. $^\dagger$Against Direct LLM, head-to-head outcomes comprise $\TerraVsDirectLLMCatBWins$ TERRA wins, $\TerraVsDirectLLMCatBTies$ ties, and $\TerraVsDirectLLMCatBLosses$ losses, reflecting parametric fluency on generic summaries versus grounded citation depth.
\end{table}

On Category~B, TERRA's faithfulness exceeds Flat Vector RAG by $\CatBTerraVsFlatRAGDelta$ (paired Wilcoxon $q=\CatBTerraVsFlatRAGAdjQ$; the bootstrap CI of the mean difference, [$\CatBTerraVsFlatRAGBootCILow$, $\CatBTerraVsFlatRAGBootCIHigh$], includes zero while the Hodges--Lehmann median difference, $\CatBTerraVsFlatRAGHLMedian$ [$\CatBTerraVsFlatRAGHLLow$, $\CatBTerraVsFlatRAGHLHigh$], does not), HippoRAG by $\CatBTerraVsHippoRAGDelta$ ($q=\CatBTerraVsHippoRAGAdjQ$, $d_z=\CatBTerraVsHippoRAGDz$) and GraphRAG by $\CatBTerraVsGraphRAGDelta$ ($q=\CatBTerraVsGraphRAGAdjQ$). Head-to-head, TERRA wins $\TerraVsFlatRAGCatBWins$, $\TerraVsGraphRAGCatBWins$ and $\TerraVsHippoRAGCatBWins$ of 100 queries against Flat, GraphRAG and HippoRAG, with $\TerraVsFlatRAGCatBTies$, $\TerraVsGraphRAGCatBTies$ and $\TerraVsHippoRAGCatBTies$ ties. A no-retrieval Direct LLM scores higher on Category~B ($\DirectLLMCatBFaithfulness$ vs.\ $\TerraCatBFaithfulness$; $\Delta=\CatBTerraVsDirectLLMDelta$, $q=\CatBTerraVsDirectLLMAdjQ$; TERRA wins $\TerraVsDirectLLMCatBWins$, ties $\TerraVsDirectLLMCatBTies$, loses $\TerraVsDirectLLMCatBLosses$), consistent with dual judges rewarding fluent parametric summaries; TERRA's advantage is grounded citation provenance (Section~\ref{subsec:citation-correctness}) and retrieval coverage, not raw judged faithfulness.

\noindent\textbf{Contextualizing the Direct LLM Comparison.} Against Direct LLM, TERRA is significantly better on Category~A ($\CatATerraVsDirectLLMDelta$, $q=\CatATerraVsDirectLLMAdjQ$) and on the combined in-domain slice ($\CombinedTerraVsDirectLLMDelta$ [$\CombinedTerraVsDirectLLMBootCILow$, $\CombinedTerraVsDirectLLMBootCIHigh$], $q=\CombinedTerraVsDirectLLMAdjQ$), significantly worse on Category~B ($\CatBTerraVsDirectLLMDelta$, $q=\CatBTerraVsDirectLLMAdjQ$) and, after certified-safety clamping, on Category~D ($\CatDTerraVsDirectLLMDelta$, $q=\CatDTerraVsDirectLLMAdjQ$). The combined result is driven by Category~A, where Direct LLM's ungrounded answers are scored down by the non-Llama judges (its Category~A score is $\DirectLLMCatAFaithfulness$, versus $\DirectLLMCatAFaithfulnessLegacy$ under the earlier single Llama judge).

\subsubsection{ROUGE-L Lexical Trade-off and Recall Disentanglement}
Standard lexical metrics such as ROUGE-L F1 enforce a severe length penalty via candidate-length precision in the denominator ($\text{Precision} = \text{LCS}/|\text{cand}|$). Table~\ref{tab:rouge_tradeoff} details the lexical precision-recall trade-off and mean word counts across Category B.

\begin{table}[htbp]
\caption{Category B Lexical Alignment and Answer Verbosity Trade-off.}
\label{tab:rouge_tradeoff}
\centering
\resizebox{0.98\columnwidth}{!}{%
\begin{tabular}{lcccc}
\toprule
\textbf{Pipeline} & \textbf{ROUGE-L P} & \textbf{ROUGE-L R} & \textbf{ROUGE-L F1} & \textbf{Mean Words} \\
\midrule
Human Ground Truth & $1.000$ & $1.000$ & $1.000$ & $70.8$ \\
Direct LLM (No RAG) & $\DirectLLMCatBRougePrecision$ & $\mathbf{\DirectLLMCatBRougeRecall}$ & $\DirectLLMCatBRougeFOne$ & $\DirectLLMCatBMeanWords$ \\
Flat Vector RAG & $\mathbf{\FlatRAGCatBRougePrecision}$ & $\FlatRAGCatBRougeRecall$ & $\mathbf{\FlatRAGCatBRougeFOne}$ & $\FlatRAGCatBMeanWords$ \\
GraphRAG-Adapted & $\GraphRAGCatBRougePrecision$ & $\GraphRAGCatBRougeRecall$ & $\GraphRAGCatBRougeFOne$ & $\GraphRAGCatBMeanWords$ \\
HippoRAG-Adapted & $\HippoRAGCatBRougePrecision$ & $\HippoRAGCatBRougeRecall$ & $\HippoRAGCatBRougeFOne$ & $\HippoRAGCatBMeanWords$ \\
\textbf{TERRA GraphRAG} & $\TerraCatBRougePrecision$ & $\TerraCatBRougeRecall$ & $\TerraCatBRougeFOne$ & $\TerraCatBMeanWords$ \\
\bottomrule
\end{tabular}%
}
\vspace{1mm}
\raggedright \footnotesize \textit{Note (computation):} F1 values are the mean of per-query F1 scores computed individually per response--reference pair; they are not the harmonic mean of the macro P and R column averages and will not reproduce exactly from those column values due to the nonlinearity of the harmonic mean. \textit{Note (reference set):} All system P/R/F1 values were computed against researcher-constructed reference answers averaging $70.8$ words via deterministic tokenized LCS matching; the canonical reference set is \texttt{terra\_350\_benchmark.json} (SHA256: \texttt{8144fa81...}). The Human Ground Truth row (P/R/F1 = 1.000) reflects each reference evaluated against itself by definition.
\end{table}

Crucially, evaluating lexical alignment exposes the inadequacy of surface n-gram metrics like ROUGE-L for judicial QA. On standard ROUGE-L F1, Flat Vector RAG scores higher ($\mathbf{\FlatRAGCatBRougeFOne}$) than TERRA ($\TerraCatBRougeFOne$), with HippoRAG at $\HippoRAGCatBRougeFOne$ and GraphRAG at $\GraphRAGCatBRougeFOne$. This divergence reflects two factors: (1)~Because reference holdings average $70.8$ words, candidate-length precision ($\text{Precision} = \text{LCS}/|\text{cand}|$) heavily penalizes detailed legal syntheses ($\TerraCatBRougePrecision$ for TERRA vs.\ $\mathbf{\FlatRAGCatBRougePrecision}$ for Flat RAG), where TERRA provides structured chronological trajectories rather than brief summaries; and (2)~On benchmark-wide \textbf{ROUGE-L Recall}, TERRA's conservative firewall falsely refuses \TerraCatBFalseRefusals\ Category~B queries (scored as zero overlap), yielding an overall recall of $\TerraCatBRougeRecall$, compared to $\FlatRAGCatBRougeRecall$ for Flat RAG, $\GraphRAGCatBRougeRecall$ for GraphRAG, and $\HippoRAGCatBRougeRecall$ for HippoRAG (and $\mathbf{\DirectLLMCatBRougeRecall}$ for Direct LLM). On answered queries, TERRA reaches \TerraCatBAnsweredRougeRecall\ (averaging \TerraCatBAnsweredMeanWords\ words), comparable to Flat RAG's answered-only \FlatRAGCatBAnsweredRougeRecall; its lower all-query value reflects 16 safety refusals scored as zero. Crucially, n-gram overlap cannot distinguish valid modern doctrine from overruled precedent: ROUGE-L rewards verbatim citation of dead law, whereas TERRA's core strengths lie in verified citation graph recovery (Table~\ref{tab:retrieval_eval}) and dual-judge semantic faithfulness ($\TerraCatBFaithfulness$, FDR $q < 0.05$).

\subsection{Factual Precision, Safety Firewalling, and Adversarial Deception}
\label{subsec:safety-and-factual}

Table~\ref{tab:overall_summary} presents the comprehensive comparative performance across all five pipelines and four benchmark categories over the canonical \numUniqueQueries-query benchmark ($N=\numTotalEvals$ evaluations across \numUniqueQueries\ unique queries $\times$ \numComparativePipelines\ comparative pipelines).

\begin{table*}[t]
\caption{Comprehensive Benchmark Performance Breakdown ($N=\numTotalEvals$ Evaluations Across \numUniqueQueries\ Unique Queries $\times$ \numComparativePipelines\ Comparative Pipelines, with 95\% Bootstrap CIs).}
\label{tab:overall_summary}
\centering
\resizebox{0.98\textwidth}{!}{%
\begin{tabular}{lccccccc}
\toprule
\textbf{Pipeline} & \textbf{Overall Faithfulness [95\% CI]} & \textbf{Answer-Only Faithfulness [95\% CI]} & \textbf{Cat A Factual} & \textbf{Cat B Evolutionary} & \textbf{Cat C Out-of-Context Refusal} & \textbf{Cat D Adversarial Refusal} & \textbf{Median Latency} \\
\midrule
Direct LLM (No RAG) & $\DirectLLMOverallFaithfulness$ $\DirectLLMOverallCI$ & $\DirectLLMAnswerOnlyFaithfulness$ $\DirectLLMAnswerOnlyCI$ & $\DirectLLMCatAFaithfulness$ & $\DirectLLMCatBFaithfulness$ & $\DirectLLMCatCRejection$ & $\DirectLLMCatDRejection$ & $\DirectLLMLatencyMedianSeconds\text{ s}$ \\
Flat Vector RAG & $\FlatRAGOverallFaithfulness$ $\FlatRAGOverallCI$ & $\FlatRAGAnswerOnlyFaithfulness$ $\FlatRAGAnswerOnlyCI$ & $\FlatRAGCatAFaithfulness$ & $\FlatRAGCatBFaithfulness$ & $\FlatRAGCatCRejection$ & $\FlatRAGCatDRejection$ & $\FlatRAGLatencyMedianSeconds\text{ s}$ \\
GraphRAG-Adapted & $\GraphRAGOverallFaithfulness$ $\GraphRAGOverallCI$ & $\GraphRAGAnswerOnlyFaithfulness$ $\GraphRAGAnswerOnlyCI$ & $\GraphRAGCatAFaithfulness$ & $\GraphRAGCatBFaithfulness$ & $\mathbf{\GraphRAGCatCRejection}$ & $\GraphRAGCatDRejection$ & $\GraphRAGLatencyMedianSeconds\text{ s}$ \\
HippoRAG-Adapted & $\mathbf{\HippoRAGOverallFaithfulness}$ $\mathbf{\HippoRAGOverallCI}$ & $\HippoRAGAnswerOnlyFaithfulness$ $\HippoRAGAnswerOnlyCI$ & $\mathbf{\HippoRAGCatAFaithfulness}$ & $\HippoRAGCatBFaithfulness$ & $\HippoRAGCatCRejection$ & $\HippoRAGCatDRejection$ & $\mathbf{\HippoRAGLatencyMedianSeconds\text{ s}}$ \\
\textbf{TERRA GraphRAG} & $\TerraOverallFaithfulness$ $\TerraOverallCI$ & $\mathbf{\TerraAnswerOnlyFaithfulness}$ $\mathbf{\TerraAnswerOnlyCI}$ & $\TerraCatAFaithfulness$ & $\mathbf{\TerraCatBFaithfulness}$ & $\mathbf{\TerraCatCRejection}$ & $\mathbf{\TerraCatDRejection}$ & $\TerraLatencyMedianSeconds\text{ s}$ \\
\bottomrule
\end{tabular}%
}
\vspace{-4mm}
\raggedright \scriptsize \textbf{Notes:} Bold denotes top performance among retrieval-augmented (RAG) pipelines; Direct LLM (No RAG) serves as an ungrounded parametric baseline reference.
\end{table*}

\input{tables/table_sensitivity_raw_vs_clamped}
\input{tables/table_refusal_breakdown}

\begin{figure}[htbp]
    \centering
    \resizebox{0.98\columnwidth}{!}{\input{figure_3_safety_rejection}}
    \caption{Safety refusal rate comparison across Category C (Out-of-Context) and Category D (Adversarial Deception). Bars show system-flagged refusals (Table~\ref{tab:refusal_breakdown}, Panel~A): TERRA flags 100\% of out-of-context and \TerraCatDFlaggedRefusalRate{} (\TerraCatDFlaggedRefusals) of adversarial queries; Direct LLM flags none.}
    \label{fig:safety_rejection}
\end{figure}

\enlargethispage{14\baselineskip}
Key empirical findings include:
\begin{itemize}
    \item \textbf{Dual-Metric Faithfulness Leadership:} Distinguishing overall faithfulness (awarding perfect scores to epistemic refusals) from answer-only faithfulness (evaluating generated legal arguments) is critical. TERRA attains the highest overall faithfulness ($\TerraOverallFaithfulnessMean$, 95\% CI $\TerraOverallCI$) ahead of HippoRAG-Adapted ($\HippoRAGOverallFaithfulnessMean$, $\HippoRAGOverallCI$), whose confidence intervals overlap, and the highest answer-only faithfulness ($\TerraAnswerOnlyFaithfulnessMean$ vs.\ $\HippoRAGAnswerOnlyFaithfulnessMean$ for HippoRAG); answer-only scores are conditional on different answered subsets (TERRA answers $\TerraAnsweredQueries$/350 queries, HippoRAG \HippoRAGAnswerOnlyCount, Flat \FlatRAGAnswerOnlyCount, GraphRAG \GraphRAGAnswerOnlyCount). Flagged refusals total $\TerraFlaggedRefusals$/350 for TERRA against $\HippoRAGFlaggedRefusals$ for HippoRAG, $\FlatRAGFlaggedRefusals$ for Flat and $\GraphRAGFlaggedRefusals$ for GraphRAG ($\TerraEpistemicRefusals$, $\HippoEpistemicRefusals$, $\FlatEpistemicRefusals$ and $\GraphEpistemicRefusals$ counting unflagged textual refusals).
    \item \textbf{Evolutionary Multi-Hop Performance (Category B):} On multi-decade precedent trajectories, TERRA achieves $\TerraCatBFaithfulness$ faithfulness, significantly outperforming all evaluated retrieval baselines: HippoRAG ($\HippoRAGCatBFaithfulness$, $\Delta=\CatBTerraVsHippoRAGDelta$, $q=\CatBTerraVsHippoRAGAdjQ$), GraphRAG ($\GraphRAGCatBFaithfulness$, $\Delta=\CatBTerraVsGraphRAGDelta$, $q=\CatBTerraVsGraphRAGAdjQ$), and Flat Vector RAG ($\FlatRAGCatBFaithfulness$, $\Delta=\CatBTerraVsFlatDelta$, $q=\CatBTerraVsFlatRAGAdjQ$). A non-retrieval Direct LLM scores higher on Category~B ($\DirectLLMCatBFaithfulness$, $\Delta=\CatBTerraVsDirectLLMDelta$, $q=\CatBTerraVsDirectLLMAdjQ$; TERRA win rate $\TerraVsDirectLLMCatBWinRate$), consistent with dual judges rewarding fluent ungrounded parametric prose over grounded retrieval constraints, whereas TERRA's retrieval evaluation shows stronger Category~B citation coverage and full-precedent-path recovery.
    \item \textbf{Factual Lookup Performance (Category A):} On Category~A, TERRA scores $\TerraCatAFaithfulness$ against HippoRAG $\HippoRAGCatAFaithfulness$, Flat $\FlatRAGCatAFaithfulness$ and Direct LLM $\DirectLLMCatAFaithfulness$; it is significantly higher than Direct ($\CatATerraVsDirectLLMDelta$) and Flat ($\CatATerraVsFlatRAGDelta$, $q=\CatATerraVsFlatRAGAdjQ$) and statistically indistinguishable from GraphRAG ($\CatATerraVsGraphRAGDelta$, $q=\CatATerraVsGraphRAGAdjQ$) and HippoRAG ($\CatATerraVsHippoRAGDelta$, $q=\CatATerraVsHippoRAGAdjQ$). The safety gate falsely refuses $\TerraCatAFalseRefusals$/100 Category~A queries (and $\TerraCatBFalseRefusals$/100 on Category~B; see Limitation~\ref{item:falserefusal}).
    \item \textbf{Adversarial Deception Robustness (Category D):} On deceptive prompts embedding authentic titles with fabricated premises (e.g., maritime quotas in \textit{Smith v. Allwright}), TERRA achieves $\mathbf{\TerraCatDRejection}$ explicit safety refusals alongside $\TerraCatDRefutationRate$ verified substantive refutations ($\TerraCatDRefutations$), delivering $\mathbf{\TerraCatDSafeHandlingRate}$ total safe handling ($\TerraCatDSafeHandling$; Table~\ref{tab:refusal_breakdown}, Panel~B). In paired binary contingency testing across the 75 matched adversarial queries, TERRA demonstrates $100\%$ monotonic containment over all baseline-safe completions ($c=0$ across all four comparative baselines: every query safely resolved by any baseline is also safely resolved by TERRA). On total safe handling, TERRA handled more traps safely than each baseline, with every discordant pair favouring TERRA: GraphRAG-Adapted ($\GraphRAGCatDSafeHandling \to \TerraCatDSafeHandling$, $b=6, c=0$, exact two-tailed McNemar $p = \TerraVsGraphRAGCatDSafeMcNemarP$), Flat Vector RAG ($\FlatRAGCatDSafeHandling$, $b=5, c=0$, $p = \TerraVsFlatRAGCatDSafeMcNemarP$), HippoRAG-Adapted ($\HippoRAGCatDSafeHandling$, $b=3, c=0$, $p = \TerraVsHippoRAGCatDSafeMcNemarP$), and Direct LLM ($\DirectLLMCatDSafeHandling$, $b=2, c=0$, $p = \TerraVsDirectLLMCatDSafeMcNemarP$). These eight Category~D McNemar comparisons (four safe-handling, four explicit-refusal) are secondary tests outside the pre-specified $M=20$ family and are Holm-adjusted as one family; after adjustment no safe-handling difference is significant (GraphRAG-Adapted: adjusted $p = \TerraVsGraphRAGCatDSafeMcNemarHolmP$; Flat Vector RAG: $\TerraVsFlatRAGCatDSafeMcNemarHolmP$). On explicit refusal, TERRA refused more traps than Direct LLM ($\TerraCatDRejection$ vs.\ $\DirectLLMCatDRejection$, $b=48, c=0$, $p = \TerraVsDirectLLMCatDRefusalMcNemarP$; Holm-adjusted $p = \TerraVsDirectLLMCatDRefusalMcNemarHolmP$) and had the highest audited refusal rate of the evaluated systems (vs.\ Flat RAG $\FlatRAGCatDRejectionRate$, $b=21, c=11$, $p = \TerraVsFlatRAGCatDRefusalMcNemarP$; HippoRAG $\HippoRAGCatDRejectionRate$, $b=22, c=8$, $p = \TerraVsHippoRAGCatDRefusalMcNemarP$; GraphRAG $\GraphRAGCatDRejectionRate$, $b=26, c=10$, $p = \TerraVsGraphRAGCatDRefusalMcNemarP$); after Holm adjustment none of these three retrieval-baseline differences is significant (adjusted $p = \TerraVsFlatRAGCatDRefusalMcNemarHolmP$, $\TerraVsHippoRAGCatDRefusalMcNemarHolmP$, $\TerraVsGraphRAGCatDRefusalMcNemarHolmP$). Under calibrated dual-budget traversal and fast-path grounding, fatal constraint violations were observed on only $\TerraCatDFatalErrors$, with $\TerraCatDHedges$ uncommitted hedges. The only two remaining fatal errors (queries D033 and D034) are reciprocal procedural misattributions between closely coupled apportionment landmarks (\textit{Baker v. Carr} [1962] and \textit{Reynolds v. Sims} [1964]), where the generator transposed the threshold justiciability holding with the substantive apportionment standard. Category~D labels come from a calibrated semantic verifier served by the same model family as the generator (Limitation~\ref{item:samefamily}).
    \item \textbf{Out-of-Context Safety (Category C):} On queries demonstrably outside indexed U.S. Supreme Court case law (spanning natural sciences, foreign law, municipal ordinances, contemporary software, and corporate accounting), TERRA achieves a $\mathbf{\TerraCatCRejectionRate}$ safety refusal rate ($\mathbf{\TerraCatCRejection}$), tying GraphRAG ($\mathbf{\GraphRAGCatCRejection}$) and outperforming Flat RAG ($\FlatRAGCatCRejection$) and Direct LLM ($\DirectLLMCatCRejection$).
\end{itemize}

\subsection{Standalone Traffic Cop Classifier Evaluation and Defense-in-Depth}
\label{subsec:router-standalone-eval}

While the component-level ablation in Table~\ref{tab:ablation} establishes the system-level operational necessity of query routing (Ablation~3 in the pilot: forcing all queries into multi-hop GraphRAG traversal causes a $\AblationRouterLatencySlowdownFactor\times$ latency penalty alongside $\AblationThreeCatAFalseRefusalRate$ false refusals on factual lookups), evaluating the router purely as a component in an end-to-end pipeline leaves its intrinsic classification quality unquantified. Furthermore, Second Review Issue~\#6 noted that query intent classification, domain boundary pre-filtering, and routing error mechanisms required formal, standalone empirical characterization.

To resolve this question, we evaluate the Traffic Cop Router as an independent, standalone binary intent classifier across all $N=\RouterTotalQueries$ canonical benchmark queries. All evaluations were executed deterministically at temperature $T=0.0$ using Google Gemma 3 12B Instruction-tuned (\texttt{\RouterModel}) served via the Databricks Model Serving API under strict schema validation (\texttt{RoutingDecision}), with certified zero fallback contamination ($0/350$). The complete specification, prompt template, and Pydantic schema are published in \texttt{docs/ROUTER\_SPECIFICATION.md}.

\input{tables/table_router_eval}

\subsubsection{Standalone Intent Classification Performance}
As reported in Table~\ref{tab:router_eval} (Panel~A), the Traffic Cop Router achieves high standalone classification fidelity:
\begin{itemize}
    \item \textbf{High Overall Accuracy and Balanced $F_1$:} Across the full 350-query suite, the router achieves an overall accuracy of $\mathbf{\RouterAccuracy}$ ($\RouterCorrectCount$) and a macro-averaged $F_1$ score of $\mathbf{\RouterMacroFOne}$ (macro precision $\RouterMacroPrecision$, macro recall $\RouterMacroRecall$). On the factual class (Gold EASY, Category~A), the classifier achieves $\RouterEasyFOne$ $F_1$ with perfect recall ($\RouterEasyRecall$, $100/100$). On complex inquiries (Gold HARD, Categories~B, C, and D), it achieves $\RouterHardFOne$ $F_1$ with $\RouterHardPrecision$ precision ($240/240$). All \RouterErrorCount\ misroutes are Category~D queries sent to the fast path, which in the canonical run still recorded retrieval and NLI grading on every query; all ten were safely handled (8 flagged refusals, for which the grader did not pass the context, and 2 unflagged textual refusals, D068 and D070, whose context the grader passed; none affirmed a fatal error).
    \item \textbf{Asymmetric Error Profile Favoring Safety:} The directional error profile reflects deliberate architectural alignment. The False-HARD rate on factual queries is certified at $\mathbf{\RouterFalseHardRate}$ ($0/100$), ensuring that single-point factual lookups are never needlessly penalized by multi-hop graph expansion or context flooding. Across all complex queries, the False-EASY rate is $\mathbf{\RouterFalseEasyRate}$ ($10/250$), with zero false-EASY routes on Category~B multi-hop evolutionary chains ($0/100$) and Category~C out-of-context unanswerable probes ($0/75$).
    \item \textbf{Standalone Inference Latency Profile:} In isolated execution, the Gemma 3 12B router operates with a median wall-clock latency of $\mathbf{\RouterLatencyMedianMs\text{ ms}}$ and a mean of $\RouterLatencyMeanMs\text{ ms}$ (p90 of $\RouterLatencyPNinetyMs\text{ ms}$, spanning from a minimum of $\RouterLatencyMinMs\text{ ms}$ to an un-cached cold start maximum of $\RouterLatencyMaxMs\text{ ms}$). This lightweight footprint enables responsive front-line complexity triage before invoking heavier 70B parameter models.
\end{itemize}

\subsubsection{Mechanistic Analysis of Category D Misroutes and Defense-in-Depth}
As disaggregated in Table~\ref{tab:router_eval} (Panel~B), all 10 queries classified as EASY by the router belong exclusively to Category~D (Adversarial Deception Traps; $\RouterCatDFalseEasyRate$, $\RouterCatDFalseEasyCount/75$). A forensic examination of these 10 queries (\texttt{D021}, \texttt{D039}, \texttt{D060}, \texttt{D062}, \texttt{D068}, \texttt{D069}, \texttt{D070}, \texttt{D072}, \texttt{D073}, \texttt{D074}) reveals a uniform syntactic structure: each query is formulated as a single-sentence constitutional inquiry referencing a single landmark authority (e.g., query \texttt{D021}: \textit{``Did the Supreme Court in Mapp v. Ohio establish that police custodial interrogations must be preceded by four-point constitutional rights warnings?''}). Because these queries lack explicit multi-hop syntactic connectors (such as comparative doctrine tracing across decades), the SLM router classifies them based on surface-level syntactic simplicity as single-case lookups.

Crucially, routing these adversarial queries to the fast path does \textit{not} induce generative hallucination or credulous acceptance of false premises. In the canonical run the fast route (the Grounded Fast Path, Issue~\#12) recorded retrieval and NLI grading on all 110 queries it handled, so its generations were conditioned on retrieved context rather than ungrounded parametric completion. Forensic auditing of the downstream completions under independent dual judging records a $\mathbf{\RouterCatDSafeHandlingRate}$ safe-handling rate ($10/10$):
\begin{itemize}
    \item \textbf{Safe Epistemic Refusals ($8/10$, $80.0\%$):} On 8 queries (\texttt{D021}, \texttt{D039}, \texttt{D060}, \texttt{D062}, \texttt{D069}, \texttt{D072}, \texttt{D073}, \texttt{D074}), the NLI grader did not pass the retrieved context and the system refused deterministically, without generating an answer.
    \item \textbf{Unflagged Textual Refusals ($2/10$, $20.0\%$):} On 2 queries (\texttt{D068} regarding the Lottery Case and \texttt{D070} regarding the oleomargarine tax), the model generated unflagged textual refusals asserting absence or insufficiency of retrieved context without triggering the refusal flag (the grader had passed their retrieved context), scoring 1.000 by dual judges.
    \item \textbf{Zero Fatal Violations ($0/10$, $0.0\%$):} Across all 10 misrouted adversarial cases, zero fatal constraint violations were observed.
\end{itemize}
These ten cases are consistent with downstream checks limiting the cost of a misroute: even when an adversarial input slips past the front-line small language model router via deceptive single-sentence syntax, the downstream grader and generator handled all ten safely in the canonical run. With only ten queries and judge-based scoring, this is a small-sample observation, not evidence that downstream mechanisms reliably prevent hallucination.

\subsection{Stage Latency Breakdown and Fast-Path Routing Efficiency}
\label{subsec:latency-distribution}

Because mean latency metrics in multi-hop systems are heavily skewed by extreme outliers, Table~\ref{tab:latency_percentiles} reports the complete latency distribution across pipeline stages and execution paths.

\begin{table}[htbp]
\caption{Pipeline Execution Latency Percentiles (Milliseconds). Panel~A: Gross wall-clock time (end-to-end including any server-imposed retry sleep). Panel~B: Net time (retry sleep subtracted). TERRA fast route ($N=110$): 100 Category~A queries plus 10 Category~D queries routed EASY by the traffic-cop; in the canonical run the fast route recorded router inference, retrieval, NLI grading and generation (retrieval and grading times are non-zero on all 110 queries). TERRA hard route ($N=240$): 100 Category~B + 75 Category~C + 65 Category~D queries; hard path executes router + ChromaDB retrieval + BFS traversal + NLI grading + generator. Baselines ran sequentially (5--6 Sept.); TERRA ran with four concurrent workers (14 Sept.), so comparisons are indicative only. Forced-hard-path latency (Ablation~3) is in Table~\ref{tab:ablation} (median $\AblationThreeLatencyMedianSeconds$~s).}
\label{tab:latency_percentiles}
\centering
\resizebox{0.98\columnwidth}{!}{%
\begin{tabular}{lrrrrr}
\toprule
\textbf{Execution Path} & \textbf{p25} & \textbf{Median} & \textbf{p75} & \textbf{p95} & \textbf{Mean $\pm$ SD} \\
\midrule
\multicolumn{6}{l}{\textit{Panel A: Gross latency (ms)}} \\
\midrule
Direct LLM (No RAG) & $\DirectLLMLatencyPTwentyFiveMs$ & $\DirectLLMLatencyMedianMs$ & $\DirectLLMLatencyPSeventyFiveMs$ & $\DirectLLMLatencyPNinetyFiveMs$ & $\DirectLLMLatencyMeanAndStdMs$ \\
Flat Vector RAG & $\FlatRAGLatencyPTwentyFiveMs$ & $\FlatRAGLatencyMedianMs$ & $\FlatRAGLatencyPSeventyFiveMs$ & $\FlatRAGLatencyPNinetyFiveMs$ & $\FlatRAGLatencyMeanAndStdMs$ \\
GraphRAG-Adapted & $\GraphRAGLatencyPTwentyFiveMs$ & $\GraphRAGLatencyMedianMs$ & $\GraphRAGLatencyPSeventyFiveMs$ & $\GraphRAGLatencyPNinetyFiveMs$ & $\GraphRAGLatencyMeanAndStdMs$ \\
HippoRAG-Adapted & $\HippoRAGLatencyPTwentyFiveMs$ & $\mathbf{\HippoRAGLatencyMedianMs}$ & $\HippoRAGLatencyPSeventyFiveMs$ & $\mathbf{\HippoRAGLatencyPNinetyFiveMs}$ & $\HippoRAGLatencyMeanAndStdMs$ \\
\textbf{TERRA Overall} & $\TerraLatencyPTwentyFiveMs$ & $\TerraLatencyMedianMs$ & $\TerraLatencyPSeventyFiveMs$ & $\TerraLatencyPNinetyFiveMs$ & $\TerraLatencyMeanAndStdMs$ \\
\midrule
TERRA Fast Route ($N=110$; Cat A + 10$\times$Cat D) & $\TerraFastPathLatencyPTwentyFiveMs$ & $\TerraFastPathLatencyMedianMs$ & $\TerraFastPathLatencyPSeventyFiveMs$ & $\TerraFastPathLatencyPNinetyFiveMs$ & $\TerraFastPathLatencyMeanAndStdMs$ \\
TERRA Hard Route ($N=240$; Cat B + C + 65$\times$Cat D) & $\TerraHardPathLatencyPTwentyFiveMs$ & $\TerraHardPathLatencyMedianMs$ & $\TerraHardPathLatencyPSeventyFiveMs$ & $\TerraHardPathLatencyPNinetyFiveMs$ & $\TerraHardPathLatencyMeanAndStdMs$ \\
\midrule
\multicolumn{6}{l}{\textit{Panel B: Net latency, retry sleep subtracted (ms; baselines unaffected)}} \\
\midrule
\textbf{TERRA Overall (net)} & --- & $\TerraNetLatencyMedianMs$ & --- & $\TerraNetLatencyPNinetyFiveMs$ & --- \\
TERRA Fast Route (net, $N=110$) & $\TerraFastPathNetLatencyPTwentyFiveMs$ & $\TerraFastPathNetLatencyMedianMs$ & $\TerraFastPathNetLatencyPSeventyFiveMs$ & $\TerraFastPathNetLatencyPNinetyFiveMs$ & $\TerraFastPathNetLatencyMeanAndStdMs$ \\
TERRA Hard Route (net, $N=240$) & $\TerraHardPathNetLatencyPTwentyFiveMs$ & $\TerraHardPathNetLatencyMedianMs$ & $\TerraHardPathNetLatencyPSeventyFiveMs$ & $\TerraHardPathNetLatencyPNinetyFiveMs$ & $\TerraHardPathNetLatencyMeanAndStdMs$ \\
\midrule
\multicolumn{6}{l}{\textit{Secondary block: Category-level TERRA gross latency (ms)}} \\
\midrule
TERRA Category A & $\TerraCatALatencyPTwentyFiveMs$ & $\TerraCatALatencyMedianMs$ & --- & $\TerraCatALatencyPNinetyFiveMs$ & $\TerraCatALatencyMeanAndStdMs$\textsuperscript{a} \\
TERRA Category B & $\TerraCatBLatencyPTwentyFiveMs$ & $\TerraCatBLatencyMedianMs$ & --- & $\TerraCatBLatencyPNinetyFiveMs$ & $\TerraCatBLatencyMeanAndStdMs$\textsuperscript{a} \\
\bottomrule
\end{tabular}%
}
\vspace{1pt}
{\footnotesize \textsuperscript{a}~Category-level gross latency across all Category~A ($N=100$) and Category~B ($N=100$) queries; p75 omitted to limit column width.}
\end{table}

Route-stratified, TERRA's fast route ($N=\RouterFastPathCount$) has a median of \TerraFastRouteMedianS\ and its hard route ($N=\RouterHardPathCount$) \TerraHardRouteMedianS; overall, the median is \TerraLatencyMedianSeconds~s gross (\TerraNetLatencyMedianSeconds~s net of retry sleeps) and the 95th percentile \TerraLatencyPNinetyFiveSeconds~s (\TerraNetLatencyPNinetyFiveSeconds~s net). Baseline medians are \DirectLLMLatencyMedianSeconds~s (Direct), \FlatRAGLatencyMedianSeconds~s (Flat), \GraphRAGLatencyMedianSeconds~s (GraphRAG) and \HippoRAGLatencyMedianSeconds~s (HippoRAG); TERRA is therefore slower than every baseline. TERRA ran with four concurrent workers on a shared serving endpoint and \TerraThrottledRows/350 of its queries hit HTTP~429 back-off (\TerraThrottleSleepSeconds~s in total); the baselines ran sequentially with no throttling, so latencies are not from a like-for-like harness.

\subsection{Statistical Significance and Multiple Testing Control}
\label{subsec:statistical-significance}

To evaluate statistical significance, we execute paired Wilcoxon signed-rank tests ($\alpha=0.05$) and report paired Cohen's $d_z$ effect sizes, bootstrap 95\% confidence intervals, raw $p$-values, FDR $q$-values, and exact paired outcome counts ($n^+ / n^0 / n^-$). Following \cite{hoenig2001abuse}, retrospective post-hoc power is omitted as an uninformative monotonic transformation of $p$-values; statistical power is characterized prospectively for future study design via paired-mean MDES approximations ($d_z = 0.283$ at $N=100$, $d_z = 0.200$ at $N=200$, $1-\beta=0.80$), while discrete refusal is evaluated via exact Fisher and McNemar tests. Table~\ref{tab:power_analysis} reports paired evaluations across Category~B and combined in-domain queries.

\begin{table}[htbp]
\caption{Statistical Significance, Benjamini-Hochberg FDR Control ($M=20$ Family), and Exact Paired Outcomes Across Paired Evaluations.}
\label{tab:power_analysis}
\centering
\resizebox{0.98\columnwidth}{!}{%
\begin{tabular}{lccccc}
\toprule
\textbf{Comparison} & \textbf{Mean Delta [95\% CI]} & \textbf{Cohen's $d_z$} & \textbf{Wilcoxon $p$} & \textbf{FDR $q$} & \textbf{Paired ($n^+/n^0/n^-$)} \\
\midrule
TERRA vs. Direct LLM (Cat B) & $\CatBTerraVsDirectLLMDelta$ $\CatBTerraVsDirectLLMCI$ & $\CatBTerraVsDirectLLMDz$ & $\CatBTerraVsDirectLLMWilcoxonP$ & $\CatBTerraVsDirectLLMAdjQ$ & $\TerraVsDirectLLMCatBWins/\TerraVsDirectLLMCatBTies/\TerraVsDirectLLMCatBLosses$ \\
TERRA vs. Direct LLM (Combined A+B) & $\CombinedTerraVsDirectLLMDelta$ $\CombinedTerraVsDirectLLMCI$ & $\CombinedTerraVsDirectLLMDz$ & $\CombinedTerraVsDirectLLMWilcoxonP$ & $\CombinedTerraVsDirectLLMAdjQ$ & $\TerraVsDirectLLMCombinedWins/\TerraVsDirectLLMCombinedTies/\TerraVsDirectLLMCombinedLosses$ \\
TERRA vs. Direct LLM (Cat C Refusal)$^\dagger$ & $+100.0\%$ Diff. & --- & $\mathbf{\CatCTerraVsDirectLLMFisherP}$ & $\mathbf{\CatCTerraVsDirectLLMAdjQ}$ & $75/75 \text{ vs. } 0/75$ \\
\midrule
TERRA vs. Flat RAG (Cat B) & $\mathbf{\CatBTerraVsFlatRAGDelta}$ $\CatBTerraVsFlatRAGCI$ & $\mathbf{\CatBTerraVsFlatRAGDz}$ & $\mathbf{\CatBTerraVsFlatRAGWilcoxonP}$ & $\mathbf{\CatBTerraVsFlatRAGAdjQ}$ & $\TerraVsFlatRAGCatBWins/\TerraVsFlatRAGCatBTies/\TerraVsFlatRAGCatBLosses$ \\
TERRA vs. Flat RAG (Combined A+B) & $\mathbf{\CombinedTerraVsFlatRAGDelta}$ $\CombinedTerraVsFlatRAGCI$ & $\CombinedTerraVsFlatRAGDz$ & $\mathbf{\CombinedTerraVsFlatRAGWilcoxonP}$ & $\mathbf{\CombinedTerraVsFlatRAGAdjQ}$ & $\TerraVsFlatRAGCombinedWins/\TerraVsFlatRAGCombinedTies/\TerraVsFlatRAGCombinedLosses$ \\
\midrule
TERRA vs. GraphRAG (Cat B) & $\mathbf{\CatBTerraVsGraphRAGDelta}$ $\CatBTerraVsGraphRAGCI$ & $\mathbf{\CatBTerraVsGraphRAGDz}$ & $\mathbf{\CatBTerraVsGraphRAGWilcoxonP}$ & $\mathbf{\CatBTerraVsGraphRAGAdjQ}$ & $\TerraVsGraphRAGCatBWins/\TerraVsGraphRAGCatBTies/\TerraVsGraphRAGCatBLosses$ \\
TERRA vs. GraphRAG (Combined A+B) & $\mathbf{\CombinedTerraVsGraphRAGDelta}$ $\CombinedTerraVsGraphRAGCI$ & $\CombinedTerraVsGraphRAGDz$ & $\mathbf{\CombinedTerraVsGraphRAGWilcoxonP}$ & $\mathbf{\CombinedTerraVsGraphRAGAdjQ}$ & $\TerraVsGraphRAGCombinedWins/\TerraVsGraphRAGCombinedTies/\TerraVsGraphRAGCombinedLosses$ \\
\midrule
TERRA vs. HippoRAG (Cat B) & $\mathbf{\CatBTerraVsHippoRAGDelta}$ $\CatBTerraVsHippoRAGCI$ & $\mathbf{\CatBTerraVsHippoRAGDz}$ & $\mathbf{\CatBTerraVsHippoRAGWilcoxonP}$ & $\mathbf{\CatBTerraVsHippoRAGAdjQ}$ & $\TerraVsHippoRAGCatBWins/\TerraVsHippoRAGCatBTies/\TerraVsHippoRAGCatBLosses$ \\
TERRA vs. HippoRAG (Combined A+B) & $\CombinedTerraVsHippoRAGDelta$ $\CombinedTerraVsHippoRAGCI$ & $\CombinedTerraVsHippoRAGDz$ & $\CombinedTerraVsHippoRAGWilcoxonP$ & $\CombinedTerraVsHippoRAGAdjQ$ & $\TerraVsHippoRAGCombinedWins/\TerraVsHippoRAGCombinedTies/\TerraVsHippoRAGCombinedLosses$ \\
\bottomrule
\end{tabular}%
}
\vspace{1mm}
\raggedright \footnotesize $^\dagger$Evaluated via Fisher's Exact Test on Out-of-Context safety refusal ($75/75$ vs. $0/75$, complete empirical separation). Paired outcomes report TERRA wins ($n^+$), ties ($n^0$), and losses ($n^-$); Category C reports contingency refusal counts ($75/75\text{ vs. }0/75$). FDR correction was applied jointly across the complete pre-specified family of $M=20$ hypothesis tests (4 comparative baselines $\times$ 5 evaluation slices: Categories A, B, C, D, and Combined A+B); primary multi-hop evolutionary and safety comparisons are reported above. The complete, unredacted 20-test family with sorted $p$-values and exact paired outcomes is reported in Appendix~\ref{sec:appendix-fdr} (Table~\ref{tab:full_fdr_family}) for independent verification.
\end{table}

\enlargethispage{2\baselineskip}
Statistical testing confirms that TERRA's Category B reasoning gains remain statistically significant across all retrieval baselines under Benjamini-Hochberg FDR correction ($M=20$ family, Table~\ref{tab:full_fdr_family}):
\begin{itemize}
    \item \textbf{Against Flat Vector RAG:} TERRA achieves a $\mathbf{\CatBTerraVsFlatRAGDelta}$ gain on Category B ($p = \CatBTerraVsFlatRAGWilcoxonP$, FDR $q = \CatBTerraVsFlatRAGAdjQ$, $d_z = \CatBTerraVsFlatRAGDz$) and $\mathbf{\CombinedTerraVsFlatRAGDelta}$ on combined in-domain queries ($p = \CombinedTerraVsFlatRAGWilcoxonP$, FDR $q = \CombinedTerraVsFlatRAGAdjQ$).
    \item \textbf{Against GraphRAG-Adapted:} TERRA achieves a $\mathbf{\CatBTerraVsGraphRAGDelta}$ advantage on Category B ($p = \CatBTerraVsGraphRAGWilcoxonP$, FDR $q = \CatBTerraVsGraphRAGAdjQ$, $d_z = \CatBTerraVsGraphRAGDz$) and $\mathbf{\CombinedTerraVsGraphRAGDelta}$ on combined queries ($p = \CombinedTerraVsGraphRAGWilcoxonP$, FDR $q = \CombinedTerraVsGraphRAGAdjQ$).
    \item \textbf{Against HippoRAG-Adapted:} TERRA achieves a $\mathbf{\CatBTerraVsHippoRAGDelta}$ advantage on Category B ($p = \CatBTerraVsHippoRAGWilcoxonP$, FDR $q = \CatBTerraVsHippoRAGAdjQ$, $d_z = \CatBTerraVsHippoRAGDz$).
    \item \textbf{Safety Refusal Robustness:} On Category C, TERRA's $100\%$ refusal ($75/75$) demonstrates decisive refusal advantage over Direct LLM ($0/75$, Fisher's exact $p = \CatCTerraVsDirectLLMFisherP$, FDR $q = \CatCTerraVsDirectLLMAdjQ$).
    \item \textbf{Against Direct LLM (Non-Retrieval Reference):} On Category~B, Direct LLM achieves higher faithfulness than TERRA ($\Delta = \CatBTerraVsDirectLLMDelta$, $p = \CatBTerraVsDirectLLMWilcoxonP$, FDR $q = \CatBTerraVsDirectLLMAdjQ$, Table~\ref{tab:full_fdr_family}, Rank~8), whereas on combined in-domain queries (Categories~A+B), TERRA leads by $\mathbf{\CombinedTerraVsDirectLLMDelta}$ ($p = \CombinedTerraVsDirectLLMWilcoxonP$, FDR $q = \CombinedTerraVsDirectLLMAdjQ$), driven by Category~A ($\CatATerraVsDirectLLMDelta$, $q = \CatATerraVsDirectLLMAdjQ$).
\end{itemize}

\subsection{Temporal Parametric Leakage Stress Test}
\label{subsec:temporal-leakage}

While retrospective legal retrieval requires evaluating questions strictly under historical jurisprudence as of temporal cutoff $T$, modern generative language models possess extensive pretraining corpora containing post-$T$ legal history. Consequently, models risk emitting \textit{temporal parametric leakage}---anachronistically importing future doctrinal shifts, statutory overhauls, or later overrulings into historical analyses. To rigorously quantify and stress-test this vulnerability, we preregistered and executed a dedicated 50-query temporal leakage probe across five historical Supreme Court horizons ($T \in \{1890, 1930, 1950, 1970, 1980\}$), spanning seven major doctrinal domains and balanced across case prominence and multi-hop reasoning depths. Six probes (12\%) targeting historical cases decided between 1857 and 1905 (evaluated at cutoffs $T=1890$ and $T=1930$) have target nodes with zero out-edges in the EEG due to the corpus citation-parsing boundary. Arc evolution cases for these entries are historically verified against primary legal sources but are not EEG graph edges. All leakage metrics are computed identically across all 50 probes regardless of arc status.

\input{tables/table_temporal_leakage_probe}

\subsubsection{Experimental Design and Metric Decoupling}
The probe evaluates \TemporalTotalEvaluations\ responses generated across four comparative architectures: Direct LLM, Flat Vector RAG, $\text{TERRA}_{\text{ungated}}$ ($L_{\text{gen}}$), and $\text{TERRA}_{\text{gated}}$ ($L_{\text{final}}$). To separate baseline capabilities from prompt compliance, each architecture is tested under two distinct regimes: \textbf{Condition A} (standard task prompt without explicit temporal constraint, measuring default end-to-end containment) and \textbf{Condition B} (explicit instruction to strictly adhere to year $T$ and ignore future precedents, measuring prompt-assisted containment). All candidate generations were evaluated under a preregistered 3-tier severity rubric: Level~0 (clean adherence to pre-$T$ doctrine), Level~1 (non-prejudicial historical mention or temporal disclaimer), and Level~2 (fatal substantive reliance on post-$T$ law). The primary metric, \textbf{Substantive Post-$T$ Leakage Rate (SPTLR)}, captures the proportion of Level~2 fatal failures.

Critically, evaluating temporal containment solely via SPTLR creates a severe vulnerability to metric distortion: systems that simply refuse to answer generate zero post-$T$ text, appearing perfectly ``leak-free'' despite failing the underlying legal task. To eliminate this confound, we decouple containment from legal problem-solving by measuring: (1)~\textbf{Refusal Rate}, isolating epistemic or safety rejections; and (2)~\textbf{Historical Legal Correctness Rate (HLCR)}, assessing whether the system correctly articulates the operative contemporaneous rule as of year $T$. Finally, we measure the \textbf{NLI Gate Leakage Interception Rate (NGLIR)}, defined as the proportion of ungated candidate leaks intercepted by the entailment gate on the conditioning denominator of leaking candidates ($D$).

\subsubsection{Evaluator Reliability, Concordance, and Consensus Protocol}
Responses were independently evaluated by two distinct model-based judges: \texttt{qwen3-next-80b} and \texttt{gpt-oss-120b}. Across $N=\TemporalValidPairs$ complete pairs, the evaluators achieved $\TemporalRawConcordance$ raw concordance and $100\%$ monotonic inclusion (every leak identified by GPT-OSS was a strict subset of Qwen's flags). However, the quadratic weighted Cohen's $\kappa$ was $\TemporalQwk$, appearing considerably lower than the benchmark-wide agreement ($\kappa \approx 0.70$). Forensic diagnosis confirms that this divergence is an artifact of the \textit{Kappa Paradox}~\cite{feinstein1990high}: GPT-OSS exhibited a severe $91.7\%$ base-rate skew toward Level~0 ($289/315$ clean), which mechanically inflated expected chance agreement $P_e$ to $\approx 60\%$, deflating $\kappa$ despite high raw agreement. Furthermore, qualitative analysis revealed an evaluator threshold asymmetry: Qwen operated as a hyper-vigilant lexical scanner (flagging any passing historical reference as Level~1 or Level~2), whereas GPT-OSS evaluated holistic legal operability (classifying responses as clean if post-$T$ references did not alter the substantive legal outcome).

To guarantee objective reproducibility, all evaluations were reconciled via a pre-specified, deterministic 5-tier consensus protocol: (1)~safe refusals automatically resolve to Level~0; (2)~unanimous agreement is preserved; (3)~lexical breach of forbidden historical phrases assigns Level~2; (4)~substantive quotation of post-$T$ law assigns Level~2; and (5)~subtle mentions resolve conservatively to $\max(s_1, s_2) = \text{Level~1}$. Crucially, tiers 3--5 bias strictly \textit{toward} flagging leakage, ensuring that consensus adjudication is conservative against TERRA.

\subsubsection{Empirical Missingness Mechanism and Sensitivity Bounds}
Among the $400$ evaluations, $85$ GPT-OSS calls suffered provider-side safety filter terminations. To determine whether this missingness introduced bias, we audited failure distribution across queries and domains. Missingness was strongly non-random: failures clustered heavily in Intellectual Property/Patents ($32.1\%$, $18/56$), Antitrust ($23.2\%$, $13/56$), and Commerce ($23.2\%$, $13/56$). Furthermore, Direct LLM accounted for $30$ of the $85$ missing evaluations, concentrated on queries with high potential for blatant modern citation. Because safety filters disproportionately suppressed Direct LLM's most egregious leakage instances, complete-case Direct LLM leakage ($28.0\%$ in Condition A) is an underestimate. To maintain absolute statistical honesty, Table~\ref{tab:temporal_leakage_probe} reports full sensitivity bounds $[SPTLR_{\text{lower}}, SPTLR_{\text{upper}}]$ encompassing all $2^{85}$ missing-data assignments, alongside complete $(b, c)$ discordant-pair reporting for all hypotheses.

\subsubsection{Primary Hypothesis $H_1$: Metric Decoupling and Refusal Masking}
Discordant pairs are defined in standard McNemar form: $b$ denotes query pairs where $\text{TERRA}_{\text{gated}}$ is clean (Level~0 or Level~1) while the baseline exhibits Level~2 substantive leakage; $c$ denotes query pairs where $\text{TERRA}_{\text{gated}}$ exhibits Level~2 leakage while the baseline is clean. For each comparison, Table~\ref{tab:temporal_leakage_probe} reports both the un-imputed complete-case pairs ($N_{\text{cc}}$, where both judges completed evaluation under the 5-tier consensus protocol) and the full $N=50$ sensitivity bounds ($[b_{\min}, b_{\max}], [c_{\min}, c_{\max}]$ adding bounded missing-pair additions across all $2^{85}$ feasible permutations).

Preregistered Primary Hypothesis $H_1$ predicted that under default prompting (Condition A), $\text{TERRA}_{\text{gated}}$ would achieve lower raw SPTLR than Flat Vector RAG. As reported in Table~\ref{tab:temporal_leakage_probe}, this hypothesis is \textbf{not supported}. Under complete cases ($N_{\text{cc}}=36$), consensus yields $b=0, c=3$ (one-sided exact McNemar $p = \HoneSptlrPval$; two-sided $p=0.2500$, or $p=0.0625$ if missing records are zero-imputed at $N=50$). Bounding over all missingness permutations yields $b \in [0, 5], c \in [7, 8]$. Flat Vector RAG exhibits a lower raw leakage interval ($\FlatRagSptlrA$) than $\text{TERRA}_{\text{gated}}$ ($\TerraGatedSptlrA$). 

However, forensic examination of discordant pairs reveals the exact mechanistic origin of this result: Flat RAG produced zero leaks on queries where TERRA was clean ($b=0$) because Flat RAG suffered a massive $\FlatRagRefusalA$ refusal collapse ($23/50$ queries), completely abdicating on nearly half the dataset. When a retrieval system refuses to answer, it mechanically cannot leak post-$T$ doctrine. Conversely, TERRA attempted and successfully engaged with these difficult historical inquiries (refusal rate $\TerraGatedRefusalA$), incurring $c=3$ complete-case Level-2 classifications on complex multi-jurisdictional tests where Flat RAG simply refused.

When evaluating substantive legal utility via Historical Legal Correctness Rate (HLCR), TERRA significantly outperforms Flat Vector RAG: $\text{TERRA}_{\text{gated}}$ achieves $\mathbf{\TerraGatedHlcrA}$ consensus correctness compared to only $\FlatRagHlcrA$ for Flat RAG ($b=20, c=1$; exact McNemar $p = \HoneHlcrPvalFormatted$, $p < 0.0001$). Thus, TERRA effectively doubles reliable historical legal competence, and its apparent raw SPTLR parity in Condition A is solely an artifact of Flat RAG's refusal volume masking.

\subsubsection{Confirmatory Hypotheses $H_2$--$H_4$ and Directional Stability}
On baseline comparisons where models actively generate answers, TERRA demonstrates robust containment advantage:
\begin{itemize}
    \item \textbf{$H_2$ (TERRA vs. Direct LLM in Condition A):} Complete cases ($N_{\text{cc}}=32$) yield $b=2, c=4$. Bounding over the 18 missing Direct LLM evaluations yields $b \in [11, 16]$ and $c \in [5, 6]$ ($p \in [\HtwoPvalMin, \HtwoPvalMax]$). Because $c_{\max} = 6 < b_{\min} = 11$, \textbf{Directional Stability} is formally certified: across all $2^{85}$ missing-data permutations, TERRA unconditionally generates more clean historical answers than unaugmented parametric generation, reaching nominal significance ($p = \HtwoPvalMin$).
    \item \textbf{$H_3$ (TERRA vs. Flat Vector RAG in Condition B):} Complete cases ($N_{\text{cc}}=29$) yield $b=5, c=1$. Under explicit temporal prompting, Flat RAG's refusal rate drops to $30.0\%$, exposing its susceptibility to retrieval-induced leakage ($\FlatRagSptlrB$). Full bounds span $b \in [10, 13], c \in [4, 8]$ ($p \in [\HthreePvalMin, \HthreePvalMax]$), achieving certified Directional Stability ($c_{\max} = 8 < b_{\min} = 10$).
    \item \textbf{$H_4$ (TERRA vs. Direct LLM in Condition B):} Complete cases ($N_{\text{cc}}=31$) yield $b=1, c=1$. Even when Direct LLM is explicitly instructed to adhere to year $T$, parametric weights cause frequent anachronistic leakage ($\DirectLlmSptlrB$). Full bounds span $b \in [8, 11], c \in [2, 6]$ ($p \in [\HfourPvalMin, \HfourPvalMax]$), maintaining Directional Stability ($c_{\max} = 6 < b_{\min} = 8$).
\end{itemize}

\subsubsection{Gate Ablation ($H_5, H_6$) and Architectural Attribution (NGLIR)}
Evaluating the NLI gate ablation ($H_5, H_6$) indicates that marginal bulk leakage reduction over raw EEG retrieval is statistically underpowered at $N=50$ ($b \in [0, 2], c \in [1, 1]$ for Condition A; $b \in [2, 5], c \in [0, 4]$ for Condition B; complete cases $b=0, c=1$ and $b=2, c=0$). Rather than acting as a high-volume statistical filter, forensic auditing confirms that the NLI gate functions as a \textbf{targeted qualitative safety firewall}. Across the probe, the gate executed six critical interventions: intercepting one substantive Level-2 leak in Condition B (\textit{United States v. Jackalow}, converting an anachronistic maritime jurisdictional test into a safe refusal) and five ungrounded hallucinations where the generator attempted to fabricate non-existent pre-$T$ precedents.

\subsection{Adversarial Generalization Under Weak Parametric Memory: Obscure Overruled Precedents}
\label{subsec:obscure-generalization}
Table~\ref{tab:obscure_overruled_eval} reports the obscure suite (Section~\ref{subsec:benchmark-design}).
\input{tables/table_obscure_overruled_eval}
\paragraph{Parametric collapse.} Direct LLM safely handled \ObscDirectSafeSummary\ of the obscure queries, with \ObscDirectFatalAff\ answers affirming the overruled rule as current law and \ObscDirectFatalUnw\ reciting it without warning, compared with \ObscCrossDirectFamousSafe\ on the famous Category~D1 subset (descriptive Fisher $p\,\ObscCrossDirectFisherPRel$). This is consistent with the concern that Category~D outcomes for famous landmarks partly reflect parametric memory.
\paragraph{Preregistered primary test.} TERRA-Unlabeled versus Flat Vector RAG was not rejected (discordant pairs $b=\ObscPrimaryDiscordantB$, $c=\ObscPrimaryDiscordantC$; exact McNemar $p\,\ObscPrimaryMcNemarPRel$; gate \ObscPrimaryGate), so the fixed-sequence procedure leaves all secondary contrasts unconfirmed.
\paragraph{Source of the advantage.} Descriptively, Full TERRA exceeded TERRA-Unlabeled ($b=\ObscSecFullVsUnlabeledB$, $c=\ObscSecFullVsUnlabeledC$; exact $p\,\ObscSecFullVsUnlabeledPRel$). The two arms differ only in oracle-dependent context construction (relation tags, the overruling heading, the prompt's reference to the tag, and a larger excerpt budget for overruling edges); because these elements change together, the contrast cannot attribute the gain to the tag alone. We read the result as showing that TERRA's safety on obscure precedents depends on curated overruling annotations available to the system, not that unlabeled traversal fails: with $N=\ObscCohortSize$ the study cannot distinguish TERRA-Unlabeled from Flat Vector RAG.
\paragraph{Robustness.} The magnitude depends on the judge: Judge~1 scored Full TERRA at \ObscJudgeOneFullSafePct\ and TERRA-Unlabeled at \ObscJudgeOneUnlabeledSafePct\ (primary $p\,\ObscJudgeOnePrimaryPRel$), Judge~2 scored \ObscJudgeTwoFullSafePct\ and \ObscJudgeTwoUnlabeledSafePct\ (primary $p\,\ObscJudgeTwoPrimaryPRel$), and the combined severity rule tracks the stricter judge. The primary contrast is non-significant under both. Excluding items with empty TERRA context ($N=\ObscSubsetEmptyContextN$) gives primary $p\,\ObscSubsetEmptyContextPRel$ and Full TERRA \ObscSubsetEmptyContextFullSafePct; excluding items whose context lacks the overruler ($N=\ObscSubsetOverrulerMissingN$), $p\,\ObscSubsetOverrulerMissingPRel$ and \ObscSubsetOverrulerMissingFullSafePct; explicit-overrule items only ($N=\ObscSubsetExplicitOverruleN$), $p\,\ObscSubsetExplicitOverrulePRel$ and \ObscSubsetExplicitOverruleFullSafePct.
\paragraph{Cross-suite contrast (descriptive).} Full TERRA's safe handling was \ObscCrossFullObscureSafe\ on obscure precedents versus \ObscCrossFullFamousSafe\ on the famous subset (risk difference \ObscCrossFullRiskDiff, 95\% CI \ObscCrossFullNewcombeCI); Direct LLM and TERRA-Unlabeled fell from \ObscCrossDirectFamousSafe\ and \ObscCrossUnlabeledFamousSafe\ to \ObscCrossDirectObscureSafe\ and \ObscCrossUnlabeledObscureSafe. With $N=\ObscFamousCohortSize$ and $N=\ObscCohortSize$ these contrasts are descriptive; a non-significant difference does not establish invariance. The famous-subset rates use a stricter dual-judge rubric on nineteen items and are not comparable to the \TerraCatDSafeHandlingRate\ safe-handling rate of Table~\ref{tab:refusal_breakdown}.
\paragraph{Reliability.} Inter-judge agreement was Cohen's $\kappa=\ObscKappa$ over \ObscKappaCleanEvals\ judgments; evidence quotes could not be verified verbatim for \ObscQwenUnverifiedQuotes\ (Judge~1) and \ObscGptOssUnverifiedQuotes\ (Judge~2) judgments, with labels retained. Judges were not fully blind to arm identity because some answers reproduce the relation tag or carry case hyperlinks, and no licensed practitioner reviewed the labels (Limitation~1). A blinded spot-check by one author (not a licensed practitioner) of \ObscSpotN\ items stratified by judge disagreement, primary-test discordance and random selection matched the combined class on \ObscSpotAgreeMatch\ of \ObscSpotAgreeN\ items where the judges agreed but on only \ObscSpotDisagreeMatch\ of \ObscSpotDisagreeN\ where they disagreed. All \ObscSpotMismatchN\ combined-class mismatches rated the answer safer than the severity rule did (Full TERRA \ObscSpotMismatchFull, Flat Vector RAG \ObscSpotMismatchFlat, Direct LLM \ObscSpotMismatchDirect, TERRA-Unlabeled \ObscSpotMismatchUnlabeled), so the rule is conservative but not uniformly across arms, which may affect arm contrasts; the check is not a validation of either judge.

\subsection{External Benchmark Generalization and Anti-Circularity Audit (CaseFacts)}
\label{subsec:casefacts-generalization}

While the canonical \numUniqueQueries-query benchmark rigorously probes multi-hop temporal evolution within indexed Supreme Court jurisprudence, evaluating exclusively on researcher-constructed queries introduces potential risks of circular alignment or latent benchmark overfitting. To establish out-of-distribution empirical validity, we evaluate TERRA against an independent, third-party benchmark: \textbf{CaseFacts} \cite{putta2026casefactsbenchmarklegalfactchecking}, a legal fact-checking dataset constructed from U.S. Supreme Court decisions by an external research team.

\subsubsection{Anti-Circularity and Graph Substrate Audit}
To formally falsify the hypothesis of circular alignment between graph construction and benchmark evaluation, we executed an exhaustive topological audit of the 20,000-opinion Event Evolution Graph:
\begin{itemize}
    \item \textbf{Curated Backbone Seed Density vs. Canonical Benchmark Footprint:} We draw a fundamental scientific distinction between corpus-construction seed density and benchmark query footprint. During offline graph construction, exactly $34$ landmark decisions out of the $20,000$-node corpus ($\CaseFactsSeedDensity$, $\CaseFactsSeedDensityRaw$) received curated procedural reasoning traces. In contrast, the canonical 350-query benchmark touches a total footprint of $284$ distinct nodes ($\CaseFactsFootprintDensity$, $\CaseFactsFootprintNodes$) across its target cases and multi-hop paths (100 in Category A, 151 in Category B, 0 in Category C, and 58 in Category D, with 25 nodes of multi-category overlap). Consequently, $\mathbf{\CaseFactsUntouchedSubstrate}$ of the graph substrate ($\CaseFactsUntouchedNodes$ nodes) remained completely unvisited during canonical evaluation.
    \item \textbf{External Precedent Curation Provenance:} The 19 curated landmark overrulings encoded in the backbone are a strict subset ($19 \subset 56$) of the landmark overruling decisions cataloged by the official Congressional Research Service (CRS) \textit{Constitution Annotated}, confirming that curation reflects established constitutional history rather than benchmark-fitted selections.
    \item \textbf{Chain-Length Dependence:} The Category~B gold chains span 3 to 5 nodes (2 to 4 citation hops): \CatBQueriesLenThree\ chains have 3 nodes, \CatBQueriesLenFour\ have 4 and \CatBQueriesLenFive\ have 5. Full-path recovery by Full TERRA falls with chain length (\CatBTerraPathRecallLenThree, \CatBTerraPathRecallLenFour\ and \CatBTerraPathRecallLenFive\ respectively), so performance is not invariant across chain lengths.
\end{itemize}

\subsubsection{Standardized External Splits and Protocol Parity}
From the complete CaseFacts repository, we established two standardized, auditable evaluation splits:
\begin{enumerate}
    \item \textbf{Track 1: In-Corpus Precedent Classification Benchmark ($N=\CaseFactsTrackOneN$):} Comprises 67 Supported claims (true historical holdings), 45 Refuted claims (inverted or false holdings), and 8 Overruled claims (doctrinally invalidated precedents) anchored in pre-1980 decisions ($1850 \le T \le 1980$).
    \item \textbf{Track 2: Dedicated Modern Refusal Control ($N=\CaseFactsTrackTwoN$):} Comprises 50 modern Supreme Court decisions decided between 1982 and 2023 (e.g., \textit{Citizens United v. FEC}, \textit{United States v. Windsor}, \textit{Fulton v. City of Philadelphia}, \textit{Epic Systems Corp. v. Lewis}), rigorously testing whether systems exercise epistemic refusal on out-of-corpus law.
\end{enumerate}

All 510 evaluations ($360$ Track 1, $150$ Track 2 across Direct LLM, Flat Vector RAG, and Full TERRA) were executed under strict payload parity ($T=0.0$, max tokens 4096, identical LLaMA 3.3 70B generator) and evaluated by independent, non-LLaMA dual judges (\texttt{qwen3-next-80b-a3b-instruct} and \texttt{gpt-oss-120b}). Evaluators achieved $100.0\%$ concordance on parsed classification labels with certified zero LLaMA fallback contaminations.

\input{tables/table_casefacts_eval}

\subsubsection{Empirical Findings: In-Corpus Precedent Verification (Track 1)}
As reported in Table~\ref{tab:casefacts_eval} (Panel A), Full TERRA demonstrates decisive generalization advantages on external precedent verification:
\begin{itemize}
    \item \textbf{Macro-$F_1$ Gains:} Full TERRA achieves $\mathbf{\CaseFactsTerraMacroFOne}$ Macro-$F_1$, outperforming Flat Vector RAG ($\CaseFactsFlatMacroFOne$, a $\CaseFactsTerraLiftFlat$ lift) and Direct LLM ($\CaseFactsDirectMacroFOne$, a $\CaseFactsTerraLiftDirect$ lift). This advantage persists despite TERRA's NLI entailment gate converting 25 borderline in-corpus queries into out-of-corpus safety refusals when retrieved context was legally incomplete.
    \item \textbf{Overruled Class Diagnostic Recall ($n=8$):} In strict compliance with Fix 1, the Overruled class is evaluated as a diagnostic finding with exact Clopper-Pearson 95\% confidence intervals rather than an unconstrained headline metric. Full TERRA correctly identifies $\mathbf{\CaseFactsTerraOverruled}$ Overruled cases (accuracy $50.0\%$, 95\% CI: $\CaseFactsTerraOverruledCI$), double the recall of Flat Vector RAG ($\CaseFactsFlatOverruled$, 95\% CI: $\CaseFactsFlatOverruledCI$) and Direct LLM ($\CaseFactsDirectOverruled$, 95\% CI: $\CaseFactsDirectOverruledCI$).
    \item \textbf{Topological Edge Nuance and Non-Recall Audit:} Graph verification confirms that 2 of the 8 pairs (\textit{Wolf $\to$ Mapp}, \textit{Olmstead $\to$ Katz}) possess direct \texttt{OVERRULES} edges in the EEG, while 6 pairs reflect in-graph citation paths and doctrinal succession. Forensic auditing of the 4 non-recalled Overruled cases demonstrates that none resulted from citation fabrication or hallucination: 2 were safely intercepted by the NLI entailment gate (\textit{Plessy}, \textit{Schenck}) due to incomplete contextual entailment; 1 was a defensible label discordance under 5/5 faithful citations (\textit{Cicenia v. Lagay}, where the model accurately proved the proposition was refuted by \textit{Miranda}); and 1 was an EEG retrieval depth limitation (\textit{Chimel v. California}, where the modifying overruling authority \textit{United States v. Robinson} lay beyond the 2-hop traversal boundary; Section~\ref{subsec:limitations}).
\end{itemize}

\subsubsection{Empirical Findings: Modern Refusal Control and Hallucination Suppression (Track 2)}
On Track 2 (Table~\ref{tab:casefacts_eval}, Panel B), Full TERRA establishes complete epistemic safety on out-of-corpus jurisprudence:
\begin{itemize}
    \item \textbf{100\% Out-of-Corpus Discrimination:} TERRA correctly classifies $\mathbf{\CaseFactsTerraRefusalTrackTwo}$ ($100.0\%$) of modern cases as out-of-corpus, outperforming Flat Vector RAG ($\CaseFactsFlatRefusalTrackTwo$, $92.0\%$) and Direct LLM ($\CaseFactsDirectRefusalTrackTwo$, $80.0\%$).
    \item \textbf{Dual Refusal Modality:} Of the 50 queries, $\CaseFactsTerraHardRefusalTrackTwo$ ($92.0\%$) were short-circuited directly by the NLI entailment gate into safety refusals, while $\CaseFactsTerraDisclosedTrackTwo$ ($8.0\%$) were fully disclosed historical analyses where the model explicitly noted that the target decision was post-1980, evaluated pre-1980 doctrine without anachronism, and concluded with \texttt{OUTSIDE PRECEDENT CORPUS}.
    \item \textbf{Zero Hallucination / Zero Phantom Citations:} Forensic examination of all 50 Track 2 completions confirmed a certified $\mathbf{\CaseFactsTerraHallucinationTrackTwo}$ ($0.0\%$) hallucination rate, with zero phantom citations and zero fabricated precedents, whereas Direct LLM hallucinated substantive answers on $\CaseFactsDirectHallucinationTrackTwo$ ($20.0\%$) and Flat RAG on $\CaseFactsFlatHallucinationTrackTwo$ ($8.0\%$).
\end{itemize}

\subsection{Standardized External Benchmark Evaluation on CLERC (NAACL 2025)}
\label{subsec:clerc-generalization}

While the internal \numUniqueQueries-query benchmark rigorously isolates specific algorithmic mechanisms (Sections~\ref{subsec:safety-and-factual}--\ref{subsec:ablation-study}) and CaseFacts establishes external fact-checking and overruling recognition (Section~\ref{subsec:casefacts-generalization}), evaluating on researcher-crafted queries leaves open whether the architecture generalizes to standardized legal retrieval and generative drafting benchmarks designed by independent natural language processing teams. To address this question, we evaluate TERRA on a stratified sample of \textbf{CLERC} (\textit{Case Law Evaluation and Retrieval Corpus}; Findings of NAACL 2025) \cite{hou2024clerc}, a standardized legal benchmark constructed by Johns Hopkins University and collaborating institutions.

\subsubsection{Four-Layer Scientific Evaluation Architecture}
\label{subsubsec:four-layer-eval}
To establish methodological clarity and prevent confounding between internal mechanism isolation, standardized retrieval benchmarks, and external generalization scope, our empirical validation adheres to a structured four-layer evaluation architecture:
\begin{enumerate}
    \item \textbf{Layer 1: Internal Controlled Benchmark ($N=\numUniqueQueries$):} Isolates fine-grained algorithmic mechanisms (temporal citation graphs, procedural reasoning traces, NLI entailment gate) across four controlled experimental categories (Sections~\ref{subsec:safety-and-factual}--\ref{subsec:ablation-study}).
    \item \textbf{Layer 2: Standardized Precedent Retrieval and Generative Legal Drafting (CLERC Sample, NAACL 2025):} Evaluates general-purpose precedent case retrieval (Task 1A) and unredacted judicial analysis generation (Task 1B) on external, uncurated judicial opinions from a stratified sample of CLERC ($N=120$), complemented by an engineered out-of-corpus refusal stress test ($N=50$).
    \item \textbf{Layer 3: External Temporal Fact-Checking and Overruling Recognition (CaseFacts, ACL 2026):} Evaluates temporal proposition classification, doctrinal overruling recognition, and anti-circularity under third-party adjudication ($N=120$ in-corpus, $N=50$ modern refusal controls; Section~\ref{subsec:casefacts-generalization}).
    \item \textbf{Layer 4: Licensed Legal Practitioner Evaluation:} Parked as essential future work (Section~\ref{subsec:limitations}, Limitation~1) due to resource constraints; all evaluations at Layers 1--3 rely on reproducible algorithmic and non-LLaMA dual-judge protocols.
\end{enumerate}

\subsubsection{Benchmark Protocol, Provenance Hardening, and Anti-Leakage Audit}
\label{subsubsec:clerc-protocol}
The evaluation protocol is divided into two distinct tracks:
\begin{itemize}
    \item \textbf{Track 1: Standardized CLERC Precedent Retrieval and Generative Drafting Sample ($N=\ClercNQueries$):} To evaluate on external common-law contexts while maintaining reproducible execution, we extracted a stratified sample of 120 authentic pre-1980 legal drafting contexts from the 1,498 candidate instances in the CLERC test split. The sample was drawn pre-hoc (prior to observing any model outputs) with a fixed random seed (\texttt{SEED=42}), balanced equally between 60 Direct Precedent Citation queries and 60 Indirect Synthesis queries across five historical eras (1850--1880: $11.7\%$; 1881--1910: $20.0\%$; 1911--1940: $26.7\%$; 1941--1960: $21.7\%$; 1961--1980: $20.0\%$).
    \item \textbf{Track 2: TERRA-Specific Modern and Cross-Jurisdictional Refusal Stress Test ($N=\ClercNRefusal$):} Rather than a native CLERC benchmark task, Track 2 serves as a dedicated architectural safety stress test comprising 50 negative control queries (39 post-1980 modern Supreme Court decisions and 11 non-SCOTUS federal circuit/statutory queries) to rigorously evaluate epistemic boundary enforcement and anti-hallucination guardrails under out-of-corpus pressure.
\end{itemize}

\paragraph{Provenance Audit and Seed Purge}
To guard against latent circularity, we executed an automated cryptographic provenance audit. Out of 120 Track 1 target cases, exactly zero ($0.0\%$, 0/120) overlap with the 34 hand-curated landmark seed nodes that received bespoke reasoning traces during graph construction; all overlapping queries were purged during split generation. Track 1 evaluates on a certified $100.0\%$ uncurated graph substrate. Crucially, this exclusion specifically guarantees that evaluation is not artificially aided by the bespoke reasoning-trace seeds created during system development; it does not preclude legitimate retrieval over ordinary authentic Supreme Court opinions in the 20,000-opinion corpus, which serves as the uniform, shared knowledge substrate for all evaluated retrieval systems. As a documented limitation and scope distinction, the broader corpus-wide substantive holding extraction cascade (Tier~1 syllabi and Tier~2 headnotes covering $51.1\%$ of the 20,000 opinions) was not separately screened, as it constitutes the uniform, automated knowledge substrate available to all in-corpus GraphRAG retrievers.

\paragraph{Pre-Flight Gate Smoke Test}
Following the fatal-interception safeguards developed in prior evaluation phases, the batch harness enforces a mandatory pre-flight assertion over three known-refusal probes (Probe 1: Modern Post-1980 Authority; Probe 2: Fabricated Holding Trap; Probe 3: Non-SCOTUS Circuit Rule). The harness verifies that the NLI Smart Grader successfully short-circuits execution before launching the full evaluation batch.

\paragraph{Independent Non-LLaMA Dual Judging}
All completions across Direct LLM, Flat Vector RAG, and Full TERRA were executed under strict payload parity ($T=0.0$, max tokens 4096, identical LLaMA 3.3 70B generator) and judged by independent, non-LLaMA dual evaluators (\texttt{qwen3-next-80b-a3b-instruct} and \texttt{gpt-oss-120b}). Exact citation metrics (Citation Precision and Citation Recall) were computed via deterministic parse-tree resolution using \texttt{eyecite}.

\input{tables/table_clerc_eval}

\subsubsection{Empirical Findings: Precedent Retrieval and Analysis Generation (Track 1)}
As detailed in Table~\ref{tab:clerc_eval} (Panel A), Full TERRA demonstrates decisive structural advantages across retrieval and generation while surfacing a clear architectural boundary in open-ended generative drafting:
\begin{itemize}
    \item \textbf{Precedent Retrieval Comparison (Task 1A):} On the evaluated CLERC sample, Full TERRA achieved slightly higher Recall@10 ($\ClercTerraRecallTen$ vs. $\ClercFlatRecallTen$) and nDCG@10 ($\ClercTerraNDCG$ vs. $\ClercFlatNDCG$) than Flat Vector RAG, alongside an MRR@10 of $\ClercTerraMRR$ vs. $\ClercFlatMRR$ for dense retrieval. While directional improvements are observed from chronological citation traversal and holding re-ranking, the absolute margin on single-turn benchmark queries remains modest ($\ClercTerraRecallTenDelta$ percentage points Recall@10).
    \item \textbf{Generative Analysis and Lexical Fidelity (Task 1B):} In judicial analysis drafting, TERRA attains $\ClercTerraRougeL$ ROUGE-L and $\ClercTerraBleu$ BLEU-4, reflecting structural adherence to judicial reasoning styles without degenerating into verbatim repetition. Deterministic \texttt{eyecite} parse-tree extraction demonstrates that TERRA attains the highest Citation Precision among evaluated systems ($\text{CP} = \ClercTerraCitePrecision$, vs. $\ClercFlatCitePrecision$ for Flat Vector RAG and $\ClercDirectCitePrecision$ for Direct LLM) alongside a Citation Recall of $\text{CR} = \ClercTerraCiteRecall$ and Citation $F_1$ of $\ClercTerraCiteFOne$ (compared to Flat RAG's $\ClercFlatCiteRecall$ and Direct LLM's $\ClercDirectCiteRecall$, where Direct LLM's higher recall is driven by prolific ungrounded citation dumping).
    \item \textbf{Zero Case Caption Fabrication:} Forensic auditing of all citations extracted across the 120 Track 1 completions confirms that no fabricated case captions or phantom citations were observed in the audited TERRA outputs ($0/120$, $\ClercTerraFabricationTrackOne$). Every citation generated by TERRA corresponded to a verified historical Supreme Court authority ($\ClercTerraInCorpusCitePct$) or an authentic external context case ($\ClercTerraExternalCitePct$). When the NLI entailment gate fired to intercept drafting queries ($53/120$ cases), forensic auditing confirmed zero fabricated case citations among the short-circuited completions ($0/53$). By contrast, unaugmented Direct LLM suffered a $\ClercDirectFabricationTrackOne$ phantom citation rate, frequently inventing non-existent case captions and reporter volumes.
    \item \textbf{Disclosed Limitation: Doctrinal Misapplication Under Retrieval Horizon Truncation:} However, when generation proceeded on incomplete retrieval subgraphs ($67/120$ cases), independent dual judging flagged a doctrinal misapplication rate of $\ClercTerraMisapplicationTrackOne$ ($\ClercTerraMisapplicationOrTrackOne$ under Logical OR union, $\ClercTerraMisapplicationAndTrackOne$ under Unanimous AND consensus; compared to $\ClercFlatMisapplicationTrackOne$ for Flat Vector RAG and $\ClercDirectMisapplicationTrackOne$ for Direct LLM). Case-by-case auditing reveals that these errors are not phantom fabrications, but instances of doctrinal misapplication where the generator attempted to reason by analogy from retrieved historical authorities when the gold controlling precedent lay beyond the 2-hop traversal radius. This establishes a fundamental empirical principle: \textit{authentic evidence does not equate to sufficient evidence}. Grounding generations on authentic graph subgraphs completely eliminates phantom case inventions, but generative reasoning remains vulnerable to doctrinal error whenever relevant authority is truncated by the graph retrieval horizon. This finding directly motivates our in-depth examination of precedent-chain path recovery and multi-hop retrieval depth in Section~\ref{subsec:retrieval-quality}.
\end{itemize}

\subsubsection{Empirical Findings: Refusal Control on Negative Probes (Track 2)}
On Track 2 (Table~\ref{tab:clerc_eval}, Panel B), Full TERRA establishes robust epistemic boundary discipline on external out-of-corpus queries:
\begin{itemize}
    \item \textbf{Total Out-of-Corpus Discrimination:} TERRA achieves $\mathbf{\ClercTerraRefusalTrackTwo}$ total safe refusal on out-of-corpus queries, outperforming Flat Vector RAG ($\ClercFlatRefusalTrackTwo$) and Direct LLM ($\ClercDirectRefusalTrackTwo$) by a $3.3\times$ margin. We emphasize that Track 2 functions specifically as an engineered epistemic boundary and refusal stress test rather than an overall task utility metric. While high safe refusal ($\ClercTerraRefusalTrackTwo$) demonstrates effective containment of out-of-corpus queries, real-world deployment requires balancing safe refusal with retrieval coverage.
    \item \textbf{Hard NLI Interception and Gate Safety:} The NLI entailment gate successfully intercepted $\mathbf{\ClercTerraHardRefusalTrackTwo}$ of out-of-corpus queries prior to generation. Forensic auditing confirmed zero fabricated case citations among the 86 gate-triggering events across both tracks (0/53 in Track 1, 0/33 in Track 2), confirming that gate interception reliably halts generation before unverified citations can be introduced.
    \item \textbf{Fabrication vs. Misapplication on Negative Probes:} Full TERRA recorded a certified $\mathbf{\ClercTerraFabricationTrackTwo}$ case caption fabrication rate on Track 2 (zero phantom case names across all 50 negative controls), compared to $\ClercDirectFabricationTrackTwo$ citation fabrication in Direct LLM. On the remaining 17 cases where generation proceeded without gate short-circuit (representing \ClercTerraMisapplicationOrTrackTwo\ across all 50 Track-2 queries; 100.0\% of generated completions), both judges unanimously flagged doctrinal misapplication ($\ClercTerraMisapplicationTrackTwo$), as the model attempted to answer modern post-1980 or statutory questions using historical pre-1980 SCOTUS authorities.
\end{itemize}

\subsection{Direct Evidence Retrieval Quality and Precedent-Chain Coverage}
\label{subsec:retrieval-quality}

A critical question in retrieval-augmented legal architectures is whether downstream generative advantages stem from superior evidence retrieval or merely reflect the generator's ability to synthesize plausible analyses from average context. Furthermore, Second Review Issue~\#11 noted that while TERRA's canonical graph expansion operates with breadth $k=2$ and depth $d=2$, Category~B queries comprise 3-to-5 node evolutionary arcs (spanning 2 to 4 citation hops), requiring explicit evaluation of retrieval success as a function of precedent chain length.

The core judicial benchmark encompasses $N=350$ queries across four operational strata: Category~A ($N=100$ single-point factual lookups), Category~B ($N=100$ multi-hop evolutionary chains), Category~C ($N=75$ out-of-context unanswerable refusal probes where target authorities are absent from the corpus, evaluated for gate interception in Section~\ref{subsec:safety-and-factual}), and Category~D ($N=75$ adversarial temporal overruling pairs). To rigorously isolate evidence-retrieval quality on complex multi-hop common-law reasoning without conflating single-point lookups (which TERRA intentionally routes through its parametric Fast Path) or unanswerable negative probes (where no target nodes exist in the corpus), we evaluate raw evidence retrieval on the $N=\CatBRetrievalQueries$ Category~B multi-hop trajectories ($\CatBGoldPrecedents$ gold landmark precedents) prior to generator invocation.

Because comparative architectures deliver heterogeneous retrieval representations---Flat Vector RAG returns a linear ranked list of document chunks, GraphRAG delivers unranked modularity community reports, HippoRAG yields unranked Personalized PageRank (PPR) activation clusters, and TERRA produces a 2-hop directed EEG subgraph structured by temporal edges and holding headnotes---linear ranked-list metrics such as MRR or nDCG cannot be applied uniformly across community and subgraph representations without imposing arbitrary ranking assumptions. We therefore evaluate deterministic dual-resolution entity recall (exact CAP Node ID and canonical case title) and full precedent trajectory recovery in Table~\ref{tab:retrieval_eval}. All retrieval metrics are evaluated directly on the canonical stored \texttt{context\_full} fields from the canonical v2 evaluation file (\texttt{terra\_canonical\_v2\_eval.json}, SHA256 prefix \texttt{95e775e2}) that were delivered to the LLaMA-3.3-70B-Instruct generator in the primary dual-judge benchmark, so that retrieval is scored on the context the generator actually received.

\input{tables/table_retrieval_eval}

\subsubsection{Overall Precedent-Chain Recovery Performance}
As reported in Table~\ref{tab:retrieval_eval} (Panel A), comparative retrievers exhibit fundamental divergences in their ability to capture multi-hop common-law trajectories:
\begin{itemize}
    \item \textbf{Dense Semantic Search Horizon Truncation:} Flat Vector RAG achieves only $\CatBFlatMacroNodeRecall$ macro node recall ($\CatBFlatMicroNodeRecall$ micro) and captures the complete precedent trajectory on only $\CatBFlatFullPathRecall$ of queries ($\CatBFlatFullPathCount$). While dense embeddings successfully identify one or two lexically adjacent landmark cases, they almost completely fail to recover intermediate stepping-stone authorities across multi-decade evolutionary spans.
    \item \textbf{Static Community and Undirected Diffusion Limitations:} GraphRAG-Adapted achieves $\CatBGraphRAGMacroNodeRecall$ node recall and $\CatBGraphRAGFullPathRecall$ full-path recovery ($\CatBGraphRAGFullPathCount$), bounded by the fact that static Louvain modularity communities group cases topically without preserving chronological precedence. HippoRAG-Adapted achieves substantially higher node recall ($\CatBHippoRAGMacroNodeRecall$), but attains only $\CatBHippoRAGFullPathRecall$ full-chain path recall ($\CatBHippoRAGFullPathCount$). Because Personalized PageRank diffuses undirected probability mass across dense citation clusters, it frequently activates high-degree peripheral precedents rather than following the strict temporal citation spine.
    \item \textbf{Directional Precedent Spine Traversal:} Full TERRA achieves $\CatBTerraMacroNodeRecall$ macro node recall ($\CatBTerraMicroNodeRecall$ micro) and recovers the complete evolutionary chain on $\CatBTerraFullPathCount$/100 queries, a $\CatBTerraVsFlatNodeRecallDelta$ node-recall lift and $\CatBTerraVsFlatPathRecallFactor$ path-recovery factor over Flat Vector RAG (which recovers $\CatBFlatFullPathCount$/100), and $\CatBTerraVsHippoNodeRecallDelta$ and $\CatBTerraVsHippoPathRecallFactor$ over HippoRAG. By traversing chronologically directed citation edges and re-ranking via substantive holding headnotes, TERRA preserves the continuous chain of legal authority necessary for sound doctrinal synthesis.
\end{itemize}

\subsubsection{Chain-Length Disaggregation and Hop Analysis}
Table~\ref{tab:retrieval_eval} (Panel B) disaggregates precedent discovery across evolutionary chain lengths, addressing Second Review Issue~\#11:
\begin{itemize}
    \item \textbf{3-Node Chains ($N=\CatBQueriesLenThree$, 2 Citation Hops, 1 Intermediate Node):} On 2-hop chains (e.g., \textit{Plessy} $\to$ \textit{Gaines} $\to$ \textit{Brown}), Flat Vector RAG captures the full chain in only $\CatBFlatPathRecallLenThree$ of queries ($\CatBFlatPathCountLenThree$), GraphRAG achieves $\CatBGraphRAGPathRecallLenThree$ ($\CatBGraphRAGPathCountLenThree$), and HippoRAG achieves $\CatBHippoRAGPathRecallLenThree$ ($\CatBHippoRAGPathCountLenThree$). Full TERRA achieves $\CatBTerraNodeRecallLenThree$ node recall and $\CatBTerraPathRecallLenThree$ full path recovery ($\CatBTerraPathCountLenThree$), confirming that $d=2$ graph expansion reliably spans two-hop doctrinal progressions.
    \item \textbf{4-Node Chains ($N=\CatBQueriesLenFour$, 3 Citation Hops, 2 Intermediate Nodes):} When doctrinal arcs span 3 citation hops (e.g., \textit{Strauder} $\to$ \textit{Neal} $\to$ \textit{Norris} $\to$ \textit{Batson}), baseline path recovery drops to zero for both Flat Vector RAG ($\CatBFlatPathRecallLenFour$, $\CatBFlatPathCountLenFour$) and GraphRAG ($\CatBGraphRAGPathRecallLenFour$, $\CatBGraphRAGPathCountLenFour$), and collapses to $\CatBHippoRAGPathRecallLenFour$ ($\CatBHippoRAGPathCountLenFour$) for HippoRAG. In contrast, Full TERRA sustains $\CatBTerraNodeRecallLenFour$ node recall and $\CatBTerraPathRecallLenFour$ full path recovery ($\CatBTerraPathCountLenFour$).
    \item \textbf{5-Node Chains ($N=\CatBQueriesLenFive$, 4 Citation Hops, 3 Intermediate Nodes):} On long 4-hop lineages, all baseline retrievers fail completely to recover the full precedent path ($\CatBFlatPathRecallLenFive$, $\CatBGraphRAGPathRecallLenFive$, and $\CatBHippoRAGPathRecallLenFive$; $0/5$ each). Full TERRA achieves $\CatBTerraNodeRecallLenFive$ node recall and $\CatBTerraPathRecallLenFive$ full path recovery ($\CatBTerraPathCountLenFive$). Due to the limited sample size ($N=\CatBQueriesLenFive$, exact Clopper-Pearson 95\% CI: \CatBFiveNodeClopperPearsonCI), this 5-node stratum serves as a qualitative descriptive diagnostic demonstrating topological feasibility rather than a statistically generalized population estimate.
\end{itemize}

\subsubsection{Formal Definition and Seed Positioning of Multi-Seed Anchoring}
The ability of TERRA's nominal $d=2$ graph traversal to achieve complete path recovery on 3-hop and 4-hop lineages is governed by \textit{multi-seed anchoring}. For an evolutionary query $q$, entity and legal concept extraction identifies a set of seed nodes $S_q \subset V$ with $|S_q| \ge 2$ (empirical mean of $\CatBMeanSeedsPerQuery$ seeds per query across Category~B; $\CatBTwoSeedsPct$ with $|S_q|=2$, $\CatBThreeSeedsPct$ with $|S_q|=3$, and $\CatBFourPlusSeedsPct$ with $|S_q| \ge 4$). From each seed $s \in S_q$, directed EEG traversal expands a $d=2$ graph ball $B(s, 2)$ along forward and backward citation edges with breadth $k=2$.

Multi-seed extraction is designed to anchor both extremities of a precedent trajectory, the historical origin (index 0) and the culminating modern authority (index $m-1$). Seed placement was not re-audited for the canonical run, so the following is a mechanism explanation rather than a measured result. For an evolutionary chain $v_0 \to \dots \to v_{m-1}$, the forward 2-hop ball $B(v_0, 2)$ covers $\{v_0, v_1, v_2\}$ while the reverse 2-hop ball $B(v_{m-1}, 2)$ covers $\{v_{m-1}, v_{m-2}, v_{m-3}\}$. When intermediate branching along the precedent spine remains within the $k=2$ beam, the forward and backward neighborhoods intersect at intermediate stepping-stone authorities ($v_2$ in 5-node chains), which can cover 3-hop and 4-hop doctrinal paths without unconstrained search depths ($d \ge 4$). In the canonical run, all gold nodes were recovered for $\CatBTerraPathCountLenFour$ 4-node and $\CatBTerraPathCountLenFive$ 5-node chains; high citation fan-out at intermediate nodes, exceeding the $k=2$ beam width, is a plausible but untested explanation for the 16 incompletely recovered 4-node chains.

These empirical findings confirm that TERRA's high generative win rates (Table~\ref{tab:overall_summary}) and low fabrication rates (Table~\ref{tab:clerc_eval}) are fundamentally anchored in direct evidence retrieval quality rather than ungrounded parametric fluency.

\subsubsection{Retrieval Hyperparameter Robustness and Traversal Horizon Sensitivity}
\label{subsubsec:hyperparameter-sensitivity}
To evaluate retrieval sensitivity across topological expansion horizons, Table~\ref{tab:retrieval_eval} (Panel~C) reports sensitivity sweeps across \IssueTwentySixConfigCount\ configurations ($N=\CatBRetrievalQueries$ Category~B queries) under a dual-reporting protocol:

\textit{1) Historical-Baseline Sensitivity Analysis (Pre-Issue~\#13 Engine):} Evaluating the unconstrained retrieval engine reproduces the published canonical baseline ($k=\IssueTwentyFourSeedK, d=\IssueTwentyFourBfsDepth, C=\IssueTwentyFourHubCap$) with certified 100/100 bit-identical parity across node sets, gold landmark intersections, and aggregate metrics (\IssueTwentySixHistMacro\ macro recall, \IssueTwentySixHistFullPath\ full path recall). Truncating search depth to $d=1$ reduces full-path recovery to \IssueTwentySixHistDepthOnePath\ (\IssueTwentySixHistDepthOneMacro\ macro recall) by prematurely pruning intermediate stepping-stones, while reducing context from \IssueTwentySixHistMeanChars\ to \IssueTwentySixHistDepthOneChars\ characters and offline traversal latency from \IssueTwentySixHistOfflineLatencyMS~ms to \IssueTwentySixHistDepthOneLatencyMS~ms. Expanding depth to $d=3$ yields modest path recovery gains (\IssueTwentySixHistDepthThreePath, 68/100) at exponential context scale (\IssueTwentySixHistDepthThreeChars\ characters, \IssueTwentySixHistDepthThreeLatencyMS~ms). Constraining hub caps to $C=25$ severely restricts multi-hop recovery (\IssueTwentySixHistCapTwentyFivePath\ path recall), whereas relaxing to $C=100$ achieves near-complete path recovery (\IssueTwentySixHistCapOneHundredPath, \IssueTwentySixHistCapOneHundredMacro\ macro recall) at prohibitive context scale (\IssueTwentySixHistCapOneHundredChars\ characters). Adaptive depth ($d_{\text{adapt}}$) matches static $d=2$ closely (\IssueTwentySixHistAdaptiveMacro\ macro, \IssueTwentySixHistAdaptivePath\ path recall). Single-seed initialization ($k=1^\ast$) attains \IssueTwentySixHistKOnePath\ path recovery, confirming that multi-seed boundary anchoring provides essential stepping-stone coverage, while vector-seed expansion ($k=4^\ast$) retrieves zero seeds due to the known active-index category filter constraint (\IssueTwentySixHistKFourPath\ path recall).

\textit{2) Post-Issue~\#13 Bounded Operational Diagnostic:} Enforcing production operational caps ($35{,}000$ structural characters, 40 traversed nodes, asymmetric hub budgets), expanding 2 seed nodes into citation neighborhoods halts traversal after visiting a mean of \IssueTwentySixBoundMeanVisitedNodes\ expansion centers, assembling \IssueTwentySixBoundMeanFullNodes\ unique case authorities into context. The bounded engine preserves \IssueTwentySixGoldRetentionPct\ of gold landmark precedents (\IssueTwentySixGoldRetentionCount) and achieves \IssueTwentySixBoundFullPath\ full path recall (48/63 on 3-node, 12/32 on 4-node, 1/5 on 5-node), eliminating \IssueTwentySixContextReductionPct\ of context character volume (from \IssueTwentySixHistMeanChars\ to \IssueTwentySixBoundMeanChars\ characters) and accelerating offline traversal latency by \IssueTwentySixSpeedupFactor\ (from \IssueTwentySixHistOfflineLatencyMS~ms to \IssueTwentySixBoundOfflineLatencyMS~ms).

\input{tables/table_citation_eval}

\subsection{Generative Citation Correctness, Doctrinal Support, and Interactive Hyperlink Provenance}
\label{subsec:citation-correctness}

While direct evidence retrieval evaluation (Table~\ref{tab:retrieval_eval}) rigorously establishes the retriever's ability to assemble necessary precedent authorities in the prompt context prior to generator invocation, legal validity fundamentally depends on whether the generator accurately cites and attributes doctrine in its final synthesized completions. Large language models frequently commit subtle citation errors---inventing phantom case names, citing authorities for propositions they never established, transposing holdings across contemporaneous landmarks, or hallucinating citations on unanswerable inquiries.

To resolve Issue~\#16, Table~\ref{tab:citation_eval} evaluates end-to-end citation fidelity directly in the synthesized completions across the canonical \numUniqueQueries-query benchmark suite (\numTotalEvals\ comparative evaluations across all \numComparativePipelines\ pipelines). Grounded doctrinal precision and F1 serve as primary metrics (TERRA \TerraDocPrecisionGrounded/\TerraDocFOneGrounded, HippoRAG \HippoDocPrecisionGrounded/\HippoDocFOneGrounded, Flat Vector RAG \FlatDocPrecisionGrounded/\FlatDocFOneGrounded), with the judge-weighted variant as secondary (TERRA \TerraDocPrecisionJudgeWeighted, HippoRAG \HippoDocPrecisionJudgeWeighted; weighting multiplies by the dual-judge faithfulness of each answer). TERRA's entity citation precision is \TerraEntityCitePrecision\ (earlier single-judge run \TerraEntityCitePrecisionLegacy); \CitationLinksResolved\ of \CitationLinksEmitted\ emitted links resolved (\CitationLinkResolutionRate).



\subsubsection{Category B Precedent Milestone Citation and Doctrinal Support}
As reported in Table~\ref{tab:citation_eval} (Panel~A), comparative architectures exhibit substantial divergence in their ability to cite and support multi-hop legal precedents:
\begin{itemize}
    \item \textbf{Precedent Citation Recall Leadership:} On Category~B multi-decade evolutionary trajectories ($N=\CatBRetrievalQueries$ queries, $\CatBGoldPrecedents$ gold landmark authorities), Full TERRA achieves $\mathbf{\CatBTerraMacroCitationRecall}$ macro citation recall ($\CatBTerraMicroCitationRecall$ micro), outperforming HippoRAG-Adapted ($\CatBHippoRAGMacroCitationRecall$ macro), Direct LLM ($\CatBDirectMacroCitationRecall$), Flat Vector RAG ($\CatBFlatMacroCitationRecall$), and GraphRAG-Adapted ($\CatBGraphRAGMacroCitationRecall$) by $\CatBTerraVsHippoCitationRecallDelta$ to $\CatBTerraVsGraphCitationRecallDelta$ absolute percentage points.
    \item \textbf{Doctrinal Support Precision:} When evaluating whether cited authorities factually support the asserted legal claims without holding misattribution, Full TERRA attains $\CatBTerraDoctrinalCitationPrecision$ judge-weighted Doctrinal Citation Precision: higher than HippoRAG-Adapted ($\CatBHippoRAGDoctrinalCitationPrecision$), Direct LLM ($\CatBDirectDoctrinalCitationPrecision$) and GraphRAG-Adapted ($\CatBGraphRAGDoctrinalCitationPrecision$), but lower than Flat Vector RAG ($\CatBFlatDoctrinalCitationPrecision$), which has the highest judge-weighted value.
    \item \textbf{Citation $F_1$:} Balancing precision and recall, Full TERRA's Citation $F_1$ is $\CatBTerraCitationFOne$: higher than HippoRAG-Adapted ($\CatBHippoRAGCitationFOne$; ratio $\CatBTerraVsHippoCitationFOneLift$) and Direct LLM ($\CatBDirectCitationFOne$; ratio $\CatBTerraVsDirectCitationFOneLift$), but lower than Flat Vector RAG ($\CatBFlatCitationFOne$; ratio $\CatBTerraVsFlatCitationFOneLift$).
    \item \textbf{Complete Precedent Path Citation:} On complete precedent trajectories, Full TERRA captures and cites 100\% of the gold landmark precedent milestones in $\mathbf{\CatBTerraFullPathCitationCount}$ answers ($\CatBTerraFullPathCitationRecall$), compared to only $4/100$ for HippoRAG, $2/100$ for Flat Vector RAG, $1/100$ for Direct LLM, and $0/100$ for GraphRAG---a $\mathbf{\CatBTerraVsHippoPathCitationLift}$ advantage over HippoRAG and a $\mathbf{\CatBTerraVsDirectPathCitationLift}$ increase over Direct LLM.
\end{itemize}

\subsubsection{Dynamic Interactive Hyperlink Attribution and Epistemic Negative Containment}
Table~\ref{tab:citation_eval} (Panel~B) profiles the explainable attribution mechanics across the 350-benchmark suite:
\begin{itemize}
    \item \textbf{Interactive Hyperlinked Attribution:} In accordance with Section~\ref{subsubsec:xai-provenance}, Full TERRA dynamically formats referenced opinions as interactive relative markdown links. Across benchmark completions, TERRA emitted $\mathbf{\TerraTotalMarkdownLinks}$ interactive links. Empirical auditing against the 20,000-case EEG index confirms that $\mathbf{\TerraResolvedMarkdownLinks/\TerraTotalMarkdownLinks}$ links ($\mathbf{\TerraMarkdownLinkResolutionRate}$) pass the project's link validator (which accepts a link when its target identifier exists in the index or its caption matches an indexed case title or alias), with $\mathbf{\TerraMarkdownLinkFidelityRate}$ semantic title-to-node fidelity (including authentic historical aliases such as the \textit{Slaughter-House Cases} dynamically resolving to its formal reporter caption \textit{Butchers' Benevolent Ass'n}, \texttt{83\_0036-01}). The remaining $\TerraBrokenMarkdownLinks$ links fail the validator: three carry page-padding errors in the identifiers of cases that are in the corpus (e.g.\ \texttt{301\_001-01} for \texttt{301\_0001-01}, \textit{NLRB v. Jones \& Laughlin}; queries~B073 and~D038), and one is a bare \texttt{/case/118} path (query~B075). A stricter check that requires the identifier itself to exist in the index passes 644 of 664 links ($97.0\%$); the other 16 links with a non-existent identifier were accepted on caption match. In stark contrast, all four baseline architectures emit \textbf{zero} interactive hyperlinks (\BaselinesMarkdownLinks), leaving their citations unlinked and unverifiable without manual docket lookups.
    \item \textbf{Out-of-Corpus Negative Probe Containment (Category C):} On the $N=75$ unanswerable negative probes, under the formal piecewise precision formulation (where emitting zero citations constitutes safe epistemic silence, scored as $1.00$, while hallucinating citations is scored as $0.00$), Flat Vector RAG suffered severe context bleed, hallucinating $\CatCFlatEmittedCitations$ spurious case citations across unanswerable queries ($\CatCFlatCitationPrecision$ precision). In contrast, Full TERRA achieved $\mathbf{\CatCTerraCitationPrecision}$ precision, emitting $\CatCTerraEmittedCitations$ citations across all 75 queries ($\CatCTerraCleanRefusals$ clean refusals). In the deployed run, 74 of these 75 probes had no retrieved context, so this containment largely reflects refusal on empty retrieval rather than entailment grading of retrieved evidence.
\input{tables/table_computational_budget}

\end{itemize}

\subsection{Computational Budget Profiling and Geometric Gate Calibration}
\label{subsec:computational-budget}

A rigorous comparative evaluation of complex agentic RAG frameworks requires auditing the full computational footprint of the underlying architecture. Second Review Issue~\#10 noted that multi-stage RAG pipelines deploy multiple sequential inference calls (query routing, vector retrieval, graph traversal, NLI verification, and final generation), which can compound latency, token consumption, and serving costs. To provide complete transparency and establish call parity against standard single-shot baselines, Table~\ref{tab:computational_budget} profiles total operations per query, active 70B generator invocations, end-to-end and stage latencies, effective token volumes, and modelled financial costs across all \numComparativePipelines\ comparative architectures, alongside an empirical calibration of the Tier~2 geometric pre-filter on the 20,000-case CAP index.



\input{tables/table_temporal_update_stability}

\subsubsection{Sequential Pipeline Operations and Call-Parity Accounting}
Comparative single-shot baselines (Flat Vector RAG, GraphRAG, and HippoRAG) execute a uniform two-stage retrieve-then-generate pipeline ($\FlatRAGTotalOpsPerQuery$ operations per query), uniformly issuing $\FlatRAGGenCallsFormatted$ active 70B generator calls across the $N=\numUniqueQueries$ benchmark suite (Panel~A). In contrast, TERRA implements a multi-stage sequential architecture:
\begin{equation}
\begin{aligned}
q &\xrightarrow{\text{Router}} \text{Tier 2 Gate} \xrightarrow{\text{Retrieve}} \text{EEG BFS} \\
  &\xrightarrow{\text{NLI Grader}} \text{70B Generator}
\end{aligned}
\end{equation}
Across the 350-query benchmark, TERRA performs \TerraOperationsTotalCount\ operations (\TerraOperationsPerQuery\ per query; operations counted where a stage recorded non-zero time) at a modelled \TerraDeployedCostPerThousand\ per 1,000 queries under the deployed configuration (Gemma~3 12B router; Llama~3.3 70B NLI grader and generator). The cost is modelled from token-count estimates at list prices, not metered; it prices a grader call for every Category~B--D query (twice for refused Category~C and~D queries) and none for Category~A, although Category~A recorded grading time. Pricing one grader call only where the recorded grading time was at least 1~ms gives \$4.66 per 1,000 queries instead.

These counts are measured from recorded stage times, not taken from a design model. In the canonical run, Category~A (the nominal fast path) recorded retrieval and NLI-grader time on all 100 queries and invoked the generator on \TerraCatAGenCallsFormatted; Category~B recorded all four stages on all 100 queries and invoked the generator on \TerraCatBGenCallsFormatted; Category~C recorded routing, retrieval and grading time on all 75 queries but no generation, because all 75 were refused, and for 74 of them the recorded grading time was under 1~ms and the retrieved context was empty, which is consistent with a deterministic empty-context rejection rather than an LLM grading call; Category~D recorded all three pre-generation stages on all 75 queries and invoked the generator on \TerraCatDGenCallsFormatted, the other 46 queries being refused (7 of them with an empty context). This yields \TerraCatATotalOps, \TerraCatBTotalOps, \TerraCatCTotalOps\ and \TerraCatDTotalOps\ operations on Categories~A--D (\TerraCatATotalOpsPerQuery, \TerraCatBTotalOpsPerQuery, \TerraCatCTotalOpsPerQuery\ and \TerraCatDTotalOpsPerQuery\ per query) and \TerraBenchmarkTotalOpsFormatted\ in total (\TerraOverallTotalOpsPerQuery\ per query). A design-level model in which out-of-domain queries beyond the geometric boundary are terminated before retrieval is not what the deployed run did: every Category~C query recorded retrieval time (at least 178~ms excluding routing), so the threshold analysis in Panel~B describes an offline calibration of a candidate pre-filter.

Importantly, while TERRA executes more sequential pipeline stages than single-shot baselines, it invokes the generator on only \TerraOverallGenCallsCount\ of the 350 benchmark queries (\TerraOverallGenCallsPct), a \TerraOverallGenCallsReductionPct\ reduction in active generator calls compared to the single-shot baselines ($350/350$). The reduction comes from refusals: Category~C made no generator calls (\TerraCatCGenCallsFormatted), Category~D made \TerraCatDGenCallsFormatted, and the refusals on Categories~A and~B also skip generation. It is not a reduction in total 70B-model use: the NLI grader is also a Llama~3.3 70B model and recorded at least 1~ms of grading time on 276 of the 350 queries, so TERRA's generator and grader calls together exceed the baselines' one call per query. We frame the generator-call reduction not as an alternative metric for call parity, but as one mitigating factor in aggregate serving efficiency.

\subsubsection{Candid Disclosure of Deliberate Latency and Cost Trade-Offs}
Operating a multi-stage sequential architecture entails a deliberate engineering trade-off between computational budget and epistemic reliability. On complex multi-hop evolutionary chains (Category~B), TERRA exhibits a median pipeline latency of $\TerraHardPathMedianSeconds\text{ s}$ ($\TerraCatBLatencyMeanMS\text{ ms}$ mean) and a modelled cost of $\TerraCatBTotalCostPerThousand$ per 1,000 queries ($\TerraCatBPerQueryCost$ per query), compared to $\FlatRAGLatencyMedMS\text{ ms}$ and $\FlatRAGTotalCostPerThousand$ per 1,000 queries ($\FlatRAGPerQueryCost$ per query) for Flat Vector RAG---representing a $\TerraCatBVsFlatCostRatio$ per-query cost ratio.

We candidly disclose this computational overhead as an intentional, necessary investment to resolve the structural vulnerabilities of dense semantic retrieval. As established in Section~\ref{subsec:retrieval-quality}, Flat Vector RAG achieves only $\CatBFlatFullPathRecall$ full precedent path recovery ($\CatBFlatFullPathCount$), collapsing completely on three-hop and four-hop legal trajectories. In contrast, TERRA's coordinated EEG citation traversal and procedural trace synthesis attains $\CatBTerraFullPathRecall$ complete precedent path recovery ($\CatBTerraFullPathCount$), delivering a $\CatBTerraVsFlatPathRecallFactor$ increase in trajectory recovery over Flat RAG. In high-stakes legal applications where ungrounded generative hallucinations or missing precedent stepping-stones induce malpractice liability, expending $\sim \TerraHardPathMedianSeconds\text{ s}$ and roughly $\TerraCatBPerQueryCost$ per query (modelled) to recover the full precedent path on $\CatBTerraFullPathCount$ of the Category~B queries is, in our judgement, a defensible trade-off.

\subsubsection{Stage Latency Profiling}
Panel~A reports mean stage times from the logged timing fields. For hard-route queries (Categories~B--D, 240 of 350) the logged retrieval timer is cumulative and includes routing time (the two correlate at $r=0.999$ across those queries, and the difference is always at least 178~ms), so retrieval is taken as logged retrieval time minus routing time; for fast-route queries the timers are separate. NLI grading and generation are the logged values. The stages sum exactly to the end-to-end time for Category~A and the single-shot baselines. For Categories~B--D they leave a remainder that is not attributed to any stage and whose cause was not established (mean $\TerraCatBRemainderMS\text{ ms}$, $\TerraCatCRemainderMS\text{ ms}$ and $\TerraCatDRemainderMS\text{ ms}$, respectively), so the columns do not reconcile exactly to the end-to-end means:
\begin{enumerate}
    \item \textbf{Traffic Cop Routing:} Across all categories, SLM routing consumes an average of $\TerraOverallDecoupledRoutingMS\text{ ms}$ ($\TerraCatADecoupledRoutingMS\text{ ms}$ on Category~A, $\TerraCatBDecoupledRoutingMS\text{ ms}$ on Category~B).
    \item \textbf{Graph Retrieval:} Retrieval latency, isolated from routing, averages $\TerraOverallDecoupledRetrievalMS\text{ ms}$ overall ($\TerraCatADecoupledRetrievalMS\text{ ms}$ on Category~A, $\TerraCatBDecoupledRetrievalMS\text{ ms}$ on Category~B, $\TerraCatCDecoupledRetrievalMS\text{ ms}$ on Category~C and $\TerraCatDDecoupledRetrievalMS\text{ ms}$ on Category~D). The Category~A fast path does not bypass retrieval: all 100 Category~A queries recorded retrieval time.
    \item \textbf{NLI Smart Grader Verification:} Logged grading time averages $\TerraOverallDecoupledNliMS\text{ ms}$ overall ($\TerraCatADecoupledNliMS\text{ ms}$ on Category~A, $\TerraCatBDecoupledNliMS\text{ ms}$ on Category~B, $\TerraCatCDecoupledNliMS\text{ ms}$ on Category~C and $\TerraCatDDecoupledNliMS\text{ ms}$ on Category~D). Category~C's low mean reflects its contexts: 74 of 75 Category~C queries retrieved an empty context (mean $\TerraCatCMeanRawChars$ characters across all 75) and recorded grading time under 1~ms, consistent with a deterministic rejection rather than an LLM call. Category~D contexts are longer (mean $\TerraCatDMeanRawChars$ characters, none above the 60{,}000-character clamp), and its grading time is correspondingly higher.
    \item \textbf{70B Answer Generation:} Generative synthesis averages $\TerraOverallDecoupledGenerationMS\text{ ms}$ overall. On Category~B, generation requires $\TerraCatBDecoupledGenerationMS\text{ ms}$ due to comprehensive multi-paragraph legal synthesis; on Category~C, generator latency is $\TerraCatCDecoupledGenerationMS\text{ ms}$ because no Category~C query reached the generator.
\end{enumerate}

\subsubsection{Tier 2 Geometric Domain Boundary Pre-Filter Calibration}
Table~\ref{tab:computational_budget} (Panel~B) presents empirical calibration for the Tier~2 geometric proximity pre-filter across distance thresholds $\tau \in [0.90, 1.40]$ on the 20,000-case CAP index, evaluated against $N=200$ in-domain legal queries (Categories~A and B) and $N=75$ out-of-domain probes (Category~C).

The isolated pre-filter demonstrates strong discriminative separation, achieving an AUROC of $\TerraGateIsolatedAUROC$. At aggressive thresholds ($\tau \le 1.10$), the filter rejects out-of-domain queries at high rates ($\TerraGateTPRTauZeroNineZero$ at $\tau=0.90$) but induces unacceptable false refusals on legitimate constitutional inquiries ($\TerraGateFPRTauZeroNineZero$, $\TerraGateFPRCountTauZeroNineZero$). As the threshold relaxes to $\tau=1.25$ and $\tau=1.30$, in-domain false refusals collapse to zero ($\TerraGateFPRTauOneThreeZero$, $\TerraGateFPRCountTauOneThreeZero$).

At the configured threshold $\tau=1.30$, the offline calibration shows zero in-domain false refusal and flags $\TerraGateTPRTauOneThreeZero$ ($\TerraGateTPRCountTauOneThreeZero$) of the out-of-domain probes by distance alone; the remaining $\TerraGateFNRTauOneThreeZero$ ($\TerraGateFNRCountTauOneThreeZero$) fall below the threshold. This is a calibration of a candidate pre-filter, not a description of the deployed run: in the canonical run every Category~C query recorded retrieval time (at least 178~ms excluding routing), so no query was terminated before retrieval, and all 75 Category~C queries were refused downstream, 74 of them on an empty retrieved context with a recorded grading time under 1~ms. Geometric proximity is an \textit{optional coarse pre-filter} intended to conserve retrieval and grading compute on obvious out-of-domain inputs; embedding distance does not prove legal validity or domain membership.

\subsection{Benchmark-Graph Traversal Stability Under Temporal Snapshot Expansion}
\label{subsec:temporal-update-stability}

To evaluate precedent traversal stability under corpus growth without generative confounding, we test standalone breadth-first graph traversal across four chronological snapshot expansions: $1890\to1930$ (Pair~I, Step $N=40$~yr), $1930\to1950$ (Pair~II, Step $N=20$~yr), $1950\to1970$ (Pair~III, Step $N=20$~yr), and $1970\to1980$ (Pair~IV, Step $N=10$~yr).



Across $N=\IssueNineteenTotalEvals$ qualifying Category~B observations ($\IssueNineteenUniqueQueries$ unique queries anchored at an origin precedent decided at or before Year~$T$, with at least one subsequent gold authority in $(T, T+N]$), BFS citation traversal ($k=2, d=2$, matching production parameters) was run on base subgraph $G_{\le T}$ and expanded subgraph $G_{\le T+N}$ seeded by benchmark origin precedent ($s = q[\texttt{target\_nodes}][0]$). As reported in Table~\ref{tab:temporal-update-stability}, expanding from $T$ to $T+N$ shifts grand pooled recall from $\IssueNineteenOverallRecallTMeanStd$ to $\IssueNineteenOverallRecallTNMeanStd$ ($\Delta\text{Recall} = \IssueNineteenOverallDeltaRecallMeanStd$), with $\IssueNineteenStabilityCoefficient$ of queries exhibiting non-negative recall shift ($\Delta\text{Recall} \ge 0$). On newly eligible authorities ($S_{\text{new}}$), traversal attains a grand mean Discovery Rate of $\IssueNineteenDiscoveryRateMean$, with median traversal latency under one millisecond (\IssueNineteenBFSLatencyMedianMs~ms; grand mean $0.15 \pm 0.18$~ms). Three methodological boundaries govern these findings: (1)~Because gold sequences include seed anchor $s$ (retrieved at depth 0), absolute recall reflects baseline seed inclusion; $\Delta\text{Recall}$ is invariant to this offset. (2)~Negative shifts on specific queries stem from degree caps when expanding nodes become hubs. (3)~This standalone experiment evaluates topological traversal over available directed edges (generic \texttt{CITES}) on NetworkX DiGraphs; it is decoupled from upstream vector retrieval (subject to the category filter limitation in Section~\ref{subsec:limitations}) and does not evaluate live vector index upsert, which remains future systems engineering.

\section{Discussion}
\label{sec:discussion}

The empirical evaluation validates that TERRA addresses fundamental failure modes of standard RAG architectures in knowledge-intensive, dynamic environments. This section contextualizes our primary architectural contributions, candidly examines current system limitations, and outlines future research directions.

\subsection{Architectural Differentiation and System Scope}
\label{subsec:contrib-discussion}
Compared to Microsoft Research's GraphRAG framework \cite{edge2024graphrag}---which constructs static hierarchical entity-relation communities for dataset summarization---our architecture directly addresses temporal precedent evolution and knowledge drift. While Microsoft GraphRAG requires substantial offline LLM summarization costs and remains temporally blind to overruled holdings, our dual-index Event Evolution Graph (EEG) cross-links chronological citation edges (\texttt{CITES}) with authoritative typed \texttt{OVERRULES} edges imported from CRS annotations (\IssueTwentyOneOverrulesEdges\ total, from which the 19 Category~D invalidation pairs are drawn; see Section~\ref{subsec:phase1-ingestion}) and node-level procedural reasoning traces. Furthermore, our NLI Smart Grader is designed to evaluate semantic entailment and reduce unsupported generations, although residual errors remain \cite{asai2024selfrag, zheng2023judging}. We formally delineate TERRA's architectural scope: the offline ingestion pipeline indexes full-docket Supreme Court jurisprudence (1850--1980) as a broad, high-distractor knowledge graph, while the online evaluated task focuses specifically on longitudinal precedent tracking, temporal drift mitigation, and safety firewalling over constitutional and federal public-law doctrine.

Three core design innovations drive TERRA's performance:
\begin{enumerate}
    \item \textbf{Multi-Tier Domain Boundary Safety:} Rather than relying on static allowlists, TERRA pairs lexical anchoring with downstream NLI entailment gating and an optional scale-dependent geometric proximity flag (which terminated no query in the canonical run).
    \item \textbf{Entailment-Gated Generation and XAI Attribution:} By replacing keyword heuristics with an NLI Smart Grader, TERRA operationalizes a ``prefer silence over hallucination'' principle with transparent reasoning paths, qualified by our \IssueTwentyOneAuditTargetConfirmedPct\ citation extraction precision on the underlying precedent graph (Section~\ref{subsec:phase1-ingestion}).
    \item \textbf{Citation Graphs as Temporal Reasoning Structures:} By structuring opinions as chronological citation edges (\texttt{CITES}) alongside authoritative typed \texttt{OVERRULES} relations within an Event Evolution Graph (Section~\ref{subsec:phase1-ingestion}), TERRA explicitly models chronological precedent shifts.
\end{enumerate}

\subsection{Empirical Analysis of Firewall Component Contributions and Scale Dynamics}
\label{subsec:firewall-analysis}

A critical finding from our systematic post-hoc auditing concerns the empirical division of labor within TERRA's safety architecture. The original design conceptualized the Domain Boundary Firewall as an active two-stage pre-retrieval filter where vector distance ($d_{L2} < 1.3$) flags out-of-domain (OOD) queries before graph traversal. However, systematic auditing across the preliminary 50-query Category C audit set on the initial 10,009-opinion index revealed that the vector proximity gate was non-functional at this scale ($0/50$ intercepted) because unconstrained OOD distances ($0.281 \le d_{L2} \le 0.859$) fell entirely beneath $1.3$.

Consequently, safety performance in Table~\ref{tab:overall_summary} (\TerraCatCRejectionRate\ refusal on Category C, \TerraCatDRejectionRate\ refusal and $\TerraCatDRefutationRate$\ verified refutation on Category D) was delivered by downstream refusal rather than by the proximity gate, which terminated no query: most Category~C refusals followed an empty retrieved context (74 of 75), and the remaining refusals across categories followed a non-empty context. Removing NLI in ablation (Table~\ref{tab:ablation}, Ablation 2) reduced adversarial safety to 0.0\%, so the grader is necessary in that configuration; the ablation does not, however, apportion the canonical Category~C refusals between empty-context and entailment-based refusal.

Crucially, empirical calibration on the expanded 20,000-opinion substantive corpus demonstrates that geometric pre-filtering exhibits scale- and density-dependent behavior. As reported in Table~\ref{tab:computational_budget} (Panel~B), evaluating the full $N=75$ Category~C benchmark suite against the 20k index yields an isolated AUROC of $\TerraGateIsolatedAUROC$. At the configured threshold $\tau = 1.30$, the offline calibration gives $\TerraGateFPRTauOneThreeZero$ in-domain false refusal ($\TerraGateFPRCountTauOneThreeZero$) and flags $\TerraGateTPRTauOneThreeZero$ ($\TerraGateTPRCountTauOneThreeZero$) of the out-of-domain probes by distance; the deployed run did not terminate any query before retrieval. Comparing the 10k index (where 50/50 probes fell below $1.3$) to the 20k index (where 41/75 probes exceed $1.3$) reflects an empirical observation of corpus manifold density: expanding corpus volume across seven verified constitutional domains creates denser semantic clustering, pushing extraneous out-of-domain queries farther into the embedding margin. We do not assert any generalized theoretical scaling law from these two empirical corpus snapshots; rather, we treat geometric proximity strictly as an optional coarse pre-filter that could save retrieval and NLI compute on obvious outliers. In the canonical run, 81 of the 139 refusals occurred on an empty retrieved context (74 in Category~C and 7 in Category~D) and 58 on a non-empty context (2, 16, 1 and 39 in Categories~A--D).


\subsection{Limitations}
\label{subsec:limitations}

To ensure complete scientific transparency and maintain the highest standard of academic integrity, we identify eighteen methodological limitations of the current architecture:
\begin{enumerate}
    \item \textbf{Absence of Licensed Legal Practitioner Evaluation (Most Critical Limitation for Deployment):} Due to geographical and institutional resource constraints, this study did not include evaluation by licensed legal practitioners. All faithfulness and safety assessments rely on automated LLM-as-a-judge scoring (Section~\ref{subsec:eval-metrics}) and continuous lexical metrics (ROUGE-L) evaluated against researcher-constructed reference answers cross-referenced against public legal reference sources (Cornell LII, Oyez, CRS Constitution Annotated). While the adversarial critic persona exhibits a modest, statistically significant correlation with reference lexical grounding (Pearson $r = \JudgePearsonR, p = \JudgePearsonP$), automated scoring cannot replicate practitioner judgment on subtle constitutional doctrine. Specifically, practitioners would assess legal correctness, temporal consistency, citation validity, and hallucination behavior beyond continuous metrics. Furthermore, lacking an independent labeled validation set for the NLI Smart Grader (Section~\ref{subsubsec:terra-nli-grader}), the false-non-entailment rate remains empirically unquantified beyond ablation evidence. We consider human expert evaluation of both legal output quality and NLI verifier reliability essential prerequisites for production deployment, and release the complete 350-query benchmark suite and codebase to invite independent expert replication. The author spot-check described in Section~\ref{subsec:obscure-generalization} is a single non-expert review of a stratified sample and does not substitute for licensed practitioner review.
    \item \textbf{Metric Disentanglement and Refusal-Volume Masking:} Evaluating complex RAG pipelines solely via aggregate faithfulness induces severe metric distortion: systems with high epistemic refusal volume receive top marks for non-committal answers on unanswerable queries, masking generative weaknesses. Under our dual-judge protocol, TERRA has the highest overall ($\TerraOverallFaithfulnessMean$ vs.\ $\HippoRAGOverallFaithfulnessMean$; intervals overlap) and answer-only faithfulness ($\TerraAnswerOnlyFaithfulnessMean$ vs.\ $\HippoRAGAnswerOnlyFaithfulnessMean$) with fewer flagged refusals ($\TerraFlaggedRefusals$ vs.\ $\HippoFlaggedRefusals$), but answer-only scores are conditional on different answered subsets. On multi-hop reasoning TERRA leads HippoRAG by $\CatBTerraVsHippoRAGDelta$ ($q=\CatBTerraVsHippoRAGAdjQ$). Future benchmark protocols must mandate disentangled reporting of refusal rates alongside answer-only faithfulness.
    \item \textbf{Evaluation Rubric Resolution and Likert Normalization:} Coarse evaluation scales (e.g., discrete binary pass/fail or $\{0.0, 0.8, 1.0\}$ heuristics) induce artificial ceiling compression, clustering over $75\%$ of responses at $1.0$. Implementing a 5-point Likert scale normalized to $[0, 1]$ effectively resolves this compression, exposing genuine variances in doctrinal nuance.
    \item \textbf{Jurisdictional and Temporal Corpus Boundaries:} The pre-built Event Evolution Graph is anchored in \CorpusTotalOpinions\ authentic substantive U.S. Supreme Court opinions spanning \CorpusTemporalSpan, encompassing seven doctrinal domains of federal jurisprudence (Table~\ref{tab:corpus_stats}). Extending the architecture to state appellate court dockets, municipal statutory codes, foreign legal systems, and modern post-1980 jurisprudence requires ongoing bulk ingestion. Furthermore, while Section~\ref{subsec:temporal-update-stability} empirically demonstrates graph traversal stability across temporal snapshot expansions ($T \to T+N$), live vector index insertion and streaming graph maintenance remain unaddressed engineering extensions.
    \item \label{item:hosted-api-determinism}\textbf{Hosted API Output Variability at $T=0.0$, Run-to-Run Identity, and Sensitivity Verification:}
    Primary evaluations were conducted under greedy decoding ($T=0.0$) using \texttt{databricks-meta-llama-3-3-70b-instruct} and \texttt{databricks-gemma-3-12b}. While setting $T=0.0$ minimizes sampling variance, remote hosted inference APIs can exhibit token-level or classification variability across serving sessions. Across benchmark evaluations, query \texttt{D066} exhibited an empirical classification shift (EASY $\to$ HARD), consistent with numerical or serving-level variability in hosted inference; the precise backend cause cannot be isolated from available logs. Comparing two 14-September TERRA executions across the 350-query benchmark suite, exact string-identical generated answers occurred on \TerraIdenticalAnswerFraction\ of queries overall (206/350), with pronounced category divergence: \TerraIdenticalAnswerCatA\ on Category~A (95/100) and \TerraIdenticalAnswerCatC\ on Category~C (62/75), versus \TerraIdenticalAnswerCatD\ on Category~D (39/75) and \TerraIdenticalAnswerCatB\ on Category~B (10/100). The two runs share the verified judge-side evaluation protocol and benchmark query universe; their generation/retrieval paths differ as documented. Both runs were evaluated under the standardized dual-judge protocol (\texttt{databricks-qwen3-next-80b-a3b-instruct} and \texttt{databricks-gpt-oss-120b}) using the frozen RAG prompt template (SHA256: \texttt{6d18fc14\dots}) and normalized Likert rubric. Sensitivity verification on an 80B model (\texttt{databricks-qwen3-next-80b-a3b-instruct}) across a separate 15-query subset (10 Cat~B, 5 Cat~D; distinct from the firewall stress test in Section~\ref{subsec:benchmark-design}) confirmed invariant relative rankings: TERRA achieved $1.000$ Category~B Faithfulness (vs. $0.900$ Flat RAG, $0.800$ Direct LLM) and $100\%$ Category~D safety rejection (vs. $40\%$ Flat RAG, $0\%$ Direct LLM). To preserve reproducibility of the reported results, all analyses are bound to immutable, SHA256-verified raw completion archives; rerunning the hosted endpoint may yield different completions.
    \item \textbf{Parametric Pretraining Confounders, Latent Weights, and Temporal Containment Bounds:} On continuous evaluation of Category D (Adversarial Deception), Direct LLM recorded a higher continuous Faithfulness score ($\DirectLLMCatDFaithfulness$ vs. $\TerraCatDFaithfulness$), and the paired Wilcoxon signed-rank comparison significantly favours Direct LLM ($\Delta = -0.200$, $p = 3.81 \times 10^{-5}$, FDR $q = 1.27 \times 10^{-4}$; Table~\ref{tab:full_fdr_family}, Rank 6). Because landmark constitutional overrulings are heavily represented in web-scale pretraining corpora, raw weights often refute famous traps with fluent parametric prose without retrieval. However, evaluating discrete outcome safety via the formal rubric (Section~\ref{subsec:benchmark-design}, Table~\ref{tab:refusal_breakdown}) reveals that Direct LLM achieves its $94.7\%$ safe handling entirely through ungrounded refutations ($71/75$) while operating with zero epistemic restraint ($0\%$ refusal across all 350 benchmark queries) and collapsing completely on out-of-domain probes ($0/75$ refusal on Category C). In contrast, TERRA achieves $\mathbf{\TerraCatDSafeHandlingRate}$ total safe handling ($73/75$: $\TerraCatDRejection$ explicit refusals and $\TerraCatDRefutations$ verified refutations), delivering a $+\TerraCatDRejectionRate$ explicit refusal advantage over Direct LLM ($p = \TerraVsDirectLLMCatDRefusalMcNemarP$; Holm-adjusted $p = \TerraVsDirectLLMCatDRefusalMcNemarHolmP$), no baseline-safe query handled unsafely ($c=0$), and fewer fatal constraint violations ($\TerraCatDFatalErrors$ versus $\BaselineFatalMin$--$\BaselineFatalMax$ of 75 across the four baselines) with zero uncommitted hedges ($0/75$). While the pre-specified hypothesis test in Table~\ref{tab:full_fdr_family} (Rank 6) evaluates continuous faithfulness, TERRA's true structural advantage lies in verifiable citation provenance, epistemic restraint, and temporal boundary integrity. Forensic auditing reveals that the two remaining fatal errors (queries D033 and D034) are reciprocal misattributions between coupled apportionment landmarks (\textit{Baker v. Carr} [1962] and \textit{Reynolds v. Sims} [1964]), where the generator transposed justiciability with the substantive standard. Adaptive edge-weighting to prioritize invalidation paths constitutes essential future work. Similarly, on Category~B (Evolutionary Precedents), Direct LLM achieves higher judge-scored faithfulness ($\DirectLLMCatBFaithfulness$ vs.\ $\TerraCatBFaithfulness$, $\Delta=\CatBTerraVsDirectLLMDelta$, $q=\CatBTerraVsDirectLLMAdjQ$), consistent with base model fluency on generic historical narratives, whereas TERRA's retrieval evaluation shows stronger Category~B citation coverage and full-precedent-path recovery ($\CatBTerraMacroCitationRecall$ macro citation recall vs.\ $\CatBDirectMacroCitationRecall$ for Direct LLM; Section~\ref{subsec:citation-correctness}).

    More fundamentally, we distinguish \textit{latent parametric knowledge} from \textit{operative doctrinal leakage}. Because modern foundation models are pretrained on web-scale text, weights encode post-$T$ legal history (e.g., \textit{Plessy} overruled by \textit{Brown} in 1954). TERRA does not purge latent weights; its provenance-bounded retrieval and temporal gate prevent latent memory from intruding into operative pre-$T$ reasoning. While forensic auditing (Section~\ref{subsec:temporal-leakage}) confirms suppression of Level-2 leakage in $\text{TERRA}_{\text{gated}}$, future work should explore circuit editing and machine unlearning for weight-level guarantees. Provider-side safety filtering on historical legal topics (terminating 85 evaluations in Section~\ref{subsec:temporal-leakage}) highlights the operational dependency of LLM-as-a-judge pipelines on third-party API policies, motivating our sensitivity-bounded framework. Section~\ref{subsec:obscure-generalization} tests this confounder directly: on obscure overruled precedents Direct LLM safe handling was \ObscDirectSafePct\ (\ObscCrossDirectFamousSafe\ on the famous subset), consistent with famous-landmark performance partly reflecting parametric memory.
    \item \textbf{Fixed Graph Traversal Horizon and Multi-Hop Boundary Limits:} TERRA's online retrieval operates with a calibrated 2-hop BFS graph expansion radius to bound latency and token overhead. While 2-hop traversal captures direct citing and invalidating authorities across most evolutionary arcs, it can miss subsequent modifying precedents when citation chains exceed two steps. An empirical boundary emerged during our CaseFacts audit on \textit{Chimel v. California} (query \texttt{CF\_INC\_058}; Section~\ref{subsec:casefacts-generalization}), where the modifying precedent (\textit{United States v. Robinson} [1973]) sat beyond the 2-hop boundary; the generator cited \textit{Chimel}'s 1969 wingspan doctrine but found no modifying precedent in context. This multi-hop horizon likewise accounts for the gold-precedent omission rate during drafting on CLERC ($\ClercTerraMisapplicationTrackOne$; Section~\ref{subsec:clerc-generalization}), where distant authorities prompted analogical reasoning from historical cases rather than fabrication. Developing adaptive traversal horizons along invalidation paths represents a high-priority architectural extension.
\item \textbf{Oracle Relation Labels, Exact-Title Seeding, and a Bundled Ablation:} In Full TERRA, curated \texttt{OVERRULES} edges are serialized with an explicit relation tag, dedicated heading, expanded excerpt budget, and prompt instruction; obscure benchmark pairs were drawn from these edges. The Full-TERRA advantage on that suite is therefore contingent on curated annotations, and the bundled ablation cannot isolate the tag alone. For named-case queries, seeding relied on exact caption matches (category-filtered vector retrieval returned no candidates), causing two queries to produce empty context and safe-rated refusals, one to omit the overruling case, and two to take the fast path; we evaluated TERRA as deployed and report sensitivity subsets excluding these items. The primary comparison contrasts complete systems, and Direct LLM receives no instruction block. Outcomes are judge-dependent in magnitude, the suite has thirty items, and judges were not fully blind to arm identity.

        \item \label{item:falserefusal}\textbf{False Refusals on Answerable Queries (Categories A and B):}
The safety gate produces false refusals on non-adversarial inquiries: \TerraCatAFalseRefusals/100 Category~A (Factual) and \TerraCatBFalseRefusals/100 Category~B (Evolutionary) queries receive system-level safety rejections (\texttt{safety\_rejected = True}). All \TerraCatBFalseRefusals\ Category~B false refusals originate from the NLI Smart Grader on the hard route; none occur on the fast route. Individual causes were not separately diagnosed.

\item \textbf{Latency Measurement Caveats:}
Three factors limit direct latency comparisons. (1)~\textit{Execution mode}: Baselines ran sequentially on 5--6~September~2026; TERRA ran with four concurrent workers on 14~September~2026. (2)~\textit{HTTP~429 throttling}: \TerraThrottledRows/350 TERRA queries incurred rate-limiting back-off accumulating \TerraThrottleSleepSeconds\,s of sleep; reported medians are gross with net values alongside. (3)~\textit{Bimodal routing}: The fast route contains \RouterFastPathCount\ queries (median \TerraFastRouteMedianS); the hard route contains \RouterHardPathCount\ queries (median \TerraHardRouteMedianS). Aggregate medians obscure this bimodal structure; baseline and TERRA comparisons are indicative only.

\item \label{item:samefamily}\textbf{Same-Family Generation, NLI Grading, and Category~D Verifier:}
Generation, NLI grading, and the Category~D verifier all deploy \texttt{databricks-meta-llama-3-3-70b-instruct} (same family and endpoint). Faithfulness judges are \texttt{databricks-qwen3-next-80b-a3b-instruct} and \texttt{databricks-gpt-oss-120b} (independent families). Category~D four-class labels (refusal, refutation, fatal error, hedge) may inherit shared Llama biases; cross-family verifier ablation has not been run.

\item \textbf{Unflagged Textual Refusals:}
System-level flagging (\texttt{safety\_rejected = True}) undercounts epistemic refusals. Adding textual refusals (asserting contextual absence without gate triggering) raises counts to \TerraEpistemicRefusals/350 for TERRA, \FlatEpistemicRefusals/350 for Flat Vector RAG, \GraphEpistemicRefusals/350 for GraphRAG-Adapted, and \HippoEpistemicRefusals/350 for HippoRAG-Adapted. Unflagged TERRA cases are D034 (refusal phrasing with fatal error, score 0.000), D061 (verified refutation, floored at \AmendmentMFloorScore\ under Amendment~M), and D068/D070 (fast-path textual refusals scoring 1.000). Aggregate Table~IX sensitivity results for TERRA are invariant to this distinction ($\Delta = 0.000$).

\item \textbf{Mixed-Date Evaluation Blocks:}
Not all blocks were re-evaluated under v2 dual-judge protocols; Table~\ref{tab:run_provenance} details dates and backends. Blocks lacking timestamps (pilot ablation, temporal leakage probe, CaseFacts, CLERC) carry ``date not recorded''; scores derive from earlier single-judge evaluation and are not numerically comparable to v2 evaluations.

\item \textbf{Corpus Length-Filtering Bias and Concise Landmark Exclusion:}
Filtering by length ($\ge 4{,}000$ characters) purges noise: across \IssueTwentyTotalVolumesScanned\ volumes (\IssueTwentyTotalDecisionsScanned\ cases), \IssueTwentySubFiveHundredCharCount\ records are $<500$ characters (\IssueTwentySubFiveHundredCharPct, procedural orders/dockets). However, equating length with authority induces bias: it excludes \IssueTwentyIntermediateExcludedCount\ decisions ($500$--$3{,}999$ characters, \IssueTwentyIntermediateExcludedPct), with \IssueTwentyIntermediatePageRankGeNinetyNineFive\ in the top $0.5\%$ of CAP PageRank authority. Deterministic audit of 63 landmarks reveals \IssueTwentyLandmarksExcludedByLength\ (\IssueTwentyLandmarksExcludedByLengthPct) were purged, including procedural standards (\textit{Haines}, \IssueTwentyHainesCharCount\ chars, \IssueTwentyHainesCAPPRPercentile th percentile; \textit{Diebold}, \IssueTwentyDieboldCharCount\ chars, \IssueTwentyDieboldCAPPRPercentile th percentile), competency (\textit{Dusky}, \IssueTwentyDuskyCharCount\ chars, \IssueTwentyDuskyCAPPRPercentile th percentile), tax milestones (\textit{Lucas v. Earl}, \IssueTwentyLucasCharCount\ chars), and desegregation orders (\textit{Sipuel}, \IssueTwentySipuelCharCount\ chars; \textit{Gayle}, \IssueTwentyGayleCharCount\ chars). \textit{Bolling v. Sharpe} (\IssueTwentyBollingCharCount\ chars) met the threshold but was omitted in subsampling. Future graphs should replace static floors with citation retention.

\item \textbf{Citation Extraction Precision, Same-Year Chronology, and Provenance Uncertainty:}
Automated citation extraction via \texttt{eyecite} establishes citation existence rather than doctrinal endorsement. Our stratified audit ($N=\IssueTwentyOneAuditSampleN$, seed \IssueTwentyOneAuditSeed) confirms \IssueTwentyOneAuditCitationPresentPct\ of sampled edges contain reporter strings; direct precedential reliance is verified in \IssueTwentyOneAuditTargetConfirmedPct\ (95\% CI: \IssueTwentyOneAuditTargetConfirmedCI), with \IssueTwentyOneAuditQuotationPct\ in quotations. Furthermore, \IssueTwentyOneSameYearEdges\ edges (\IssueTwentyOneSameYearEdgesPct) connect same-year opinions, with \IssueTwentyOneIntraYearReversals\ intra-year date reversals from publication artifacts and zero cross-year forward violations. The legacy claim of ``$99.91\%$ chronological validity / $0.09\%$ noise'' in earlier working drafts was untraceable to repository artifacts, mandating this audit. Citation trails provide candidate authorities with \IssueTwentyOneAuditTargetConfirmedPct\ precision rather than certified doctrinal paths.

\item \textbf{Rule-Based Doctrinal Taxonomy, Lexical Polysemy, Truncation, and Residual Fallback Bias:}
The seven-domain classification is generated by a deterministic \IssueTwentyTwoTotalLexiconPhrases-phrase rule-based lexicon over truncated texts ($\le \IssueTwentyTwoTextTruncationChars$ characters), inducing four limitations: (1)~\textit{Lexical polysemy}: Keyword matching without contextual disambiguation risks misclassification when terms appear in procedural dicta or dissenting quotes. (2)~\textit{Truncation artifacts}: Truncating texts to \IssueTwentyTwoTextTruncationChars\ characters restricts matching to introductory syllabi, occasionally missing dispositive holdings. (3)~\textit{Residual fallback bias}: The \texttt{general\_constitutional\_statutory} category acts as an unkeyed fallback for \IssueTwentyTwoGeneralConstitutionalCorpus\ decisions (\IssueTwentyTwoGeneralConstitutionalCorpusPct) lacking substantive keywords, functioning as a catch-all with \IssueTwentyTwoGeneralConstitutionalBaselineAcc\ baseline precision despite \IssueTwentyTwoFallbackValidityPct\ validity as a residual class. (4)~\textit{Absence of human expert annotation}: Taxonomy validation relies on dual automated evaluators (IEA \IssueTwentyTwoIEARawAgreement, Cohen's $\kappa=\IssueTwentyTwoIEACohensKappa$) and a blinded automated adjudicator (\IssueTwentyTwoHorvitzThompsonAccuracy\ weighted accuracy, 95\% CI: \IssueTwentyTwoHorvitzThompsonCI), not licensed practitioner annotation; classifications are not certified human ground truth.

\item \textbf{A Priori Hyperparameter Selection and Sensitivity Analysis Scope:} Core retrieval and graph expansion parameters ($k=\IssueTwentyFourSeedK$, depth $d=\IssueTwentyFourBfsDepth$, hub cap $\IssueTwentyFourHubCap$, geometric threshold $\tau=\IssueTwentyFourDistThreshold$) were fixed prior to benchmark construction and evaluation based on topological network constraints and preliminary feasibility testing on an independent \IssueTwentyFourDevQueryCount-query development suite (\IssueTwentyFourDevDate). While this frozen protocol prevents the audited forms of test-set leakage and hyperparameter overfitting, it precludes claiming global parameter optimality. Empirical sensitivity sweeps across \IssueTwentySixConfigCount\ configurations (Section~\ref{subsec:retrieval-quality}, Table~\ref{tab:retrieval_eval} Panel~C) profile performance across horizons ($k \in \{1,2,4\}$, $d \in \{1,2,3\}$, $C \in \{25,50,100\}$), confirming empirical stability around the historical configuration ($k=2, d=2, C=50$, \IssueTwentySixHistMacro\ macro recall, \IssueTwentySixHistFullPath\ full path recall) while quantifying the coverage-efficiency trade-off under bounded operational caps (\IssueTwentySixGoldRetentionPct\ landmark retention, \IssueTwentySixContextReductionPct\ context reduction).
    \item \textbf{Task-Corpus Domain Asymmetry and Private-Law Generalizability Scope:} A core methodological boundary is the intentional asymmetry between corpus composition and benchmark task scope. While the indexed repository reflects the complete 20,000-case docket spanning seven domains (\IssueTwentyFiveCorpusCombinedPrivateStatutoryPct\ private and general statutory law), the evaluated multi-hop reasoning benchmark (Category~B) evaluates constitutional and public-law precedent exclusively (\IssueTwentyFiveCatBConstitutionalPublicPct\ across five public-law domains; zero private-law chains). While the evaluation demonstrates TERRA's measured performance on the evaluated constitutional citation chains across a dense background network containing \CorpusPrivateLawOpinions\ private-law distractors, the system's performance on complex multi-hop statutory or patent reasoning remains unvalidated. Multi-hop evolutionary retrieval over complex commercial codes and specialized regulatory dockets constitutes an important open frontier for future domain-specific evaluation.
\end{enumerate}

\subsection{Future Work}
\label{subsec:future-work}
Future research will: (1)~conduct licensed practitioner evaluation across the 350-query benchmark; (2)~scale corpus ingestion to the full federal repository; and (3)~codify \texttt{OVERRULES} annotations as explicit weighted graph edges with adaptive invalidation-path traversal.

\section{Conclusion}
\label{sec:conclusion}

This paper introduced TERRA, an experimental dual-phase GraphRAG architecture addressing knowledge drift and multi-step legal reasoning on historical SCOTUS jurisprudence (\CorpusTemporalSpan). Structuring \CorpusTotalOpinions\ authentic decisions (\CorpusPrecedentEdges\ citation edges) with procedural reasoning traces, TERRA enables inference via Traffic Cop routing, BFS traversal, and an NLI grader evaluating retrieved evidence. Benchmarking across \numUniqueQueries\ queries (\numTotalEvals\ evaluations) shows higher Category~B multi-hop faithfulness than Flat Vector RAG, HippoRAG, and GraphRAG ($\Delta$ from $\CatBTerraVsHippoRAGDelta$ to $\CatBTerraVsGraphRAGDelta$; $q<0.05$), though not over Direct LLM ($\CatBTerraVsDirectLLMDelta$); higher overall faithfulness ($\TerraOverallFaithfulnessMean$ vs.\ $\HippoRAGOverallFaithfulnessMean$, overlapping CIs; $\TerraAnswerOnlyFaithfulnessMean$); and $\TerraCatDSafeHandlingRate$ adversarial safe handling ($\TerraCatDRejection$ refusals, \TerraCatDRefutations\ refutations). Constraints include \TerraCatDFatalErrors\ fatal errors, false refusals on \TerraCatAFalseRefusals\ Category~A and \TerraCatBFalseRefusals\ Category~B queries, and reliance on curated overruling annotations (Section~\ref{subsec:obscure-generalization}). Critically, because evaluations rely on automated LLM judges without licensed attorneys (Section~\ref{subsec:limitations}), human expert review remains an essential prerequisite for deployment. Citation graphs grounded with procedural traces preserve longitudinal stability (Section~\ref{subsec:temporal-update-stability}) \cite{edge2024graphrag, sun2024thinkongraph, zheng2023judging}.

%\newpage
\appendices
\section{Complete Benjamini-Hochberg FDR Hypothesis Testing Family}
\label{sec:appendix-fdr}

Table~\ref{tab:full_fdr_family} reports the pre-specified family of $M=20$ hypothesis tests across four comparative baselines and five evaluation slices ($q_{(i)} = \min_{k \ge i} \{ \frac{20 \cdot p_{(k)}}{k} \}$). Ranks 1--4 confirm post-correction advantage on Category C refusal and Category B evolutionary precedent reasoning.

\makeatletter
\global\@fptop=0pt
\global\@dblfptop=0pt
\makeatother
\begingroup
\let\oldfootnotesize\footnotesize
\renewcommand{\footnotesize}{\scriptsize}
\renewcommand{\arraystretch}{0.80}
\input{tables/table_full_fdr_family}
\endgroup

\clearpage
\begingroup
\let\oldfootnotesize\footnotesize
\renewcommand{\footnotesize}{\fontsize{5.0pt}{5.5pt}\selectfont}
\renewcommand{\IEEEbibitemsep}{1.0pt}
\bibliographystyle{IEEEtran}
\bibliography{ref}

\endgroup
\end{document}
